% Options for packages loaded elsewhere
\PassOptionsToPackage{unicode}{hyperref}
\PassOptionsToPackage{hyphens}{url}
\PassOptionsToPackage{dvipsnames,svgnames,x11names}{xcolor}
%
\documentclass[
  a4paper,
  13pt]{extreport}
\usepackage{amsmath,amssymb}
\usepackage{setspace}
\usepackage{iftex}
\ifPDFTeX
  \usepackage[T1]{fontenc}
  \usepackage[utf8]{inputenc}
  \usepackage{textcomp} % provide euro and other symbols
\else % if luatex or xetex
  \usepackage{unicode-math} % this also loads fontspec
  \defaultfontfeatures{Scale=MatchLowercase}
  \defaultfontfeatures[\rmfamily]{Ligatures=TeX,Scale=1}
\fi
\usepackage{lmodern}
\ifPDFTeX\else
  % xetex/luatex font selection
  \setmainfont[]{Times New Roman}
\fi
% Use upquote if available, for straight quotes in verbatim environments
\IfFileExists{upquote.sty}{\usepackage{upquote}}{}
\IfFileExists{microtype.sty}{% use microtype if available
  \usepackage[]{microtype}
  \UseMicrotypeSet[protrusion]{basicmath} % disable protrusion for tt fonts
}{}
\usepackage{xcolor}
\usepackage[left=3cm,right=2cm,top=2cm,bottom=2cm,headsep=1cm,footskip=1cm]{geometry}
\usepackage{color}
\usepackage{fancyvrb}
\newcommand{\VerbBar}{|}
\newcommand{\VERB}{\Verb[commandchars=\\\{\}]}
\DefineVerbatimEnvironment{Highlighting}{Verbatim}{commandchars=\\\{\}}
% Add ',fontsize=\small' for more characters per line
\newenvironment{Shaded}{}{}
\newcommand{\AlertTok}[1]{\textcolor[rgb]{1.00,0.00,0.00}{\textbf{#1}}}
\newcommand{\AnnotationTok}[1]{\textcolor[rgb]{0.38,0.63,0.69}{\textbf{\textit{#1}}}}
\newcommand{\AttributeTok}[1]{\textcolor[rgb]{0.49,0.56,0.16}{#1}}
\newcommand{\BaseNTok}[1]{\textcolor[rgb]{0.25,0.63,0.44}{#1}}
\newcommand{\BuiltInTok}[1]{\textcolor[rgb]{0.00,0.50,0.00}{#1}}
\newcommand{\CharTok}[1]{\textcolor[rgb]{0.25,0.44,0.63}{#1}}
\newcommand{\CommentTok}[1]{\textcolor[rgb]{0.38,0.63,0.69}{\textit{#1}}}
\newcommand{\CommentVarTok}[1]{\textcolor[rgb]{0.38,0.63,0.69}{\textbf{\textit{#1}}}}
\newcommand{\ConstantTok}[1]{\textcolor[rgb]{0.53,0.00,0.00}{#1}}
\newcommand{\ControlFlowTok}[1]{\textcolor[rgb]{0.00,0.44,0.13}{\textbf{#1}}}
\newcommand{\DataTypeTok}[1]{\textcolor[rgb]{0.56,0.13,0.00}{#1}}
\newcommand{\DecValTok}[1]{\textcolor[rgb]{0.25,0.63,0.44}{#1}}
\newcommand{\DocumentationTok}[1]{\textcolor[rgb]{0.73,0.13,0.13}{\textit{#1}}}
\newcommand{\ErrorTok}[1]{\textcolor[rgb]{1.00,0.00,0.00}{\textbf{#1}}}
\newcommand{\ExtensionTok}[1]{#1}
\newcommand{\FloatTok}[1]{\textcolor[rgb]{0.25,0.63,0.44}{#1}}
\newcommand{\FunctionTok}[1]{\textcolor[rgb]{0.02,0.16,0.49}{#1}}
\newcommand{\ImportTok}[1]{\textcolor[rgb]{0.00,0.50,0.00}{\textbf{#1}}}
\newcommand{\InformationTok}[1]{\textcolor[rgb]{0.38,0.63,0.69}{\textbf{\textit{#1}}}}
\newcommand{\KeywordTok}[1]{\textcolor[rgb]{0.00,0.44,0.13}{\textbf{#1}}}
\newcommand{\NormalTok}[1]{#1}
\newcommand{\OperatorTok}[1]{\textcolor[rgb]{0.40,0.40,0.40}{#1}}
\newcommand{\OtherTok}[1]{\textcolor[rgb]{0.00,0.44,0.13}{#1}}
\newcommand{\PreprocessorTok}[1]{\textcolor[rgb]{0.74,0.48,0.00}{#1}}
\newcommand{\RegionMarkerTok}[1]{#1}
\newcommand{\SpecialCharTok}[1]{\textcolor[rgb]{0.25,0.44,0.63}{#1}}
\newcommand{\SpecialStringTok}[1]{\textcolor[rgb]{0.73,0.40,0.53}{#1}}
\newcommand{\StringTok}[1]{\textcolor[rgb]{0.25,0.44,0.63}{#1}}
\newcommand{\VariableTok}[1]{\textcolor[rgb]{0.10,0.09,0.49}{#1}}
\newcommand{\VerbatimStringTok}[1]{\textcolor[rgb]{0.25,0.44,0.63}{#1}}
\newcommand{\WarningTok}[1]{\textcolor[rgb]{0.38,0.63,0.69}{\textbf{\textit{#1}}}}
\usepackage{longtable,booktabs,array}
\usepackage{calc} % for calculating minipage widths
% Correct order of tables after \paragraph or \subparagraph
\usepackage{etoolbox}
\makeatletter
\patchcmd\longtable{\par}{\if@noskipsec\mbox{}\fi\par}{}{}
\makeatother
% Allow footnotes in longtable head/foot
\IfFileExists{footnotehyper.sty}{\usepackage{footnotehyper}}{\usepackage{footnote}}
\makesavenoteenv{longtable}
\usepackage{graphicx}
\makeatletter
\def\maxwidth{\ifdim\Gin@nat@width>\linewidth\linewidth\else\Gin@nat@width\fi}
\def\maxheight{\ifdim\Gin@nat@height>\textheight\textheight\else\Gin@nat@height\fi}
\makeatother
% Scale images if necessary, so that they will not overflow the page
% margins by default, and it is still possible to overwrite the defaults
% using explicit options in \includegraphics[width, height, ...]{}
\setkeys{Gin}{width=\maxwidth,height=\maxheight,keepaspectratio}
% Set default figure placement to htbp
\makeatletter
\def\fps@figure{htbp}
\makeatother
\setlength{\emergencystretch}{3em} % prevent overfull lines
\providecommand{\tightlist}{%
  \setlength{\itemsep}{0pt}\setlength{\parskip}{0pt}}
\setcounter{secnumdepth}{-\maxdimen} % remove section numbering
\ifLuaTeX
\usepackage[bidi=basic]{babel}
\else
\usepackage[bidi=default]{babel}
\fi
\babelprovide[main,import]{vietnamese}
\ifPDFTeX
\else
\babelfont{rm}[]{Times New Roman}
\fi
% get rid of language-specific shorthands (see #6817):
\let\LanguageShortHands\languageshorthands
\def\languageshorthands#1{}
% ctu_xelatex_header.tex — header-includes cho build LaTeX CTU (xelatex).
% Ánh xạ các ký tự Unicode mà TeX Gyre Termes thiếu glyph (chỉ số dưới/trên,
% ký tự kẻ hộp) sang lệnh LaTeX tương đương để không bị rớt glyph khi biên dịch.
\usepackage{newunicodechar}
\usepackage{pifont}
\usepackage{amssymb}
\newunicodechar{✓}{\ding{51}}
\newunicodechar{✗}{\ding{55}}
\newunicodechar{≥}{\ensuremath{\geq}}
\newunicodechar{≤}{\ensuremath{\leq}}
\newunicodechar{≈}{\ensuremath{\approx}}
\newunicodechar{↔}{\ensuremath{\leftrightarrow}}
\newunicodechar{₀}{\textsubscript{0}}
\newunicodechar{₁}{\textsubscript{1}}
\newunicodechar{₂}{\textsubscript{2}}
\newunicodechar{₃}{\textsubscript{3}}
\newunicodechar{₄}{\textsubscript{4}}
\newunicodechar{₅}{\textsubscript{5}}
\newunicodechar{₆}{\textsubscript{6}}
\newunicodechar{₇}{\textsubscript{7}}
\newunicodechar{₈}{\textsubscript{8}}
\newunicodechar{₉}{\textsubscript{9}}
\newunicodechar{²}{\textsuperscript{2}}
\newunicodechar{³}{\textsuperscript{3}}
\newunicodechar{─}{\rule[0.5ex]{0.9em}{0.4pt}}
\newunicodechar{∩}{\ensuremath{\cap}}
\newunicodechar{∪}{\ensuremath{\cup}}
\newunicodechar{×}{\ensuremath{\times}}
\newunicodechar{≠}{\ensuremath{\neq}}
\newunicodechar{±}{\ensuremath{\pm}}
\newunicodechar{⊇}{\ensuremath{\supseteq}}
\newunicodechar{⊆}{\ensuremath{\subseteq}}
\newunicodechar{∈}{\ensuremath{\in}}
\newunicodechar{∉}{\ensuremath{\notin}}
\newunicodechar{→}{\ensuremath{\rightarrow}}
\newunicodechar{√}{\ensuremath{\surd}}
\newunicodechar{·}{\ensuremath{\cdot}}

% --- Bảng rộng: thu nhỏ cỡ chữ + khoảng cột, cho phép ngắt token dài để
%     bảng không tràn lề phải (overfull hbox). Áp dụng cho mọi longtable
%     (pandoc sinh toàn bộ bảng dạng longtable).
\usepackage{etoolbox}
\usepackage{seqsplit}      % ngắt được token rất dài (mã biến, số dài)
% Lưu ý: KHÔNG đổi \tabcolsep — pandoc tính bề rộng cột theo công thức
% (\columnwidth - 4\tabcolsep)*frac (tham chiếu \tabcolsep trực tiếp), nên
% giảm \tabcolsep lại làm cột RỘNG hơn và tràn nặng hơn.
\AtBeginEnvironment{longtable}{%
  \footnotesize
  \emergencystretch=3em
  \hbadness=10000          % không cảnh báo underfull do RaggedRight
}
% Cho phép TeX bẻ dòng linh hoạt hơn trong ô p{} hẹp
\setlength{\emergencystretch}{2em}

% URL/DOI dài: cho phép ngắt ở bất kỳ ký tự nào (tránh tràn lề ở danh mục TLTK).
% Kết hợp với pandoc +autolink_bare_uris để URL trần được bọc \url{}.
\usepackage{xurl}
\ifLuaTeX
  \usepackage{selnolig}  % disable illegal ligatures
\fi
\IfFileExists{bookmark.sty}{\usepackage{bookmark}}{\usepackage{hyperref}}
\IfFileExists{xurl.sty}{\usepackage{xurl}}{} % add URL line breaks if available
\urlstyle{same}
\hypersetup{
  pdflang={vi},
  colorlinks=true,
  linkcolor={black},
  filecolor={Maroon},
  citecolor={black},
  urlcolor={blue},
  pdfcreator={LaTeX via pandoc}}

\author{}
\date{}

\begin{document}

\setstretch{1.2}
\thispagestyle{empty}
\begin{center}
{\fontsize{14}{18}\selectfont\bfseries BỘ GIÁO DỤC VÀ ĐÀO TẠO\\TRƯỜNG ĐẠI HỌC CẦN THƠ\\TRƯỜNG KINH TẾ}

\vspace{1.8cm}
{\fontsize{14}{18}\selectfont\bfseries ĐỖ THÙY HƯƠNG}

{\fontsize{13}{17}\selectfont MÃ SỐ NGHIÊN CỨU SINH: P1323001}

\vspace{1.6cm}
{\fontsize{14}{18}\selectfont\bfseries CHUYÊN ĐỀ TIẾN SĨ SỐ 2}

\vspace{0.7cm}
{\fontsize{17}{23}\selectfont\bfseries XÂY DỰNG MÔ HÌNH NGHIÊN CỨU VỀ ẢNH HƯỞNG CỦA QUỐC TẾ HÓA ĐẾN HIỆU QUẢ HOẠT ĐỘNG KINH DOANH CỦA CÁC DOANH NGHIỆP Ở CHÂU Á}

\vspace{1.5cm}
{\fontsize{14}{18}\selectfont\bfseries NGÀNH QUẢN TRỊ KINH DOANH}

{\fontsize{14}{18}\selectfont\bfseries MÃ NGÀNH: 9340101}

\vspace{1.1cm}
{\fontsize{13}{17}\selectfont\bfseries NGƯỜI HƯỚNG DẪN KHOA HỌC}

{\fontsize{13}{17}\selectfont\bfseries PGS.TS. PHAN ANH TÚ}

\vspace{1.3cm}
{\fontsize{14}{18}\selectfont\bfseries CẦN THƠ, NĂM 2026}
\end{center}
\newpage

\hypertarget{lux1eddi-cam-ux111oan}{%
\subsection{LỜI CAM ĐOAN}\label{lux1eddi-cam-ux111oan}}

Tôi xin cam đoan chuyên đề tiến sĩ số 2 với tiêu đề ``Xây dựng mô hình
nghiên cứu về ảnh hưởng của quốc tế hóa đến hiệu quả hoạt động kinh
doanh của các doanh nghiệp ở Châu Á'' là công trình nghiên cứu của riêng
tôi dưới sự hướng dẫn khoa học của PGS.TS. Phan Anh Tú, Trường Đại học
Cần Thơ. Các kết quả phân tích, mô hình, bảng biểu, hình vẽ trình bày
trong chuyên đề đều trung thực và có nguồn gốc rõ ràng. Tài liệu tham
khảo được trích dẫn theo chuẩn APA 7th. Những trích dẫn từ các nghiên
cứu khác đều được nêu rõ trong danh mục tài liệu tham khảo. Tôi xin chịu
trách nhiệm hoàn toàn về nội dung khoa học của chuyên đề này.

\emph{Cần Thơ, ngày \ldots{} tháng \ldots{} năm 2026}

\emph{Nghiên cứu sinh}

Đỗ Thùy Hương

\begin{center}\rule{0.5\linewidth}{0.5pt}\end{center}

\hypertarget{tuxf3m-tux1eaft}{%
\subsection{TÓM TẮT}\label{tuxf3m-tux1eaft}}

Chuyên đề này xây dựng mô hình nghiên cứu lý thuyết và thực nghiệm về
ảnh hưởng của quốc tế hóa đến hiệu quả hoạt động kinh doanh của các
doanh nghiệp ở các quốc gia châu Á trong giai đoạn 2009--2025. Khung lý
thuyết tích hợp bốn tầng, Uppsala Model (Johanson \& Vahlne, 1977,
2009), lý thuyết dựa trên nguồn lực (Barney, 1991; Wernerfelt, 1984), lý
thuyết thể chế (North, 1990; Khanna \& Palepu, 2010; Peng, 2003) và lý
thuyết thượng tầng quản trị (Hambrick \& Mason, 1984; Hambrick, 2007),
kết hợp với Khung chấp nhận số nền tảng (Banalieva \& Dhanaraj, 2019;
Verhoef et al., 2021; Stallkamp \& Schotter, 2021) làm nền tảng lý
thuyết. Lược khảo năm dạng hàm quan hệ quốc tế hóa, hiệu quả (I--P) gồm
tuyến tính, chữ U ngược, chữ S (S-curve), chữ M (M-curve), và gánh nặng
bắt buộc (forced penalty), cùng năm phân tích tổng hợp lớn giai đoạn
1980--2024, cho thấy khoảng trống thực nghiệm trong việc kiểm định đồng
thời phi tuyến và điều tiết đa tầng ở châu Á. Hệ giả thuyết H1--H6 được
phát triển: H1 phi tuyến dạng chữ U ngược hoặc chữ S; H2 chỉ số năng lực
công nghệ (TCI) điều tiết tích cực; H3a--H3b chỉ số chấp nhận số (DAI)
điều tiết phụ thuộc chế độ thể chế: có điều kiện ở Emerging/Frontier
(H3a, dương trong OLS nhưng phụ thuộc bối cảnh và có thể null trong kiểm
định nhân quả) và có điều kiện/null ở Advanced (H3b), phân biệt với TCI;
H4a--H4c đặc điểm nhà quản trị điều tiết (kinh nghiệm, học vấn, giới
tính, tách biệt dấu kỳ vọng); H5 thể chế nội địa điều tiết theo phổ thể
chế ICRV sáu nhóm; H6 tính không đồng nhất theo thời gian. Tám mô hình
thực nghiệm M0--M7 được thiết lập, từ cơ sở tuyến tính đến mô hình điều
tiết ba chiều cấp luận án. Chiến lược nhận dạng đa tầng và kế hoạch kiểm
định độ vững toàn diện được thiết kế cho mẫu gộp 88.869 doanh nghiệp ở
50 nền kinh tế (gồm Nhật Bản) châu Á và Thái Bình Dương, bao gồm 103 cặp
quốc gia-năm qua 14 mốc khảo sát WBES giai đoạn 2009--2025 (hồi quy
chính M2 N=81.022, M5 N=79.080). Bốn đóng góp chính: (i) khung tích hợp
4 tầng kết hợp Khung chấp nhận số nền tảng chưa có tương đương cho châu
Á; (ii) tách bạch TCI và DAI theo Mô hình Áp dụng số và Năng lực phụ
thuộc bối cảnh (CDCM); (iii) mô hình M7 điều tiết ba chiều; (iv) phân
loại nhánh Advanced thành tiên tiến đổi mới sáng tạo dẫn dắt (ICRV Nhóm
I) và tiên tiến tài nguyên dẫn dắt (ICRV Nhóm II). Ba bản thảo thực
nghiệm đồng hành, P3 Việt Nam (mục tiêu JED), P4 Singapore (mục tiêu
MIR), P5 Trung Quốc (mục tiêu IJOEM), cung cấp bằng chứng quốc gia cụ
thể kiểm định H1--H6 và xác nhận tính khả nghiệm của mô hình đề xuất.

Từ khóa: quốc tế hóa; hiệu quả hoạt động kinh doanh; châu Á; mô hình phi
tuyến; điều tiết; áp dụng số; thể chế; lý thuyết thượng tầng quản trị;
ICRV.

\begin{center}\rule{0.5\linewidth}{0.5pt}\end{center}

\hypertarget{abstract}{%
\subsection{ABSTRACT}\label{abstract}}

This specialty essay constructs a theoretical and empirical research
model on the effect of internationalization on firm performance across
Asian economies during 2009--2025. The integrated theoretical framework
combines four layers, the Uppsala Model (Johanson \& Vahlne, 1977,
2009), the Resource-Based View (Barney, 1991; Wernerfelt, 1984),
Institutional Theory (North, 1990; Khanna \& Palepu, 2010; Peng, 2003),
and Upper Echelons Theory (Hambrick \& Mason, 1984; Hambrick, 2007),
together with the Foundational Digital Adoption Framework (Banalieva \&
Dhanaraj, 2019; Verhoef et al., 2021; Stallkamp \& Schotter, 2021). A
review of five functional forms of the internationalization--performance
(I--P) relationship, linear, inverted-U, S-curve, M-curve, and forced
penalty, together with five major meta-analyses covering 1980--2024
reveals an empirical gap in the simultaneous testing of nonlinearity and
multi-level moderation in Asia. Six hypotheses (H1--H6) are developed:
H1 nonlinear inverted-U or S-curve relationship; H2 Technological
Capability Index (TCI) positively moderating; H3a--H3b Digital Adoption
Index (DAI) moderating contingent on institutional regime: conditional
for Emerging/Frontier (H3a), positive under OLS but subject to selection
bias (P3 IV: β = 0.018, p = .942), conditional or null for Advanced
(H3b), and distinguishable from TCI; H4a--H4c top manager
characteristics moderating (experience, education, and gender with
distinct directional expectations); H5 home-country institutions
moderating along the ICRV six-regime institutional spectrum; H6 temporal
heterogeneity. Eight empirical models M0--M7 are specified, from linear
baseline to three-way moderation. A multi-level identification strategy
and comprehensive robustness plan are designed for a pool of 88,869
firms across 50 Asian and Pacific economies (including Japan; 96,415
ICRV-classified across 52 source economies), spanning 103 country-year
pairs across 14 WBES survey waves (2009--2025; main regressions M2
N=81,022, M5 N=79,080). The four main contributions are: (i) an
integrated 4-tier framework combined with the Foundational Digital
Adoption Framework with no equivalent for Asia; (ii) the separation of
TCI and DAI within the Country--Digital--Capability Moderation (CDCM);
(iii) the M7 three-way moderation model; (iv) sub-grouping of the
Advanced regime into innovation-driven (ICRV Group I) and
resource-driven (ICRV Group II) types. Three companion empirical papers,
P3 Vietnam (target JED), P4 Singapore (target MIR), P5 China (target
IJOEM), provide country-specific evidence testing H1--H6 and validating
the proposed model structure.

Keywords: internationalization; firm performance; Asia; nonlinear
models; moderation; digital adoption; institutions; Upper Echelons;
ICRV.

\begin{center}\rule{0.5\linewidth}{0.5pt}\end{center}

\renewcommand{\contentsname}{MỤC LỤC}
\setcounter{tocdepth}{3}
\tableofcontents

\hypertarget{danh-mux1ee5c-bux1ea3ng}{%
\subsection{DANH MỤC BẢNG}\label{danh-mux1ee5c-bux1ea3ng}}

Bảng 2.1. Mỗi tầng lý thuyết, câu hỏi phụ, biến số, giả thuyết, bằng
chứng neo đậu

Bảng 2.2. Năm dạng hàm quốc tế hóa và hiệu quả, cơ chế, và bối cảnh xuất
hiện

Bảng 2.3. So sánh bằng chứng bốn bản thảo đồng hành P2--P5

Bảng 2.4. Mapping biến đầy đủ: khái niệm lý thuyết, biến WBES, kỳ vọng
dấu, giả thuyết

Bảng 2.5. Hệ giả thuyết H1--H6, tóm tắt cơ chế, biến, kỳ vọng, và bằng
chứng neo đậu

Bảng 2.6. CDCM, Dự đoán tác động DAI theo 3 chiều ngữ cảnh

Bảng 2.7. Tám mô hình M0--M7, cấu trúc, giả thuyết kiểm định, và biến bổ
sung

Bảng 2.8. Cấu trúc mẫu gộp theo nhóm ICRV và mốc khảo sát

Bảng 2.9. Đo lường TCI theo mã WBES

Bảng 2.10. Đo lường DAI theo mã WBES

Bảng 2.11. Phân loại 50 nền kinh tế theo ICRV, đặc điểm và tiêu chí

Bảng 2.12. Kế hoạch kiểm định độ vững, sáu nhóm

Bảng 2.13. So sánh đặc điểm phương pháp và phạm vi lý thuyết của các
khung tham chiếu chính và CĐ2

\begin{center}\rule{0.5\linewidth}{0.5pt}\end{center}

\hypertarget{danh-mux1ee5c-huxecnh}{%
\subsection{DANH MỤC HÌNH}\label{danh-mux1ee5c-huxecnh}}

Hình 2.1. Khung khái niệm tích hợp CĐ2, tổng quan

Hình 2.2. Khung khái niệm tích hợp CĐ2, chi tiết với bằng chứng neo đậu

\begin{center}\rule{0.5\linewidth}{0.5pt}\end{center}

\hypertarget{danh-mux1ee5c-tux1eeb-viux1ebft-tux1eaft}{%
\subsection{DANH MỤC TỪ VIẾT
TẮT}\label{danh-mux1ee5c-tux1eeb-viux1ebft-tux1eaft}}

\begin{longtable}[]{@{}
  >{\raggedright\arraybackslash}p{(\columnwidth - 2\tabcolsep) * \real{0.5000}}
  >{\raggedright\arraybackslash}p{(\columnwidth - 2\tabcolsep) * \real{0.5000}}@{}}
\toprule\noalign{}
\begin{minipage}[b]{\linewidth}\raggedright
Viết tắt
\end{minipage} & \begin{minipage}[b]{\linewidth}\raggedright
Diễn giải
\end{minipage} \\
\midrule\noalign{}
\endhead
\bottomrule\noalign{}
\endlastfoot
AIC & Akaike Information Criterion \\
BIC & Bayesian Information Criterion \\
BEE & Business Environment and Enterprise (WBES module) \\
BREADY & Business REAdy diagnostic (WBES module) \\
CDCM & Country--Digital--Capability Moderation, Mô hình Áp dụng số và
Năng lực phụ thuộc bối cảnh \\
DAI & Digital Adoption Index, Chỉ số áp dụng số \\
FDI & Foreign Direct Investment, Đầu tư trực tiếp nước ngoài \\
FE & Fixed Effects, Hiệu ứng cố định \\
FSTS & Foreign Sales to Total Sales ratio, Tỉ lệ doanh thu xuất khẩu
trên tổng doanh thu \\
GMM & Generalized Method of Moments \\
HC1/HC3 & Heteroskedasticity-Consistent Standard Errors (Long \& Ervin,
2000) \\
I--P & Internationalization--Performance, Quan hệ quốc tế hóa và hiệu
quả \\
IB & International Business, Kinh doanh quốc tế \\
ICRV & Institutional Context Regime Variation, Phân loại chế độ thể
chế \\
ICT & Information and Communications Technology \\
ISO & International Organization for Standardization \\
IV & Instrumental Variable, Biến công cụ \\
LOWESS & Locally Weighted Scatterplot Smoothing \\
MNE & Multinational Enterprise, Doanh nghiệp đa quốc gia \\
OECD & Organisation for Economic Co-operation and Development \\
OLS & Ordinary Least Squares, Bình phương tối thiểu thông thường \\
PICS3 & WBES schema thế hệ 1 (2009--2013) \\
R\&D & Research and Development, Nghiên cứu và phát triển \\
RBV & Resource-Based View, Quan điểm dựa trên nguồn lực \\
RESET & Regression Specification Error Test \\
ROA & Return on Assets, Lợi nhuận trên tổng tài sản \\
ROS & Return on Sales, Lợi nhuận trên doanh thu \\
SE & Standard Errors, Sai số chuẩn \\
SIDS & Small Island Developing States, Các quốc đảo nhỏ đang phát
triển \\
SME & Small and Medium-sized Enterprise, Doanh nghiệp vừa và nhỏ \\
SOE & State-Owned Enterprise, Doanh nghiệp nhà nước \\
TCI & Technological Capability Index, Chỉ số năng lực công nghệ \\
UE & Upper Echelons Theory, Lý thuyết cấp cao \\
VIF & Variance Inflation Factor \\
WBES & World Bank Enterprise Surveys, Khảo sát doanh nghiệp Ngân hàng
Thế giới \\
WGI & Worldwide Governance Indicators \\
\end{longtable}

\begin{center}\rule{0.5\linewidth}{0.5pt}\end{center}

\hypertarget{phux1ea7n-2-nux1ed9i-dung-chuyuxean-ux111ux1ec1}{%
\subsection{Phần 2: NỘI DUNG CHUYÊN
ĐỀ}\label{phux1ea7n-2-nux1ed9i-dung-chuyuxean-ux111ux1ec1}}

\hypertarget{ux111ux1eb7t-vux1ea5n-ux111ux1ec1}{%
\subsubsection{2.1 Đặt vấn đề}\label{ux111ux1eb7t-vux1ea5n-ux111ux1ec1}}

\hypertarget{giux1edbi-thiux1ec7u}{%
\paragraph{2.1.1 Giới thiệu}\label{giux1edbi-thiux1ec7u}}

Quốc tế hóa là chiến lược phát triển quan trọng của doanh nghiệp trong
hơn nửa thế kỷ qua, nhưng kết quả mà nó đem lại không đồng nhất. Nền
tảng lý thuyết của dòng nghiên cứu này bắt nguồn từ các công trình kinh
điển về doanh nghiệp đa quốc gia và đầu tư trực tiếp nước ngoài: lý
thuyết vòng đời sản phẩm (Vernon, 1966), lý thuyết lợi thế độc quyền của
doanh nghiệp (Hymer, 1976; Caves, 1971), và lý thuyết nội bộ hóa chi phí
giao dịch (Buckley \& Casson, 1976), về sau được tổng hợp trong khung
chiết trung OLI (Dunning, 1988). Trên nền đó, tài liệu kinh doanh quốc
tế (IB) đã ghi nhận năm dạng hàm khác nhau trong quan hệ giữa cường độ
quốc tế hóa và hiệu quả hoạt động kinh doanh của doanh nghiệp: tuyến
tính (Hsu \& Boggs, 2003), chữ U ngược (Gomes \& Ramaswamy, 1999; Hitt,
Hoskisson \& Kim, 1997), chữ S (S-curve; Lu \& Beamish, 2004;
Contractor, Kundu \& Hsu, 2003; Ruigrok \& Wagner, 2003), chữ M
(M-curve; Riahi-Belkaoui, 1998), và gánh nặng bắt buộc (forced penalty;
Glaum \& Oesterle, 2007). Sự đa dạng này phản ánh tính phụ thuộc ngữ
cảnh của hiệu ứng quốc tế hóa: cùng một mức cường độ xuất khẩu có thể
tạo ra lợi ích hoặc gánh nặng tùy theo năng lực doanh nghiệp, thể chế
quốc gia, và giai đoạn phát triển số (Bausch \& Krist, 2007; Kirca et
al., 2012; Marano et al., 2016; Wu, Fan \& Chen, 2022). Thực vậy, các
nghiên cứu gần đây sử dụng chiến lược nhận diện chặt chẽ hơn cho thấy
dạng chữ U ngược có thể suy yếu hoặc biến mất trong các mẫu gộp xuyên
quốc gia (Pisani, Garcia-Bernardo \& Heemskerk, 2020), càng củng cố lập
luận rằng chính dị biệt bối cảnh chi phối kết quả quốc tế hóa và hiệu
quả.

Châu Á trong giai đoạn 2009--2025 là phòng thí nghiệm tự nhiên độc đáo
cho nghiên cứu quan hệ quốc tế hóa và hiệu quả. Khu vực này tồn tại đồng
thời ít nhất sáu chế độ thể chế khác nhau theo khung phân loại ICRV:
nhóm tiên tiến đổi mới sáng tạo dẫn dắt (ICRV Nhóm I) gồm Singapore,
Hong Kong SAR, Hàn Quốc, Đài Loan, Israel; nhóm tiên tiến tài nguyên dẫn
dắt (ICRV Nhóm II) gồm Saudi Arabia, Qatar, Kuwait, Bahrain, Brunei,
Oman; nhóm trung bình cao (ICRV Nhóm III) gồm Trung Quốc, Malaysia, Thái
Lan, Kazakhstan, Armenia, Georgia; nhóm chuyển đổi thu nhập trung
bình--thấp (ICRV Nhóm IV, \texttt{Lower\_mid\_transition}) gồm Việt Nam,
Ấn Độ, Indonesia, Philippines, Mông Cổ, Bangladesh, Pakistan; nhóm đang
nổi (ICRV Nhóm V, \texttt{Emerging}) gồm 17 nền kinh tế (gồm Sri Lanka,
Jordan); và nhóm quốc đảo nhỏ (ICRV Nhóm VI, \texttt{SIDS\_small}) gồm 8
nền (7 Pacific + Timor-Leste). Tổng cộng 50 nền kinh tế (gồm Nhật Bản)
châu Á--Thái Bình Dương tạo thành 103 cặp quốc gia-năm qua 14 mốc khảo
sát.

Đồng thời, châu Á đang trải qua chuyển đổi số chưa từng có trong giai
đoạn 2018--2025. Sự bùng nổ của hạ tầng số và trí tuệ nhân tạo từ năm
2023 đã tái định hình cơ chế qua đó năng lực số tác động lên quốc tế hóa
và hiệu quả (Bharadwaj et al., 2013; Banalieva \& Dhanaraj, 2019;
Verhoef et al., 2021; Yang, Zhao \& Wei, 2025).

Khung lý thuyết hiện hành về quan hệ quốc tế hóa và hiệu quả còn phân
tán. Các tổng quan lớn thường tập trung vào một hoặc hai khía cạnh mà ít
có khung tích hợp đa tầng cho phép giải thích đầy đủ tính không đồng
nhất quan sát được ở châu Á. Đặc biệt, khung áp dụng số nền tảng
(Foundational Digital Adoption Framework) chỉ mới được đề xuất gần đây
và chưa được tích hợp đầy đủ với bốn tầng lý thuyết kinh điển.

Kết quả phân tích mô tả trong Chuyên đề Tiến sĩ số 1 đã cung cấp bức
tranh thực trạng đa chiều dựa trên 88.869 doanh nghiệp WBES xuyên 50 nền
kinh tế (gồm Nhật Bản), làm rõ ba khuôn mẫu then chốt: phân tán năng
suất tăng đơn điệu theo chế độ từ Nhóm I đến Nhóm VI; tính không đồng
nhất nội bộ trong nhóm Advanced; và khuôn mẫu adaptation đặc trưng của
các quốc đảo nhỏ Thái Bình Dương. Đồng thời, phân tích đó phát hiện
CDCM, Mô hình Áp dụng số và Năng lực phụ thuộc bối cảnh: tác động của
DAI lên hiệu quả phụ thuộc vào chế độ thể chế, giai đoạn quốc tế hóa, và
mức độ bão hòa số của nền kinh tế. Chuyên đề 2 này tiếp nối bằng cách
xây dựng mô hình lý thuyết và thực nghiệm để giải thích cơ chế đằng sau
các khuôn mẫu đó.

Bốn bản thảo thực nghiệm đồng hành cung cấp bằng chứng quốc gia cụ thể
kiểm định khả năng của khung đề xuất. P2 Trung Quốc SMEs (Đỗ \& Phan,
2026c, JFAR đã công bố) phân tích 4.290 quan sát doanh nghiệp và năm
thuộc các doanh nghiệp nhỏ và vừa Trung Quốc và tìm thấy dạng chữ U
ngược bậc hai với điểm uốn tối ưu 47,8\% FSTS, cung cấp tiền thân thực
nghiệm cho H1 và cơ sở để so sánh với P5. P3 Việt Nam (Đỗ \& Phan,
2026a, mục tiêu JED) phân tích 2.958 doanh nghiệp WBES ba sóng (2009,
2015, 2023) xác nhận chữ U ngược trong cả ba sóng với điểm uốn 46,2\%
(2009), 39,3\% (2015), 41,6\% (2023), 39,7\% (tổng hợp); TCI nhân quả
qua biến công cụ (β = 1,639, F bậc một = 22,1); DAI theo giai đoạn và
phụ thuộc chọn lựa (IV null, β = 0,018). P4 Singapore (Mar, Đỗ \& Phan,
2026, mục tiêu MIR) phân tích 623 doanh nghiệp WBES 2023, cho thấy điểm
uốn hàm ý ở FSTS khoảng 88,6\% và DAI hoạt động như nguồn lực mở rộng
tình huống với hệ số DAI×FSTS² = +3,119 (p = 0,005). P5 Trung Quốc (Đỗ
\& Phan, 2026b, mục tiêu IJOEM) phân tích 4.544 doanh nghiệp WBES hai
sóng (2012, 2024) xác nhận chữ U ngược bền vững cấu trúc với điểm uốn
49,4\% (2012) và 47,2\% (2024); kết quả này nhất quán với bằng chứng
trước đó về quan hệ quốc tế hóa và hiệu quả và hiệu ứng vùng miền ở
doanh nghiệp Trung Quốc (Chen \& Tan, 2012; Xiao et al., 2013).

\hypertarget{mux1ee5c-tiuxeau}{%
\paragraph{2.1.2 Mục tiêu}\label{mux1ee5c-tiuxeau}}

Mục tiêu chung: xây dựng mô hình nghiên cứu lý thuyết và thực nghiệm về
ảnh hưởng của quốc tế hóa đến hiệu quả hoạt động kinh doanh của các
doanh nghiệp ở các quốc gia châu Á, có khả năng kiểm định bằng dữ liệu
WBES 88.869 doanh nghiệp xuyên 50 nền kinh tế (gồm Nhật Bản) giai đoạn
2009--2025.

Các mục tiêu cụ thể:

\begin{enumerate}
\def\labelenumi{\arabic{enumi}.}
\item
  Hệ thống hóa khung lý thuyết tích hợp 4 tầng (Uppsala, RBV,
  Institutional Theory, Upper Echelons) kết hợp Foundational Digital
  Adoption Framework và CDCM.
\item
  Lược khảo các mô hình quốc tế hóa và hiệu quả trên thế giới và châu Á,
  định vị khoảng trống nghiên cứu.
\item
  Phát triển hệ giả thuyết H1--H6 dựa trên khung tích hợp, có neo đậu từ
  bằng chứng P3, P4, P5.
\item
  Mô hình tám mô hình thực nghiệm M0--M7 với khả năng kiểm định trên dữ
  liệu WBES.
\item
  Đề xuất chiến lược nhận dạng và kiểm định độ vững toàn diện.
\item
  Khẳng định tính mới về mô hình so với các khung tham chiếu hiện hành.
\end{enumerate}

\hypertarget{nux1ed9i-dung}{%
\paragraph{2.1.3 Nội dung}\label{nux1ed9i-dung}}

Chuyên đề trả lời ba câu hỏi nghiên cứu chính:

Câu hỏi 1 (Q1): Quan hệ giữa cường độ quốc tế hóa (I) và hiệu quả hoạt
động kinh doanh (P) ở các doanh nghiệp châu Á có dạng hàm gì? Tuyến
tính, phi tuyến hay phụ thuộc vào ngữ cảnh?

Câu hỏi 2 (Q2): Bốn cơ chế điều tiết, năng lực công nghệ TCI, mức độ áp
dụng số DAI, đặc điểm nhà quản trị, chế độ thể chế nội địa, tác động như
thế nào lên quan hệ quốc tế hóa và hiệu quả? Có phân biệt được giữa TCI
và DAI theo CDCM hay không?

Câu hỏi 3 (Q3): Tác động quốc tế hóa và hiệu quả và cơ chế điều tiết có
thay đổi theo thời gian không, đặc biệt giữa giai đoạn 2009--2012 và
2018--2025?

\hypertarget{giux1edbi-hux1ea1n-cux1ee7a-chuyuxean-ux111ux1ec1}{%
\paragraph{2.1.4 Giới hạn của chuyên
đề}\label{giux1edbi-hux1ea1n-cux1ee7a-chuyuxean-ux111ux1ec1}}

Đối tượng nghiên cứu: doanh nghiệp đang hoạt động ở các nền kinh tế châu
Á và Thái Bình Dương đáp ứng tiêu chí mẫu của WBES.

Phạm vi không gian: 50 nền kinh tế (gồm Nhật Bản) phân theo sáu nhóm
ICRV (6+6+6+7+17+8).

Phạm vi thời gian: 2009--2026 (kèm đợt khảo sát Oman 2026 trong bộ dữ
liệu hợp nhất), bao trùm ba thế hệ lược đồ WBES, tạo thành 103 cặp nền
kinh tế--năm trên khung phân tích 50 nền (gồm Nhật Bản).

Phạm vi nội dung: mô hình hóa quan hệ quốc tế hóa và hiệu quả và bốn cơ
chế điều tiết cộng với tính không đồng nhất theo thời gian. Chuyên đề
không trình bày kết quả thực nghiệm từ toàn bộ mẫu gộp, đó là nhiệm vụ
của các bài báo tiếp theo trong chuỗi luận án. Chuyên đề này tập trung
xây dựng và thiết lập mô hình.

\hypertarget{uxfd-nghux129a}{%
\paragraph{2.1.5 Ý nghĩa}\label{uxfd-nghux129a}}

Về lý thuyết, chuyên đề đóng góp khung tích hợp 4 tầng kết hợp Khung
chấp nhận số nền tảng, chưa có khung tương đương trong tài liệu IB cho
khu vực 50 nền kinh tế (gồm Nhật Bản) bao gồm cả quốc đảo nhỏ Thái Bình
Dương. CDCM làm sắc nét sự phân biệt giữa TCI như nguồn lực nâng mặt
bằng và DAI như nguồn lực mở rộng tình huống.

Về mô hình, tám mô hình M0--M7 với điều tiết ba chiều (M7) là chuỗi mô
hình đầu tiên trong tài liệu châu Á kiểm định đồng thời phi tuyến quốc
tế hóa và hiệu quả và điều tiết đa tầng trên mẫu gộp vi mô lớn nhất từ
trước đến nay.

Về thực tiễn, kết quả cung cấp hàm ý chính sách cho doanh nghiệp Việt
Nam và châu Á mới nổi về cách tận dụng quốc tế hóa kết hợp năng lực công
nghệ và số hóa, đặc biệt trong giai đoạn AI bùng nổ.

\begin{center}\rule{0.5\linewidth}{0.5pt}\end{center}

\hypertarget{phux1b0ux1a1ng-phuxe1p-nghiuxean-cux1ee9u}{%
\subsubsection{2.2 Phương pháp nghiên
cứu}\label{phux1b0ux1a1ng-phuxe1p-nghiuxean-cux1ee9u}}

Chuyên đề tiếp cận theo lối tổng quan tài liệu tích hợp (Torraco, 2005)
cho khung lý thuyết và phát triển giả thuyết suy diễn (deductive
hypothesis development) cho hệ giả thuyết. Thiết lập mô hình thực nghiệm
tuân theo Wooldridge (2010) và Long \& Ervin (2000). Thiết kế nhận dạng
và kiểm định độ vững theo Greene (2018). Hệ giả thuyết được neo đậu vào
bằng chứng thực nghiệm từ P3, P4, P5 theo nguyên tắc nghiên cứu dựa trên
hiện tượng (Meyer et al., 2017).

Phương pháp ước lượng chính: OLS với HC1 sai số chuẩn vững (Long \&
Ervin, 2000) làm chuẩn, hai chiều hiệu ứng cố định (country × year), và
biến công cụ (2SLS) để nhận dạng nhân quả cho TCI và DAI. Kiểm định phi
tuyến áp dụng Lind--Mehlum (2010) U-test để xác nhận chữ U ngược thực sự
và kiểm định Paternoster để so sánh điểm uốn giữa các nhóm ICRV và giai
đoạn thời gian. Phần mềm: Stata 18 (OLS, IV, margins) kết hợp R (LOWESS,
ggplot2, CI điểm uốn).

\begin{center}\rule{0.5\linewidth}{0.5pt}\end{center}

\hypertarget{kux1ebft-quux1ea3-vuxe0-thux1ea3o-luux1eadn}{%
\subsubsection{2.3 Kết quả và thảo
luận}\label{kux1ebft-quux1ea3-vuxe0-thux1ea3o-luux1eadn}}

\hypertarget{tux1ed5ng-quan-tuxe0i-liux1ec7u-vux1ec1-quan-hux1ec7-quux1ed1c-tux1ebf-huxf3a-vuxe0-hiux1ec7u-quux1ea3-kinh-doanh}{%
\paragraph{2.3.1 Tổng quan tài liệu về quan hệ quốc tế hóa và hiệu quả
kinh
doanh}\label{tux1ed5ng-quan-tuxe0i-liux1ec7u-vux1ec1-quan-hux1ec7-quux1ed1c-tux1ebf-huxf3a-vuxe0-hiux1ec7u-quux1ea3-kinh-doanh}}

\hypertarget{nux103m-dux1ea1ng-huxe0m-cux1ee7a-quan-hux1ec7-quux1ed1c-tux1ebf-huxf3a-vuxe0-hiux1ec7u-quux1ea3}{%
\subparagraph{Năm dạng hàm của quan hệ quốc tế hóa và hiệu
quả}\label{nux103m-dux1ea1ng-huxe0m-cux1ee7a-quan-hux1ec7-quux1ed1c-tux1ebf-huxf3a-vuxe0-hiux1ec7u-quux1ea3}}

Tài liệu IB đã ghi nhận không dưới năm dạng hàm khác nhau trong quan hệ
giữa cường độ quốc tế hóa và hiệu quả hoạt động kinh doanh của doanh
nghiệp. Sự đa dạng này không phải ngẫu nhiên mà phản ánh tính phụ thuộc
ngữ cảnh của hiệu ứng quốc tế hóa.

Dạng hàm thứ nhất, tuyến tính. Các nghiên cứu sớm như Hsu và Boggs
(2003) và Geringer và cộng sự (1989) tìm thấy mối quan hệ tuyến tính
dương: hiệu quả tăng đơn điệu theo mức độ quốc tế hóa. Logic nền tảng là
lợi ích quy mô, đa dạng hóa doanh thu, và tiếp cận thị trường lớn hơn.
Đây là kết quả đặc thù của các mẫu với FSTS thấp hoặc phân tán hẹp.

Dạng hàm thứ hai, chữ U ngược. Hitt, Hoskisson và Kim (1997) và Gomes và
Ramaswamy (1999) đề xuất mô hình chữ U ngược: hiệu quả tăng ở giai đoạn
quốc tế hóa ban đầu đến trung bình do lợi ích quy mô và học hỏi, nhưng
giảm sau điểm uốn tối ưu do chi phí điều phối đa thị trường vượt lợi
ích. Điểm uốn thực nghiệm thường nằm trong khoảng 30--60\% FSTS (Marano
et al., 2016). Đây là dạng hàm phổ biến nhất trong các tổng quan định
lượng.

Dạng hàm thứ ba, đường cong chữ S ba giai đoạn (S-curve). Lu và Beamish
(2004) và Contractor, Kundu và Hsu (2003) phát triển lý thuyết ba giai
đoạn: giai đoạn đầu hiệu quả giảm do chi phí học tập và thiết lập; giai
đoạn giữa hiệu quả tăng do lợi ích quy mô và đa dạng hóa; giai đoạn quá
mức hiệu quả giảm do chi phí phối hợp. Dạng chữ S có hệ số β₁(FSTS)
\textgreater{} 0, β₂(FSTS²) \textless{} 0, β₃(FSTS³) \textgreater{} 0.
Bằng chứng hỗ trợ mạnh từ các MNE lớn, nhưng kém ổn định hơn với doanh
nghiệp vừa và nhỏ.

Dạng hàm thứ tư, đường cong chữ M (M-curve). Riahi-Belkaoui (1998) đề
xuất dạng chữ M với hai điểm uốn, phản ánh tính không đồng nhất trong
lợi ích quốc tế hóa theo loại thị trường. Bằng chứng thực nghiệm hạn chế
và khó tái lập.

Dạng hàm thứ năm, gánh nặng bắt buộc (forced penalty). Glaum và Oesterle
(2007) ghi nhận dạng hàm tuyến tính âm trong các nền kinh tế bắt buộc
phải quốc tế hóa do thị trường nội địa quá nhỏ (Briguglio, 1995;
Bertram, 2006). Đây là khuôn mẫu đặc thù của các quốc đảo nhỏ Thái Bình
Dương: các doanh nghiệp phải xuất khẩu để tồn tại nhưng thiếu năng lực
cạnh tranh, tạo ra gánh nặng bắt buộc, hiệu quả không tăng theo cường độ
xuất khẩu.

Tất cả năm dạng hàm có thể được thống nhất dưới khung chi phí-lợi ích
(Contractor, 2012): P(FSTS) = B(FSTS) − C(FSTS), trong đó B(·) là hàm
lợi ích và C(·) là hàm chi phí. Sự khác biệt giữa các dạng hàm nằm ở
hình dạng của B(·) và C(·) tùy theo ngữ cảnh thể chế, năng lực doanh
nghiệp, và giai đoạn số hóa.

Bảng 2.2. Năm dạng hàm quốc tế hóa và hiệu quả, cơ chế, và bối cảnh xuất
hiện.

\begin{longtable}[]{@{}
  >{\raggedright\arraybackslash}p{(\columnwidth - 6\tabcolsep) * \real{0.1923}}
  >{\raggedright\arraybackslash}p{(\columnwidth - 6\tabcolsep) * \real{0.3654}}
  >{\raggedright\arraybackslash}p{(\columnwidth - 6\tabcolsep) * \real{0.1538}}
  >{\raggedright\arraybackslash}p{(\columnwidth - 6\tabcolsep) * \real{0.2885}}@{}}
\toprule\noalign{}
\begin{minipage}[b]{\linewidth}\raggedright
Dạng hàm
\end{minipage} & \begin{minipage}[b]{\linewidth}\raggedright
Tác giả tiêu biểu
\end{minipage} & \begin{minipage}[b]{\linewidth}\raggedright
Cơ chế
\end{minipage} & \begin{minipage}[b]{\linewidth}\raggedright
Bối cảnh ICRV
\end{minipage} \\
\midrule\noalign{}
\endhead
\bottomrule\noalign{}
\endlastfoot
Tuyến tính (+) & Hsu \& Boggs (2003) & Quy mô + học hỏi & FSTS thấp, tất
cả nhóm \\
Chữ U ngược & Hitt et al.~(1997); Gomes \& Ramaswamy (1999) & Chi phí
điều phối \textgreater{} lợi ích ở FSTS cao & Nhóm III--IV \\
Chữ S ba giai đoạn (S-curve) & Lu \& Beamish (2004); Contractor et
al.~(2003) & Ba giai đoạn: học tập, quy mô, quá mức & MNE lớn, tất cả
nhóm \\
Chữ M (M-curve) & Riahi-Belkaoui (1998) & Lợi ích không đồng nhất & Bằng
chứng hạn chế \\
Forced penalty & Glaum \& Oesterle (2007) & Bắt buộc quốc tế hóa, thiếu
năng lực & Nhóm VI (SIDS) \\
\end{longtable}

\hypertarget{nux103m-tux1ed5ng-quan-ux111ux1ecbnh-lux1b0ux1ee3ng-lux1edbn-19802024}{%
\subparagraph{Năm tổng quan định lượng lớn
(1980--2024)}\label{nux103m-tux1ed5ng-quan-ux111ux1ecbnh-lux1b0ux1ee3ng-lux1edbn-19802024}}

Năm tổng quan định lượng lớn cung cấp toàn cảnh về quan hệ quốc tế hóa
và hiệu quả.

Bausch và Krist (2007) phân tích 36 nghiên cứu (41 mẫu, 1980--2005) và
tìm thấy hiệu ứng trung bình dương nhỏ, nhưng với biến động cao (SD =
0,21), cho thấy các biến điều tiết quan trọng hơn hiệu ứng trung bình.
Phương pháp: hồi quy phân tích tổng hợp với 8 biến điều tiết gồm quy mô
doanh nghiệp, cường độ R\&D, loại ngành, GDP quốc gia, khoảng cách văn
hóa, phạm vi địa lý, và hai thước đo hiệu quả. Phát hiện chính: không
tồn tại một hiệu ứng trung bình đơn nhất; các biến điều tiết thể chế và
văn hóa giải thích lượng phương sai lớn hơn đáng kể so với biến điều
tiết cấp doanh nghiệp. Hạn chế: thiếu đại diện châu Á (chỉ 8/36 nghiên
cứu), không kiểm định phi tuyến có hệ thống.

Kirca và cộng sự (2012) tổng hợp 152 mẫu độc lập từ 141 nghiên cứu, tìm
thấy mối quan hệ dương trung bình nhưng với điều tiết đáng kể từ loại
hình quốc tế hóa, thước đo hiệu quả, và thuộc tính doanh nghiệp. Cường
độ R\&D là biến điều tiết cấp doanh nghiệp mạnh nhất (β\_meta = +0,12, p
\textless{} 0,001). Hạn chế: thiếu bằng chứng từ châu Á và thị trường
mới nổi; không có biến số; không bao gồm SIDS hay nền kinh tế cận biên.

Marano và cộng sự (2016) phân tích 359 nghiên cứu với cách tiếp cận tập
trung vào thể chế. Phát hiện chính: chất lượng thể chế quốc gia xuất xứ
có thể dịch chuyển điểm uốn ±15\% FSTS, đây là nền tảng quan trọng nhất
ủng hộ phổ thể chế ICRV trong H5. Hạn chế: không phân tách TCI và DAI;
không có dữ liệu sau 2015; không bao gồm quốc đảo nhỏ. Bổ sung cho
Marano, Beugelsdijk và cộng sự (2018) tổng hợp bằng chứng về khoảng cách
văn hóa trong quá trình quốc tế hóa và cho thấy hệ quả của quốc tế hóa
phụ thuộc một cách hệ thống vào khoảng cách văn hóa và thể chế giữa thị
trường nguồn và đích; phát hiện này bổ trợ cho phổ thể chế ICRV của CĐ2
ở chỗ khoảng cách và chất lượng thể chế là hai chiều bối cảnh khác nhau
cùng định hình quan hệ, không trùng lặp.

Wu, Fan và Chen (2022) tổng hợp 20 năm bằng chứng từ doanh nghiệp đa
quốc gia thị trường mới nổi, tìm thấy điểm uốn thấp hơn so với doanh
nghiệp từ nền kinh tế phát triển khoảng 12--18\% FSTS do khoảng cách
năng lực hấp thụ. Hạn chế: không kiểm định giai đoạn số hóa; không sử
dụng dữ liệu vi mô WBES; không phân biệt TCI và DAI.

Schwens và cộng sự (2018) tập trung vào doanh nghiệp vừa và nhỏ và tìm
thấy dạng U (không phải chữ U ngược) trong một số bối cảnh, chi phí quốc
tế hóa ban đầu lớn so với năng lực, nhưng phục hồi sau khi vượt ngưỡng
thiết lập. Hạn chế: mẫu nhỏ (87 nghiên cứu); không bao phủ châu Á Thái
Bình Dương đầy đủ; không kiểm định DAI như biến điều tiết.

Kết luận tổng quan: hiệu ứng trung bình quốc tế hóa và hiệu quả là dương
nhỏ (r khoảng 0,04--0,07) nhưng rất không đồng nhất (SD = 0,15--0,21).
Cả năm tổng quan đều thiếu: (i) bao phủ đầy đủ châu Á và Thái Bình
Dương; (ii) kiểm định đồng thời phi tuyến và điều tiết đa tầng; (iii)
phân tách TCI và DAI; (iv) bao gồm quốc đảo nhỏ Thái Bình Dương; (v) dữ
liệu sau 2020. Chuyên đề này lấp đầy tất cả năm khoảng trống với mẫu gộp
88.869 doanh nghiệp (50 nền, gồm Nhật Bản) và khung CDCM.

Bảng 2.13. So sánh đặc điểm phương pháp và phạm vi lý thuyết của các
khung tham chiếu chính và CĐ2.

\begin{longtable}[]{@{}
  >{\raggedright\arraybackslash}p{(\columnwidth - 16\tabcolsep) * \real{0.0795}}
  >{\raggedright\arraybackslash}p{(\columnwidth - 16\tabcolsep) * \real{0.1591}}
  >{\raggedright\arraybackslash}p{(\columnwidth - 16\tabcolsep) * \real{0.0568}}
  >{\raggedright\arraybackslash}p{(\columnwidth - 16\tabcolsep) * \real{0.0568}}
  >{\raggedright\arraybackslash}p{(\columnwidth - 16\tabcolsep) * \real{0.1477}}
  >{\raggedright\arraybackslash}p{(\columnwidth - 16\tabcolsep) * \real{0.1932}}
  >{\raggedright\arraybackslash}p{(\columnwidth - 16\tabcolsep) * \real{0.1136}}
  >{\raggedright\arraybackslash}p{(\columnwidth - 16\tabcolsep) * \real{0.1250}}
  >{\raggedright\arraybackslash}p{(\columnwidth - 16\tabcolsep) * \real{0.0682}}@{}}
\toprule\noalign{}
\begin{minipage}[b]{\linewidth}\raggedright
Khung
\end{minipage} & \begin{minipage}[b]{\linewidth}\raggedright
Phi tuyến I--P
\end{minipage} & \begin{minipage}[b]{\linewidth}\raggedright
TCI
\end{minipage} & \begin{minipage}[b]{\linewidth}\raggedright
DAI
\end{minipage} & \begin{minipage}[b]{\linewidth}\raggedright
Đặc điểm QL
\end{minipage} & \begin{minipage}[b]{\linewidth}\raggedright
Thể chế (Inst.)
\end{minipage} & \begin{minipage}[b]{\linewidth}\raggedright
Temporal
\end{minipage} & \begin{minipage}[b]{\linewidth}\raggedright
Phân nhóm
\end{minipage} & \begin{minipage}[b]{\linewidth}\raggedright
SIDS
\end{minipage} \\
\midrule\noalign{}
\endhead
\bottomrule\noalign{}
\endlastfoot
Bausch \& Krist (2007) & Không (không kiểm định hệ thống) & Một phần
(R\&D proxy) & Không & Không & Có (GDP, văn hóa) & Không & Không &
Không \\
Kirca et al.~(2012) & Không & Có (R\&D proxy) & Không & Không & Một phần
& Không & Một phần (ngành) & Không \\
Marano et al.~(2016) & Có (TP ±15\% FSTS) & Không & Không & Không & Có
(cốt lõi) & Không & Không & Không \\
Banalieva \& Dhanaraj (2019) & Một phần & Không & Có (cốt lõi) & Không &
Một phần & Một phần & Không & Không \\
Wu, Fan \& Chen (2022) & Có (TP thấp hơn EMNE) & Một phần (absorptive
capacity) & Không & Không & Có & Có (20 năm) & Có (EMNE/DMNE) & Không \\
Đỗ \& Phan (2026c, JFAR) & Có (chữ U ngược bậc hai) & -- & -- & -- & --
& -- & -- & -- \\
\textbf{CĐ2 (luận án này)} & \textbf{Có (H1: chữ U ngược + chữ S)} &
\textbf{Có (H2: nhân quả IV)} & \textbf{Có (H3: CDCM phân biệt)} &
\textbf{Có (H4a--H4c)} & \textbf{Có (H5: ICRV 6 nhóm)} & \textbf{Có (H6:
3 giai đoạn)} & \textbf{Có (Nhóm I--VI)} & \textbf{Có (Nhóm VI FIP)} \\
\end{longtable}

\emph{Ghi chú: ``Có'' = kiểm định chính thức hoặc đóng góp lý thuyết rõ
ràng; Một phần = đề cập nhưng không phải trọng tâm; Không = không bao
gồm; -- = không áp dụng. FIP = Forced Internationalization Penalty. TCI
và DAI được phân tách theo tiêu chí CDCM; P2 JFAR không kiểm định CDCM
(pre-framework) nhưng cung cấp bằng chứng phi tuyến cubic tiền thân.}

\hypertarget{bux1eb1ng-chux1ee9ng-chuxe2u-uxe1-tux1eeb-ba-bux1ea3n-thux1ea3o-ux111ux1ed3ng-huxe0nh}{%
\subparagraph{Bằng chứng châu Á từ ba bản thảo đồng
hành}\label{bux1eb1ng-chux1ee9ng-chuxe2u-uxe1-tux1eeb-ba-bux1ea3n-thux1ea3o-ux111ux1ed3ng-huxe0nh}}

P3 Việt Nam (Đỗ \& Phan, 2026a, mục tiêu JED): Phân tích 2.958 doanh
nghiệp WBES ba sóng (2009, 2015, 2023) xác nhận chữ U ngược trong cả ba
sóng với kiểm định Lind--Mehlum p \textless{} 0,001 (tổng hợp). Điểm
uốn: 46,2\% (2009), 39,3\% (2015), 41,6\% (2023), 39,7\% (tổng hợp). P3
phân tách H1 thành hai mệnh đề: bước nhảy từ không xuất khẩu sang xuất
khẩu là biên năng suất chính (nhất quán với mô hình chọn lọc năng suất
xuất khẩu của Melitz, 2003); trong mẫu chỉ gồm doanh nghiệp xuất khẩu,
đường cong gần phẳng. TCI (b8+e6) dương bền vững (β = 0,179 tổng hợp, p
\textless{} 0,001); biến công cụ cho TCI nhân quả (β = 1,639, F bậc một
= 22,1). DAI phụ thuộc giai đoạn, dương năm 2009 (β = 0,175), null năm
2015, có tương tác âm năm 2023 (FSTS\_c×DAI\_z = −0,912, p = 0,043);
biến công cụ cho DAI null (β = 0,018, p = 0,942, F bậc một = 34,6), phụ
thuộc chọn lựa, không nhân quả.

\emph{Lưu ý về thiết lập mô hình P3: P3 Vietnam sử dụng chuỗi M0--M8
riêng, trong đó M8 là mô hình đầy đủ:
\(\ln(\mathrm{LP})_{it} = \alpha + \beta_1 \mathrm{FSTS}_{c} + \beta_2 \mathrm{FSTS}_{c}^{2} + \beta_3 \mathrm{TCI}_{z} + \beta_4 \mathrm{DAI}_{z} + \beta_5(\mathrm{FSTS}_{c} \times \mathrm{DAI}_{z}) + \beta_6(\mathrm{FSTS}_{c}^{2} \times \mathrm{DAI}_{z}) + \boldsymbol{\gamma}\mathbf{X} + \delta_s + [\lambda_t] + \varepsilon_{it}\).
Mô hình đầy đủ M0--M8 và bảng định nghĩa biến được trình bày trong bản
thảo P3 (Mục 3.2--Mục 3.3) và luận án (Mục 4.5.1). Chuỗi M0--M7 trong
Chuyên đề 2 là các mô hình cấp độ mẫu gộp 50 nền kinh tế, khác với mô
hình cấp độ quốc gia của P3.}

P4 Singapore (Mar, Đỗ \& Phan, 2026, mục tiêu MIR): Phân tích 623 doanh
nghiệp WBES 2023. Trong nền kinh tế số bão hòa, đường cong quốc tế hóa
và hiệu quả chủ yếu dương, điểm uốn hàm ý ở FSTS khoảng 88,6\% (vùng
thưa dữ liệu; kiểm định Lind--Mehlum p = 0,303). DAI×FSTS² = +3,119 (p =
0,005), DAI là nguồn lực mở rộng tình huống, chỉ phát huy ở FSTS cao nơi
nhu cầu điều phối xuyên biên giới dày đặc. TCI dương trực tiếp (β =
0,153).

P5 Trung Quốc (Đỗ \& Phan, 2026b, mục tiêu IJOEM): Phân tích 4.544 doanh
nghiệp WBES hai sóng (2012, 2024). Chữ U ngược bền vững cấu trúc, điểm
uốn 49,4\% (2012), 47,2\% (2024), kiểm định Paternoster cho hệ số FSTS²
không có sự khác biệt đáng kể (p = 0,545). TCI tăng cường theo thời gian
(+0,260 sang +0,426, Paternoster p = 0,011).

Bảng 2.3. So sánh bằng chứng bốn bản thảo đồng hành.

\begin{longtable}[]{@{}
  >{\raggedright\arraybackslash}p{(\columnwidth - 8\tabcolsep) * \real{0.0648}}
  >{\raggedright\arraybackslash}p{(\columnwidth - 8\tabcolsep) * \real{0.2870}}
  >{\raggedright\arraybackslash}p{(\columnwidth - 8\tabcolsep) * \real{0.2037}}
  >{\raggedright\arraybackslash}p{(\columnwidth - 8\tabcolsep) * \real{0.2037}}
  >{\raggedright\arraybackslash}p{(\columnwidth - 8\tabcolsep) * \real{0.2407}}@{}}
\toprule\noalign{}
\begin{minipage}[b]{\linewidth}\raggedright
Chiều
\end{minipage} & \begin{minipage}[b]{\linewidth}\raggedright
P2 Trung Quốc SMEs (Nhóm III)
\end{minipage} & \begin{minipage}[b]{\linewidth}\raggedright
P3 Việt Nam (Nhóm IV)
\end{minipage} & \begin{minipage}[b]{\linewidth}\raggedright
P4 Singapore (Nhóm I)
\end{minipage} & \begin{minipage}[b]{\linewidth}\raggedright
P5 Trung Quốc (Nhóm III)
\end{minipage} \\
\midrule\noalign{}
\endhead
\bottomrule\noalign{}
\endlastfoot
Journal & JFAR (đã công bố 2026) & JED (mục tiêu) & MIR (mục tiêu) &
IJOEM (mục tiêu) \\
Dạng I--P & \textbf{Chữ U ngược bậc hai} (điểm uốn 47,8\%) & Chữ U ngược
xác nhận & Chủ yếu dương, TP \textasciitilde88,6\% & Chữ U ngược bền
vững \\
Điểm uốn & 47,8\% FSTS & 39--46\% FSTS & \textasciitilde88,6\% (thưa dữ
liệu) & 47--49\% FSTS \\
TCI & Dương; RBV xác nhận & Dương bền vững, nhân quả (IV) & Dương trực
tiếp & Dương, tăng theo thời gian \\
DAI & Không kiểm định (pre-CDCM) & Phụ thuộc giai đoạn; IV null & Mở
rộng tình huống (+3,119) & Kiểm soát (Bậc 1 bão hòa) \\
Temporal & Cắt ngang (không temporal) & Paternoster p \textless{} 0,001
(dịch chuyển) & Cắt ngang 2023 & Paternoster p = 0,545 (ổn định) \\
Mô hình CDCM & Tiền thân (chưa tách TCI/DAI) & Nhóm IV, trung gian &
Nhóm I, bão hòa số & Nhóm III, đang trưởng thành \\
\end{longtable}

\emph{Ghi chú: P2 (Đỗ \& Phan, 2026c, JFAR) là bản thảo tiền thân đã
công bố, tìm thấy quan hệ chữ U ngược bậc hai (điểm uốn 47,8\% FSTS, mẫu
4.290 quan sát doanh nghiệp--năm) trong nhóm doanh nghiệp nhỏ và vừa
Trung Quốc, là bằng chứng nền tảng cho dạng phi tuyến trong H1. P2 không
kiểm định CDCM, được đưa vào đây để làm rõ sự tiến hóa phương pháp từ P2
sang P5.}

Bốn bản thảo cùng xác nhận ba dự đoán cốt lõi của CDCM: TCI là nhân quả
và bền vững qua ba nền kinh tế khác nhau; DAI là phụ thuộc ngữ cảnh,
không đồng nhất; chữ U ngược bậc hai là dạng hàm phổ biến nhất cho doanh
nghiệp vừa và nhỏ châu Á với điểm uốn phụ thuộc vào ICRV.

\hypertarget{khung-luxfd-thuyux1ebft-tuxedch-hux1ee3p-4-tux1ea7ng}{%
\paragraph{2.3.2 Khung lý thuyết tích hợp 4
tầng}\label{khung-luxfd-thuyux1ebft-tuxedch-hux1ee3p-4-tux1ea7ng}}

\hypertarget{tux1ea7ng-1-luxfd-thuyux1ebft-quux1ed1c-tux1ebf-huxf3a-uppsala}{%
\subparagraph{Tầng 1: Lý thuyết quốc tế hóa
Uppsala}\label{tux1ea7ng-1-luxfd-thuyux1ebft-quux1ed1c-tux1ebf-huxf3a-uppsala}}

Mô hình Uppsala (Johanson \& Vahlne, 1977) đề xuất quá trình quốc tế hóa
của doanh nghiệp diễn ra tăng dần thông qua tích lũy kinh nghiệm và giảm
dần khoảng cách tâm lý giữa thị trường nội địa và thị trường nước ngoài.
Doanh nghiệp khởi đầu xuất khẩu sang các thị trường gần về văn hóa, ngôn
ngữ, thể chế trước khi mở rộng sang các thị trường xa hơn. Mô hình ban
đầu (1977) tập trung vào ba cơ chế: cam kết thị trường, tri thức thị
trường, và quyết định cam kết hiện tại.

Johanson và Vahlne (2009) tái định vị mô hình từ gánh nặng do là người
nước ngoài sang gánh nặng do không thuộc mạng lưới. Phiên bản này nhấn
mạnh vai trò mạng lưới quan hệ kinh doanh trong quá trình quốc tế hóa.

Mô hình Uppsala chậm áp dụng cho ba hiện tượng: doanh nghiệp quốc tế hóa
ngay từ khi thành lập (Knight \& Cavusgil, 2004); doanh nghiệp đa quốc
gia từ thị trường mới nổi với quy trình quốc tế hóa khác biệt (Mathews,
2002; Luo \& Tung, 2007); và bối cảnh số nơi Internet và nền tảng số
giảm khoảng cách tâm lý đến mức tối thiểu (Banalieva \& Dhanaraj, 2019;
Stallkamp \& Schotter, 2021).

Vai trò trong mô hình đề xuất: Uppsala cung cấp logic cho hiệu ứng phi
tuyến trong H1. Doanh nghiệp tăng cường độ xuất khẩu đối mặt với chi phí
học tập ban đầu, sau đó thu được lợi ích quy mô, nhưng vượt quá ngưỡng
nhất định lại chịu chi phí phối hợp. Bằng chứng P3 Việt Nam (điểm uốn
39--46\%) và P5 Trung Quốc (điểm uốn 47--49\%) neo đậu thực nghiệm cho
H1.

\hypertarget{tux1ea7ng-2-luxfd-thuyux1ebft-doanh-nghiux1ec7p-dux1ef1a-truxean-nguux1ed3n-lux1ef1c-rbv}{%
\subparagraph{Tầng 2: Lý thuyết doanh nghiệp dựa trên nguồn lực
(RBV)}\label{tux1ea7ng-2-luxfd-thuyux1ebft-doanh-nghiux1ec7p-dux1ef1a-truxean-nguux1ed3n-lux1ef1c-rbv}}

Lý thuyết dựa trên nguồn lực (Wernerfelt, 1984; Barney, 1991) cho rằng
lợi thế cạnh tranh bền vững xuất phát từ nguồn lực có giá trị, hiếm, khó
bắt chước, và không thay thế (VRIN). Doanh nghiệp khác nhau về hiệu quả
vì sở hữu các bộ nguồn lực khác nhau. Các phát triển gần đây mở rộng lý
thuyết dựa trên nguồn lực sang lợi thế cạnh tranh bền vững gắn với bối
cảnh môi trường, xã hội và quản trị (Bhandari, Ranta \& Salo, 2023).

Teece, Pisano và Shuen (1997) phát triển khái niệm năng lực động, khả
năng tích hợp, xây dựng, và tái cấu trúc nguồn lực để thích ứng với môi
trường thay đổi. Trong bối cảnh quốc tế hóa, năng lực động đặc biệt quan
trọng để hấp thụ tri thức từ thị trường nước ngoài và tái phân bổ nguồn
lực qua các thị trường.

Cohen và Levinthal (1990) đề xuất khái niệm năng lực hấp thụ, khả năng
doanh nghiệp nhận diện, đồng hóa, và áp dụng kiến thức bên ngoài. Doanh
nghiệp có R\&D nội tại cao có năng lực hấp thụ cao và do đó hấp thụ tri
thức từ quốc tế hóa hiệu quả hơn.

Vai trò trong mô hình đề xuất: Lý thuyết dựa trên nguồn lực cung cấp lý
thuyết cho H2 (TCI điều tiết tích cực). Doanh nghiệp có năng lực công
nghệ (R\&D, ISO, máy nhập khẩu) cao hấp thụ lợi ích quốc tế hóa tốt hơn.
TCI là biến điều tiết cấp doanh nghiệp. Bằng chứng P3 Việt Nam (TCI trực
tiếp β = 0,179, nhân quả từ IV) và P5 Trung Quốc (TCI tăng cường +0,260
sang +0,426) ủng hộ TCI là nguồn lực VRIN bền vững.

\hypertarget{tux1ea7ng-3-luxfd-thuyux1ebft-thux1ec3-chux1ebf-institutional-theory}{%
\subparagraph{Tầng 3: Lý thuyết thể chế (Institutional
Theory)}\label{tux1ea7ng-3-luxfd-thuyux1ebft-thux1ec3-chux1ebf-institutional-theory}}

North (1990) định nghĩa thể chế là các quy tắc của trò chơi, chính thức
(luật pháp, hợp đồng, quyền sở hữu) và phi chính thức (chuẩn mực, văn
hóa, mạng lưới), quy định chi phí giao dịch và rủi ro kinh doanh. Kim,
Kumar, Ramalho và Russell (2026) bổ sung khung đo lường năng lực thể chế
(institutional capacity) theo hai chiều: chiều tổ chức (nhân sự, nguồn
lực tài chính, hệ thống thông tin, quản lý) và chiều quản trị (minh
bạch, độc lập, trách nhiệm giải trình). Nghiên cứu đa quốc gia này xác
nhận mối tương quan dương vững chắc giữa năng lực thể chế và phát triển
kinh tế, cung cấp nền tảng thực nghiệm cho phân loại sáu nhóm ICRV theo
năng lực và nguồn lực thể chế được sử dụng trong mô hình đề xuất.

Khanna và Palepu (2010) áp dụng lý thuyết thể chế cho thị trường mới
nổi, đề xuất khái niệm khoảng trống thể chế, các thiếu hụt trong hệ
thống thể chế (thiếu thị trường vốn phát triển, thiếu hệ thống tư pháp
hiệu quả, thiếu chuẩn kế toán minh bạch). Doanh nghiệp ở thị trường mới
nổi phải xây dựng giải pháp nội bộ để bù đắp.

Peng (2003) và Peng, Wang, Jiang (2008) đề xuất cặp chiến lược ba chân,
chiến lược doanh nghiệp được hình thành bởi ba lực: ngành, nguồn lực, và
thể chế. Trong các nền kinh tế đang nổi và cận biên, vai trò của thể chế
đặc biệt nổi bật.

Vai trò trong mô hình đề xuất: Lý thuyết thể chế cung cấp logic cho H5
(thể chế ICRV điều tiết theo phổ thể chế). Cùng một mức quốc tế hóa,
doanh nghiệp ở Nhóm I thu được hiệu quả khác doanh nghiệp ở Nhóm V do
khác biệt về chi phí giao dịch, chất lượng thể chế, và tiếp cận thị
trường. Sáu nhóm ICRV là biến điều tiết cấp quốc gia.

\hypertarget{tux1ea7ng-4-luxfd-thuyux1ebft-thux1b0ux1ee3ng-tux1ea7ng-quux1ea3n-trux1ecb}{%
\subparagraph{Tầng 4: Lý thuyết thượng tầng quản
trị}\label{tux1ea7ng-4-luxfd-thuyux1ebft-thux1b0ux1ee3ng-tux1ea7ng-quux1ea3n-trux1ecb}}

Hambrick và Mason (1984) đề xuất lý thuyết thượng tầng quản trị với luận
điểm trung tâm: các quyết định chiến lược của doanh nghiệp phản ánh đặc
điểm nhân khẩu học của nhà quản trị cấp cao. Các đặc điểm này gồm tuổi,
giáo dục, kinh nghiệm chức năng, thâm niên, kinh nghiệm quốc tế, giới
tính.

Hambrick (2007) bổ sung hai khái niệm quan trọng: mức độ tự do của nhà
quản trị trong việc đưa ra quyết định, và áp lực công việc ảnh hưởng đến
chất lượng quyết định. Cannella, Park và Lee (2008) cho thấy đa dạng
chức năng của đội quản lý cấp cao tương quan dương với hiệu quả ở các
môi trường thay đổi nhanh.

Vai trò trong mô hình đề xuất: Lý thuyết thượng tầng quản trị cung cấp
lý thuyết cho H4a--H4c (đặc điểm nhà quản trị điều tiết). Doanh nghiệp
có nhà quản trị kinh nghiệm cao (H4a), học vấn cao (H4b) chuyển hóa quốc
tế hóa thành hiệu quả tốt hơn thông qua hai kênh: chất lượng quyết định
(tránh quốc tế hóa sớm) và năng lực điều phối (quản lý phức tạp đa thị
trường). Giới tính nhà quản trị (H4c) được kiểm định khám phá không có
dấu tiên nghiệm.

\hypertarget{khung-chux1ea5p-nhux1eadn-sux1ed1-nux1ec1n-tux1ea3ng-vuxe0-cdcm}{%
\subparagraph{Khung chấp nhận số nền tảng và
CDCM}\label{khung-chux1ea5p-nhux1eadn-sux1ed1-nux1ec1n-tux1ea3ng-vuxe0-cdcm}}

Banalieva và Dhanaraj (2019) phát triển lý thuyết nội hóa cho nền kinh
tế số, đề xuất rằng nền tảng số làm thay đổi căn bản ba yếu tố cốt lõi
của lý thuyết IB: khoảng cách tâm lý giảm đến mức tối thiểu; chuỗi giá
trị phân mảnh ở quy mô chưa từng có; doanh nghiệp quốc tế hóa ngay từ
ngày đầu thành lập trở thành chuẩn mới.

Verhoef và cộng sự (2021) đề xuất phân biệt rõ giữa hai loại năng lực:
năng lực công nghệ nội tại (TCI) gồm R\&D doanh nghiệp, đổi mới sản
phẩm, chứng nhận chất lượng ISO, bằng sáng chế, đo lường bằng đầu vào và
đầu ra của hoạt động đổi mới (Bậc 3--4); và áp dụng số nền tảng (DAI)
gồm sử dụng hạ tầng số cơ bản (website, e-commerce, e-payment), đo lường
bằng Bậc 1--2.

Stallkamp và Schotter (2021) chứng minh rằng nền tảng số cho phép doanh
nghiệp nhỏ ở thị trường mới nổi quốc tế hóa với chi phí gần bằng zero,
thay đổi căn bản logic Uppsala.

CDCM, Mô hình Áp dụng số và Năng lực phụ thuộc bối cảnh, được phát hiện
qua phân tích đa mẫu trong Chuyên đề Tiến sĩ số 1 cho thấy tác động của
DAI lên hiệu quả không đồng nhất mà phụ thuộc đồng thời vào ba chiều ngữ
cảnh.

Chiều thứ nhất là chế độ thể chế (ICRV): ở Nhóm I nơi Bậc 1--2 đã bão
hòa, DAI không tạo phần thưởng đồng nhất mà chỉ phát huy qua kênh phụ
thuộc xuất khẩu (nghịch lý bão hòa số, P4 Singapore). Ở Nhóm IV, DAI
dương nhưng phụ thuộc giai đoạn và chịu thiên lệch chọn lựa (P3 Việt
Nam).

Chiều thứ hai là mức độ quốc tế hóa (FSTS): DAI là nguồn lực mở rộng
tình huống, hiệu quả hơn ở FSTS cao nơi nhu cầu điều phối xuyên biên
giới dày đặc (P4 Singapore: DAI×FSTS² = +3,119, p = 0,005).

Chiều thứ ba là mức độ bão hòa số: khi nền kinh tế đạt bão hòa Bậc 1--2,
DAI mất tác dụng phân biệt cấp doanh nghiệp (Singapore: tỉ lệ có website
67\%); khi chưa bão hòa, DAI vẫn có tín hiệu phân biệt (Việt Nam 2023:
49,8\%).

Yang, Zhao và Wei (2025) tổng hợp bằng chứng từ Trung Quốc, Ấn Độ, Việt
Nam cho thấy năng lực số là biến điều tiết mạnh nhất cho quan hệ quốc tế
hóa và hiệu quả trong giai đoạn 2018 trở đi, nhưng chỉ khi phân tách
đúng cấp độ năng lực.

\hypertarget{tux1ed5ng-hux1ee3p-4-tux1ea7ng-luxfd-thuyux1ebft}{%
\subparagraph{Tổng hợp 4 tầng lý
thuyết}\label{tux1ed5ng-hux1ee3p-4-tux1ea7ng-luxfd-thuyux1ebft}}

Bảng 2.1. Mỗi tầng lý thuyết, câu hỏi phụ, biến số, giả thuyết, bằng
chứng neo đậu.

\begin{longtable}[]{@{}
  >{\raggedright\arraybackslash}p{(\columnwidth - 10\tabcolsep) * \real{0.1667}}
  >{\raggedright\arraybackslash}p{(\columnwidth - 10\tabcolsep) * \real{0.1667}}
  >{\raggedright\arraybackslash}p{(\columnwidth - 10\tabcolsep) * \real{0.1667}}
  >{\raggedright\arraybackslash}p{(\columnwidth - 10\tabcolsep) * \real{0.1667}}
  >{\raggedright\arraybackslash}p{(\columnwidth - 10\tabcolsep) * \real{0.1667}}
  >{\raggedright\arraybackslash}p{(\columnwidth - 10\tabcolsep) * \real{0.1667}}@{}}
\toprule\noalign{}
\begin{minipage}[b]{\linewidth}\raggedright
Tầng
\end{minipage} & \begin{minipage}[b]{\linewidth}\raggedright
Lý thuyết
\end{minipage} & \begin{minipage}[b]{\linewidth}\raggedright
Câu hỏi phụ
\end{minipage} & \begin{minipage}[b]{\linewidth}\raggedright
Biến số chính
\end{minipage} & \begin{minipage}[b]{\linewidth}\raggedright
Giả thuyết
\end{minipage} & \begin{minipage}[b]{\linewidth}\raggedright
Bằng chứng P3/P4/P5
\end{minipage} \\
\midrule\noalign{}
\endhead
\bottomrule\noalign{}
\endlastfoot
1 & Uppsala & Cường độ QTH có hiệu ứng phi tuyến không? & I (FSTS,
\(\text{FSTS}^2\)) & H1 Phi tuyến chữ U ngược & P3: TP 39--46\%; P5: TP
47--49\% \\
2 & RBV & Năng lực doanh nghiệp ảnh hưởng thế nào đến hấp thụ lợi ích I?
& TCI & H2 TCI điều tiết (+) & P3: \(\beta = 0{,}179\) nhân quả (IV);
P5: +0,260 đến +0,426 \\
Digital Adoption Lens & CDCM & Năng lực số có cơ chế khác TCI không? &
DAI & H3a--H3b DAI phụ thuộc chế độ thể chế & P4:
\(\text{DAI} \times \text{FSTS}^2 = +3{,}119\) (Nhóm I); P3: IV null
(\(\beta = 0{,}018\), Nhóm IV) \\
4 & Upper Echelons & Đặc điểm nhà quản trị tác động thế nào? & Nhà quản
trị cấp cao & H4a--H4c Quản trị điều tiết & (kiểm định trong M5, dữ liệu
WBES) \\
3 & Institutional & Thể chế nội địa định hình quan hệ ra sao? & ICRV (6
nhóm) & H5 Phổ thể chế & Phân tích mô tả CĐ1: phân tán LP tăng theo chế
độ \\
-- & (Thời gian) & Quan hệ thay đổi qua thời gian không? & Year-bucket &
H6 Không đồng nhất thời gian & P3: Paternoster p \textless{} 0,001; P5:
TCI Paternoster p = 0,011 \\
\end{longtable}

Lập luận tích hợp: quốc tế hóa là quyết định chiến lược theo logic
Uppsala (tầng 1); kết quả phụ thuộc nguồn lực doanh nghiệp, TCI
(RBV/Lall, 1992) đóng vai trò nâng mặt bằng nhân quả; được hỗ trợ bởi hạ
tầng số có điều kiện, DAI (Foundational Digital Adoption
Framework/Verhoef et al., 2021) là nguồn lực mở rộng tình huống theo
CDCM; phản ánh đặc điểm nhà quản trị, gồm định hướng quốc tế của đội ngũ
quản trị cấp cao (Upper Echelons/Hambrick \& Mason, 1984; Nielsen \&
Nielsen, 2011); bị khuôn khổ bởi thể chế nội địa (Institutional
Theory/North, 1990), phổ thể chế ICRV 6 nhóm với gánh nặng bắt buộc ở
quốc đảo nhỏ; và thay đổi theo bối cảnh thời gian, thiết kế ba giai
đoạn.

Hạn chế của khung tích hợp: khó tách hiệu ứng riêng từng tầng nếu không
có dữ liệu bảng dài; một số biến quan trọng theo Uppsala (khoảng cách
tâm lý, gắn kết mạng lưới) khó đo trong WBES; TCI và DAI có thể tương
quan; ICRV và country hiệu ứng cố định có thể đa cộng; CDCM chỉ đo ở Bậc
1--2 do hạn chế WBES.

\hypertarget{khung-khuxe1i-niux1ec7m-vuxe0-muxf4-huxecnh-nghiuxean-cux1ee9u-ux111ux1ec1-xuux1ea5t}{%
\paragraph{2.3.3 Khung khái niệm và mô hình nghiên cứu đề
xuất}\label{khung-khuxe1i-niux1ec7m-vuxe0-muxf4-huxecnh-nghiuxean-cux1ee9u-ux111ux1ec1-xuux1ea5t}}

Khung khái niệm được xây dựng theo quy trình năm bước (Whetten, 1989;
Dubin, 1978): (1) xác định khoảng trống lý thuyết; (2) lựa chọn và biện
minh lý thuyết nền; (3) phát triển giả thuyết có cơ chế nhân quả rõ
ràng; (4) thiết kế sơ đồ trực quan nhất quán với hệ giả thuyết; (5) kiểm
tra tính đồng nhất giữa mô hình, giả thuyết và chiến lược thực nghiệm.
Mô hình kết hợp năm tầng lý thuyết, Uppsala Model, lý thuyết dựa trên
nguồn lực, lý thuyết thể chế, lý thuyết thượng tầng quản trị và Khung
chấp nhận số nền tảng, vào một cấu trúc nhân quả kiểm định được, nhận
diện rõ vai trò của từng biến: \textbf{biến độc lập} (quốc tế hóa, đo
bằng FSTS và dạng phi tuyến FSTS²\_c), \textbf{biến phụ thuộc} (hiệu
quả, đo bằng ln(Labour Productivity)), \textbf{biến điều tiết} cấp doanh
nghiệp (TCI, DAI) và cấp quốc gia (ICRV regime, thời gian), \textbf{biến
kiểm soát} (quy mô, tuổi, sở hữu nước ngoài, sector FE, country FE, year
FE).

\hypertarget{huxecnh-2.1-khung-khuxe1i-niux1ec7m-tuxedch-hux1ee3p-tux1ed5ng-quan}{%
\subparagraph{Hình 2.1: Khung khái niệm tích hợp (tổng
quan)}\label{huxecnh-2.1-khung-khuxe1i-niux1ec7m-tuxedch-hux1ee3p-tux1ed5ng-quan}}

\begin{figure}
\centering
\includegraphics{figures/hinh_2_1_khung_khai_niem.png}
\caption{Hình 2.1: Khung khái niệm tích hợp CĐ2}
\end{figure}

\emph{Hình 2.1.} Khung khái niệm tích hợp Chuyên đề 2 (H1--H6, năm tầng
lý thuyết).

Mũi tên liền nét biểu thị tác động trực tiếp (H1: quan hệ phi tuyến
I--P). Mũi tên đứt nét biểu thị tác động điều tiết. H1--H6 tương ứng với
các giả thuyết phát triển trong Mục 2.3.4. Biến phụ thuộc (ln(Labour
Productivity)) được xác định duy nhất ở cực phải. Biến độc lập (FSTS,
FSTS²\_c, FSTS³\_c) ở cực trái. Năm nhóm biến điều tiết phân theo hai
cấp: cấp doanh nghiệp (TCI, H2, DAI, H3, đặc điểm nhà quản trị,
H4a--H4c) và cấp quốc gia (ICRV regime, H5, giai đoạn thời gian, H6).
Biến kiểm soát (ln(emp), firm\_age, foreign\_own, ngành/country/year FE)
được đưa vào mô hình thực nghiệm nhưng không hiển thị trong sơ đồ để đơn
giản hóa trình bày. Mỗi tầng lý thuyết cung cấp cơ chế nhân quả riêng
biệt: Uppsala giải thích đường cong quốc tế hóa và hiệu quả; lý thuyết
dựa trên nguồn lực biện minh cho H2; Khung chấp nhận số nền tảng (CDCM)
biện minh cho H3; lý thuyết thượng tầng quản trị biện minh cho H4; lý
thuyết thể chế biện minh cho H5. Bằng chứng neo đậu từ P3 Việt Nam, P4
Singapore và P5 Trung Quốc xác nhận tính khả nghiệm trước khi áp dụng
cho toàn mẫu gộp 88.869 doanh nghiệp (50 nền, gồm Nhật Bản).

\hypertarget{huxecnh-2.2-cux1a1-chux1ebf-cdcm-chi-tiux1ebft}{%
\subparagraph{Hình 2.2: Cơ chế CDCM chi
tiết}\label{huxecnh-2.2-cux1a1-chux1ebf-cdcm-chi-tiux1ebft}}

\begin{figure}
\centering
\includegraphics{figures/hinh_2_2_cdcm_detail.png}
\caption{Hình 2.2: Cơ chế CDCM theo ba nhóm ICRV}
\end{figure}

\emph{Hình 2.2.} Cơ chế Country--Digital--Capability Moderation (CDCM)
theo ba nhóm ICRV đại diện.

DAI (Digital Adoption Index) không phải là nguồn lực đồng nhất mà là
\textbf{nguồn lực phụ thuộc ngữ cảnh} (context-contingent resource), vai
trò của DAI trong quan hệ quốc tế hóa và hiệu quả thay đổi theo regime
thể chế và mức bão hòa số của nền kinh tế. Nhóm I (Singapore, Advanced
I, bão hòa số cao): DAI Bậc 1+2 vận hành như nguồn lực khuếch đại tình
huống, tác động chỉ xuất hiện khi cường độ xuất khẩu đạt ngưỡng cao
(FSTS²×DAI = +3,119; P4). Nhóm IV (Việt Nam, Lower\_mid\_transition, bão
hòa số thấp; Nhóm IV trong cả hệ 5 nhóm P6 lẫn hệ 6 nhóm CD2): DAI Bậc 1
là ràng buộc đo lường; tác động dương theo OLS nhưng biến mất sau hiệu
chỉnh nội sinh biến công cụ (β = 0,018, p = 0,942; P3). Nhóm III (Trung
Quốc, Upper-Middle/Emerging transition, bão hòa số trung bình): DAI Bậc
1 phổ biến \textgreater60\%, mất khả năng phân biệt, giữ như biến kiểm
soát không phải biến điều tiết (P5). Cơ chế này giải thích tại sao kết
quả ``DAI null'' và ``DAI dương'' đều nhất quán với lý thuyết CDCM, khác
nhau ở điều kiện bối cảnh, không mâu thuẫn về lý thuyết nền tảng.

\hypertarget{bux1ea3ng-2.4-mapping-biux1ebfn-ux111ux1ea7y-ux111ux1ee7}{%
\subparagraph{Bảng 2.4: Mapping biến đầy
đủ}\label{bux1ea3ng-2.4-mapping-biux1ebfn-ux111ux1ea7y-ux111ux1ee7}}

Bảng 2.4. Mapping biến đầy đủ: khái niệm lý thuyết, biến WBES, kỳ vọng
dấu, giả thuyết.

\begin{longtable}[]{@{}
  >{\raggedright\arraybackslash}p{(\columnwidth - 10\tabcolsep) * \real{0.1034}}
  >{\raggedright\arraybackslash}p{(\columnwidth - 10\tabcolsep) * \real{0.1897}}
  >{\raggedright\arraybackslash}p{(\columnwidth - 10\tabcolsep) * \real{0.1552}}
  >{\raggedright\arraybackslash}p{(\columnwidth - 10\tabcolsep) * \real{0.1897}}
  >{\raggedright\arraybackslash}p{(\columnwidth - 10\tabcolsep) * \real{0.1552}}
  >{\raggedright\arraybackslash}p{(\columnwidth - 10\tabcolsep) * \real{0.2069}}@{}}
\toprule\noalign{}
\begin{minipage}[b]{\linewidth}\raggedright
Biến
\end{minipage} & \begin{minipage}[b]{\linewidth}\raggedright
Khái niệm
\end{minipage} & \begin{minipage}[b]{\linewidth}\raggedright
Mã WBES
\end{minipage} & \begin{minipage}[b]{\linewidth}\raggedright
Công thức
\end{minipage} & \begin{minipage}[b]{\linewidth}\raggedright
Kỳ vọng
\end{minipage} & \begin{minipage}[b]{\linewidth}\raggedright
Giả thuyết
\end{minipage} \\
\midrule\noalign{}
\endhead
\bottomrule\noalign{}
\endlastfoot
ln(LP) & Năng suất lao động & d2, l1 & ln(d2/l1) & Biến phụ thuộc &
-- \\
FSTS & Cường độ QTH & d3c & d3c/100 & \(\beta_1 > 0\) & H1 \\
\(\text{FSTS}^2\) & Phi tuyến & -- & \((\text{FSTS}_c)^2\) &
\(\beta_2 < 0\) & H1 \\
\(\text{FSTS}^3\) & Chữ S & -- & \((\text{FSTS}_c)^3\) & \(\beta_3 > 0\)
& H1 \\
TCI & Năng lực CN nội tại & b8, h8, h1, e6 & z-mean (\$\geq\$3/4 items)
& \(\beta > 0\) (trực tiếp) & H2 \\
DAI & Áp dụng số nền tảng & c22b; k33/k38 & z-mean (Bậc 1+2) & Phụ thuộc
ngữ cảnh & H3 \\
\(\text{FSTS} \times \text{TCI}\) & Điều tiết TCI & -- & tương tác &
\(\beta_\text{mod}\): phụ thuộc ngữ cảnh & H2 \\
\(\text{FSTS}^2 \times \text{DAI}\) & DAI tình huống & -- & tương tác &
\(\beta > 0\) (Nhóm I); null (Nhóm III, IV) & H3 \\
exp\_manager & Kinh nghiệm QL & b5 & năm kinh nghiệm & \(\beta > 0\)
(\(\delta_\text{exp} > 0\)) & H4a \\
educ\_manager & Học vấn QL & b7a & thứ tự & \(\beta > 0\)
(\(\delta_\text{educ} > 0\)) & H4b \\
gender\_manager & Giới tính QL & b7 & binary female=1 & Khám phá, hai
phía & H4c \\
\(\text{ICRV}_j\) & Chế độ thể chế & -- & j = I--VI, biến giả & Phổ H5 &
H5 \\
Year\_bucket & Thời gian & year & 2009--12/13--17/18--25 & Dịch chuyển
H6 & H6 \\
ln(employees) & Quy mô doanh nghiệp & l1 & ln(l1) & \(\beta > 0\) & Kiểm
soát \\
firm\_age & Tuổi doanh nghiệp & b6 & năm thành lập & \(\beta\) không xác
định & Kiểm soát \\
foreign\_own & Sở hữu nước ngoài & b2b & \$\geq\$10\% = 1 &
\(\beta > 0\) & Kiểm soát \\
\end{longtable}

Ghi chú đo lường: TCI là tổng hợp formative từ: chứng nhận ISO (b8=1),
R\&D (h8=1), đổi mới sản phẩm (h1=1), công nghệ nước ngoài có bản quyền
(e6=1). Yêu cầu ≥3/4 items không thiếu; z-chuẩn hóa trong mỗi sóng.
DAI\_z\_full (Bậc 1+2): z-mean(c22b + k33 + k38) chỉ có trong lược đồ
BEE (2019 trở đi). DAI\_z\_tier1 (Bậc 1): z-mean(c22b) có xuyên tất cả
lược đồ. FSTS được tính trung tâm theo mean trong mỗi ô quốc gia-sóng
trước khi tính FSTS² để giảm đa cộng tuyến.

\hypertarget{hux1ec7-giux1ea3-thuyux1ebft-nghiuxean-cux1ee9u-h1h6}{%
\paragraph{2.3.4 Hệ giả thuyết nghiên cứu
(H1--H6)}\label{hux1ec7-giux1ea3-thuyux1ebft-nghiuxean-cux1ee9u-h1h6}}

\hypertarget{tux1ed5ng-quan-vuxe0-logic-phuxe1t-triux1ec3n-giux1ea3-thuyux1ebft}{%
\subparagraph{Tổng quan và logic phát triển giả
thuyết}\label{tux1ed5ng-quan-vuxe0-logic-phuxe1t-triux1ec3n-giux1ea3-thuyux1ebft}}

Hệ giả thuyết H1--H6 được phát triển từ khung lý thuyết tích hợp và bằng
chứng thực nghiệm từ tài liệu, với neo đậu trong ba bản thảo đồng hành
P3--P5. Bốn nguyên tắc phát triển giả thuyết: cụ thể hóa dự đoán dấu;
neo đậu thực nghiệm từ ít nhất một bằng chứng P3/P4/P5 hoặc tổng quan
định lượng; phân biệt giả thuyết xác nhận và khám phá; liên kết với
CDCM.

H1--H3 và H5 là xác nhận (dấu rõ, ≥2 bằng chứng neo đậu từ P3/P4/P5).
H4a--H4b và H6 là hỗn hợp xác nhận/khám phá (hướng rõ nhưng độ lớn khám
phá). H4c là hoàn toàn khám phá (hai phía, không có dấu tiên nghiệm).

\hypertarget{h1-phi-tuyux1ebfn-trong-quan-hux1ec7-quux1ed1c-tux1ebf-huxf3a-vuxe0-hiux1ec7u-quux1ea3}{%
\subparagraph{H1: Phi tuyến trong quan hệ quốc tế hóa và hiệu
quả}\label{h1-phi-tuyux1ebfn-trong-quan-hux1ec7-quux1ed1c-tux1ebf-huxf3a-vuxe0-hiux1ec7u-quux1ea3}}

Cơ sở lý thuyết: Mô hình Uppsala (Johanson \& Vahlne, 1977) đề xuất quốc
tế hóa diễn ra theo ba giai đoạn với tốc độ và đặc điểm chi phí-lợi ích
khác nhau. Contractor và cộng sự (2003) và Lu và Beamish (2004) hình
thức hóa ba giai đoạn thành đường cong chữ S: giai đoạn học tập tốn kém;
giai đoạn hái quả; giai đoạn quá mức. Với các doanh nghiệp vừa và nhỏ
châu Á trong WBES, giai đoạn đầu ngắn, do đó dạng chữ U ngược (H1 bậc
hai, M1) có thể mạnh hơn chữ S đầy đủ (H1 bậc ba, M2).

Điểm uốn và chính sự tồn tại của chữ U ngược khác nhau theo chế độ thể
chế: trên khung 50 nền, chữ U ngược định hình rõ nét và tin cậy nhất ở
Nhóm IV (điểm uốn 43,0\%, trùng các ước lượng đơn quốc gia Việt Nam
39--46\% và Ấn Độ 40,7\%); ở Nhóm I đường cong duỗi thẳng gần tuyến tính
dương (điểm uốn ngoài miền dữ liệu, trùng Singapore); ở Nhóm V quan hệ
mất ý nghĩa thống kê và Nhóm VI không có chữ U ngược
(\textasciitilde50--55\%). Nhóm VI có khuôn mẫu gánh nặng bắt buộc,
không có chữ U ngược điển hình. Lưu ý: P4 Singapore (TP
\textasciitilde88,6\%) là ngoại lệ đơn quốc gia với LM p = 0,303 (không
có ý nghĩa), không đại diện cho mẫu Nhóm I trong P7.

Bằng chứng neo đậu: P3 Việt Nam (Nhóm IV): Lind--Mehlum p \textless{}
0,001 (tổng hợp), điểm uốn 39--46\%. P5 Trung Quốc (Nhóm III):
Lind--Mehlum xác nhận hai sóng, điểm uốn 47--49\%. P4 Singapore (Nhóm
I): Lind--Mehlum p = 0,303, chủ yếu dương, điểm uốn
\textasciitilde88,6\% (gần trần).

H1: Quan hệ giữa cường độ quốc tế hóa (FSTS) và hiệu quả hoạt động kinh
doanh của doanh nghiệp (ln LP) có dạng phi tuyến, với β₁(FSTS)
\textgreater{} 0 và β₂(FSTS²) \textless{} 0, phản ánh chữ U ngược ở phần
lớn các chế độ ICRV. Điểm uốn \textbf{giảm} theo phổ thể chế ICRV:
\textbf{thấp nhất ở Nhóm I (\textasciitilde28\%)}, tăng dần qua Nhóm III
(\textasciitilde47--49\%) sang Nhóm IV (\textasciitilde39--46\%) sang
Nhóm V (\textasciitilde50--55\%), với Nhóm VI có khuôn mẫu gánh nặng bắt
buộc (không có điểm uốn). Thể chế càng mạnh, doanh nghiệp đạt đỉnh hiệu
suất xuất khẩu ở mức cường độ thấp hơn.

Cơ chế biên tham gia (của H1): Bước chuyển từ không xuất khẩu (FSTS = 0)
sang xuất khẩu (FSTS \textgreater{} 0) tạo ra mức tăng năng suất lao
động đáng kể và dương, tác động tham gia thị trường xuất khẩu là động
lực chính của hiệu ứng quốc tế hóa và hiệu quả.

Cơ chế biên cường độ (của H1, gần phẳng): Trong mẫu chỉ gồm doanh nghiệp
xuất khẩu, đường cong quốc tế hóa và hiệu quả gần phẳng, tức là sau khi
đã xuất khẩu, việc tăng cường độ tiếp theo không tạo thêm nhiều lợi ích
năng suất nếu không có TCI và DAI đủ mạnh.

\hypertarget{h2-tci-ux111iux1ec1u-tiux1ebft-tuxedch-cux1ef1c-nuxe2ng-mux1eb7t-bux1eb1ng}{%
\subparagraph{H2: TCI điều tiết tích cực (nâng mặt
bằng)}\label{h2-tci-ux111iux1ec1u-tiux1ebft-tuxedch-cux1ef1c-nuxe2ng-mux1eb7t-bux1eb1ng}}

Cơ sở lý thuyết: lý thuyết dựa trên nguồn lực (Barney, 1991) cho rằng
nguồn lực VRIN tạo lợi thế cạnh tranh bền vững. TCI, năng lực công nghệ
nội tại, là nguồn lực VRIN theo Lall (1992) và tích lũy năng lực hấp thụ
theo Cohen và Levinthal (1990). Doanh nghiệp có TCI cao: hấp thụ tri
thức từ thị trường nước ngoài hiệu quả hơn; có công nghệ sản xuất tiên
tiến hơn, năng suất cao hơn ở mọi mức FSTS; và có thể duy trì hiệu quả ở
FSTS cao hơn. Cơ chế chính là nâng mặt bằng (level-shift), nâng toàn bộ
đường cong LP, không nhất thiết thay đổi vị trí điểm uốn.

Bằng chứng neo đậu: P3 Việt Nam: β(TCI) = 0,179 (p \textless{} 0,001)
bền vững 3 sóng; biến công cụ cho β(TCI) = 1,639 (p \textless{} 0,001, F
bậc một = 22,1), TCI là nhân quả. P5 Trung Quốc: β\_z(TCI) tăng từ
+0,260 (2012) sang +0,426 (2024), Paternoster z = −2,55 (p = 0,011), TCI
nâng mặt bằng tăng cường theo thời gian. P4 Singapore: β(TCI) = +0,153
(p \textless{} 0,01), không điều tiết độ cong, TCI là nâng mặt bằng
thuần túy trong Nhóm I.

H2: Chỉ số năng lực công nghệ (TCI) có mối quan hệ dương trực tiếp với
hiệu quả hoạt động kinh doanh của doanh nghiệp (β(TCI) \textgreater{} 0)
và nâng mặt bằng ln(LP) của toàn bộ đường cong quốc tế hóa và hiệu quả
(hiệu ứng nâng mặt bằng). Tác động điều tiết của TCI lên độ cong của
đường cong là câu hỏi thực nghiệm mở và được kiểm định như H2 khám phá
trong M3.

\hypertarget{h3ah3b-dai-ux111iux1ec1u-tiux1ebft-phux1ee5-thuux1ed9c-chux1ebf-ux111ux1ed9-thux1ec3-chux1ebf-cdcm}{%
\subparagraph{H3a--H3b: DAI điều tiết phụ thuộc chế độ thể chế
(CDCM)}\label{h3ah3b-dai-ux111iux1ec1u-tiux1ebft-phux1ee5-thuux1ed9c-chux1ebf-ux111ux1ed9-thux1ec3-chux1ebf-cdcm}}

Cơ sở lý thuyết: Banalieva và Dhanaraj (2019) và Verhoef và cộng sự
(2021) đề xuất DAI làm giảm khoảng cách tâm lý và chi phí giao dịch
xuyên biên giới. Tuy nhiên, CDCM cho thấy tác động này không đồng nhất
mà phụ thuộc vào ba chiều ngữ cảnh đã mô tả.

TCI hoạt động như cơ chế đào sâu nguồn lực (nội tại, tích lũy theo thời
gian, khó bắt chước), trong khi DAI hoạt động như cơ chế áp dụng phụ
thuộc (ngoại tại, phụ thuộc hệ sinh thái số; Li và cộng sự, 2022).
Coltman và cộng sự (2008) xác nhận hai cấu trúc thỏa mãn tiêu chí tổng
hợp formative riêng biệt.

Bằng chứng neo đậu: P4 Singapore: FSTS²×DAI = +3,119 (p = 0,005), DAI mở
rộng tình huống có điều kiện. P3 Việt Nam: biến công cụ null β = 0,018
(p = 0,942), OLS dương = phụ thuộc chọn lựa, không nhân quả. P5 Trung
Quốc: không điều tiết độ cong.

Bảng 2.6. CDCM, Dự đoán tác động DAI theo 3 chiều ngữ cảnh.

\begin{longtable}[]{@{}
  >{\raggedright\arraybackslash}p{(\columnwidth - 10\tabcolsep) * \real{0.2055}}
  >{\raggedright\arraybackslash}p{(\columnwidth - 10\tabcolsep) * \real{0.1370}}
  >{\raggedright\arraybackslash}p{(\columnwidth - 10\tabcolsep) * \real{0.1507}}
  >{\raggedright\arraybackslash}p{(\columnwidth - 10\tabcolsep) * \real{0.1233}}
  >{\raggedright\arraybackslash}p{(\columnwidth - 10\tabcolsep) * \real{0.2740}}
  >{\raggedright\arraybackslash}p{(\columnwidth - 10\tabcolsep) * \real{0.1096}}@{}}
\toprule\noalign{}
\begin{minipage}[b]{\linewidth}\raggedright
Chế độ (ICRV)
\end{minipage} & \begin{minipage}[b]{\linewidth}\raggedright
Mức FSTS
\end{minipage} & \begin{minipage}[b]{\linewidth}\raggedright
Bão hòa số
\end{minipage} & \begin{minipage}[b]{\linewidth}\raggedright
DAI OLS
\end{minipage} & \begin{minipage}[b]{\linewidth}\raggedright
DAI nhân quả (IV)
\end{minipage} & \begin{minipage}[b]{\linewidth}\raggedright
Cơ chế
\end{minipage} \\
\midrule\noalign{}
\endhead
\bottomrule\noalign{}
\endlastfoot
Nhóm I (Singapore) & Cao (\textgreater50\%) & Bão hòa (\textgreater65\%)
& Null trực tiếp; \(\text{FSTS}^2 \times \text{DAI} = +3{,}119\) &
(không cần IV) & Mở rộng tình huống ở nhu cầu điều phối cao \\
Nhóm III (Trung Quốc) & Trung bình (20--50\%) & Trung bình--cao
(\textgreater55\%) & Null/kiểm soát & Null (Bậc 1 mất tín hiệu phân
biệt) & Bậc 1 bão hòa \\
Nhóm IV (Việt Nam) & Thấp--trung bình (\textless40\%) & Trung bình
(\textless55\%) & Dương nhỏ (0,095--0,175) & Null (\(\beta = 0{,}018\);
F bậc một = 34,6) & Phụ thuộc chọn lựa \\
Nhóm V (Đang nổi) & Thấp (\textless20\%) & Thấp (\textless30\%) & Không
xác định & Cần IV & Hạ tầng số thiếu \\
Nhóm VI (SIDS) & Cao do bắt buộc (\textgreater40\%) & Thấp--trung & Gánh
nặng bắt buộc trội & -- & I--P gánh nặng bắt buộc che lấp tín hiệu
DAI \\
\end{longtable}

H3: Mức độ áp dụng số (DAI) có cơ chế điều tiết \textbf{phụ thuộc chế độ
thể chế} trong quan hệ quốc tế hóa và hiệu quả. DAI không hoạt động như
một phần thưởng năng suất đồng nhất mà như một nguồn lực mở rộng tình
huống theo CDCM, với hướng và cường độ khác nhau rõ ràng giữa
Emerging/Frontier và Advanced. Hai giả thuyết con phân biệt dấu kỳ vọng:

\textbf{H3a} (các nền đang nổi và chuyển đổi, Nhóm III--V): Trong các
nền kinh tế này, DAI điều tiết tích cực quan hệ FSTS sang ln(LP):
β(FSTS×DAI) \textgreater{} 0 và/hoặc β(FSTS²×DAI) \textgreater{} 0 trong
ước lượng OLS. Cơ chế: DAI giảm khoảng cách tâm lý, mở rộng tiếp cận thị
trường ảo, và tăng năng lực điều phối ở những bối cảnh mà website
presence vẫn là tín hiệu phân biệt (chưa bão hòa). Tuy nhiên, kiểm định
bằng biến công cụ tại Nhóm IV (Việt Nam) cho kết quả null (β = 0,018; p
= 0,942), OLS dương phản ánh phụ thuộc chọn lựa, không nhân quả. H3a do
đó là giả thuyết có điều kiện, cần kiểm định nhân quả để xác nhận.

\textbf{H3b} (Advanced, Nhóm I--II): Trong các nền kinh tế tiên tiến,
tác động trực tiếp không điều kiện của DAI lên ln(LP) có thể null hoặc
âm do bão hòa số và single-component measurement (chỉ website binary,
Bậc 1 proxy, không phản ánh đầy đủ năng lực số). Tuy nhiên, tác động
điều tiết điều kiện β(FSTS²×DAI) vẫn dương ở mức FSTS cao nơi nhu cầu
điều phối xuyên biên giới cao (bằng chứng neo đậu: P4 Singapore
β(FSTS²×DAI) = +3,119, p = 0,005).

TCI và DAI là hai cấu trúc phân biệt và không thể thay thế nhau trong cả
hai giả thuyết con (giả thuyết tính phân biệt CDCM, H3c: β(TCI) ≠ β(DAI)
và hai cấu trúc thỏa mãn tiêu chí discriminant validity).

\hypertarget{h4ah4c-ux111ux1eb7c-ux111iux1ec3m-nhuxe0-quux1ea3n-trux1ecb-ux111iux1ec1u-tiux1ebft}{%
\subparagraph{H4a--H4c: Đặc điểm nhà quản trị điều
tiết}\label{h4ah4c-ux111ux1eb7c-ux111iux1ec3m-nhuxe0-quux1ea3n-trux1ecb-ux111iux1ec1u-tiux1ebft}}

Cơ sở lý thuyết: lý thuyết thượng tầng quản trị (Hambrick \& Mason,
1984; Hambrick, 2007) cho rằng quyết định chiến lược phản ánh đặc điểm
nhân khẩu học của nhà quản trị cấp cao, bao gồm kinh nghiệm, học vấn, và
giới tính. Cơ chế điều tiết hoạt động qua hai kênh: chất lượng quyết
định (tránh quốc tế hóa sớm ở FSTS thấp) và năng lực điều phối (quản lý
phức tạp đa thị trường ở FSTS cao). Trong bối cảnh WBES, ba biến đại
diện có sẵn: số năm kinh nghiệm quản lý (b5), trình độ học vấn (b7a), và
giới tính nhà quản trị cấp cao (b7).

Ba giả thuyết con được tách biệt theo đặc điểm và dấu kỳ vọng:

\textbf{H4a}: Kinh nghiệm quản lý tổng thể điều tiết tích cực quan hệ
FSTS sang ln(LP): β(FSTS×exp\_manager) \textgreater{} 0. Nhà quản trị
nhiều kinh nghiệm chuyển hóa quốc tế hóa thành hiệu quả tốt hơn thông
qua khả năng nhận diện cơ hội thị trường và tránh bẫy tăng trưởng quá
mức (δ\_exp \textgreater{} 0).

\textbf{H4b}: Học vấn nhà quản trị, đại diện cho năng lực nhận thức và
khả năng tiếp thu kinh nghiệm quốc tế, điều tiết tích cực quan hệ FSTS
sang ln(LP): β(FSTS×educ\_manager) \textgreater{} 0. Nhà quản trị học
vấn cao có khung nhận thức rộng hơn để xử lý sự phức tạp thông tin trong
quản lý đa thị trường (δ\_educ \textgreater{} 0). Lưu ý: WBES (b7a)
không đo trực tiếp kinh nghiệm quốc tế; học vấn được dùng như biến đại
diện cho năng lực tư duy vượt biên giới (cognitive international
readiness).

\textbf{H4c}: Giới tính nữ của nhà quản trị cấp cao (gender\_manager =
1), kiểm định khám phá, không đặt dấu kỳ vọng tiên nghiệm; kiểm định hai
phía. Bằng chứng trong tài liệu lý thuyết thượng tầng quản trị không
đồng nhất về chiều hướng cho châu Á (Hambrick, 2007); H4c vì vậy là khám
phá.

Phạm vi kiểm định: H4a--H4b là biến đại diện (proxy) cho cơ chế lý
thuyết thượng tầng quản trị, kết quả là chỉ báo gián tiếp, không phải
kiểm định trực tiếp cấu trúc lý thuyết gốc. Hạn chế này cần được thừa
nhận rõ trong phần thảo luận kết quả.

\hypertarget{h5-thux1ec3-chux1ebf-icrv-ux111iux1ec1u-tiux1ebft-theo-phux1ed5-thux1ec3-chux1ebf}{%
\subparagraph{H5: Thể chế ICRV điều tiết theo phổ thể
chế}\label{h5-thux1ec3-chux1ebf-icrv-ux111iux1ec1u-tiux1ebft-theo-phux1ed5-thux1ec3-chux1ebf}}

Cơ sở lý thuyết: lý thuyết thể chế (North, 1990; Khanna \& Palepu, 2010)
cho rằng thể chế quy định chi phí giao dịch, rủi ro hợp đồng, và khả
năng tiếp cận thị trường, tất cả ảnh hưởng đến hiệu quả của quốc tế hóa.
Ba cơ chế độc lập giải thích tại sao thể chế tốt hơn tạo ra điểm uốn cao
hơn: (A) thực thi hợp đồng (phân biệt de jure trên giấy và de facto thực
thi thực tế; Xu, 2024), giảm chi phí đại diện trong chuỗi cung ứng quốc
tế; (B) phát triển thị trường tài chính, cho phép huy động vốn cho vốn
lưu động quốc tế; (C) chất lượng quy định, giảm gánh nặng tuân thủ cho
doanh nghiệp xuất khẩu. Kafouros và cộng sự (2023) ước tính: 1\% cải
thiện trong WGI Regulatory Quality tương đương với +2--3\% FSTS trong
điểm uốn.

Dự đoán điểm uốn theo nhóm ICRV:

\begin{longtable}[]{@{}
  >{\raggedright\arraybackslash}p{(\columnwidth - 4\tabcolsep) * \real{0.3056}}
  >{\raggedright\arraybackslash}p{(\columnwidth - 4\tabcolsep) * \real{0.5000}}
  >{\raggedright\arraybackslash}p{(\columnwidth - 4\tabcolsep) * \real{0.1944}}@{}}
\toprule\noalign{}
\begin{minipage}[b]{\linewidth}\raggedright
ICRV Nhóm
\end{minipage} & \begin{minipage}[b]{\linewidth}\raggedright
Điểm uốn dự kiến
\end{minipage} & \begin{minipage}[b]{\linewidth}\raggedright
Cơ sở
\end{minipage} \\
\midrule\noalign{}
\endhead
\bottomrule\noalign{}
\endlastfoot
Nhóm I (tiên tiến đổi mới) & \textbf{gần tuyến tính dương; điểm uốn
ngoài miền dữ liệu} & P7 50 nền: TP gộp M5 = 43,6\% (N=79.080; M2
N=81.022); trong Nhóm I đường cong không định hình (LM p = ,30), nhất
quán P4 Singapore (TP \textasciitilde88,6\%, LM p = ,303, không có ý
nghĩa) \\
Nhóm II (tiên tiến tài nguyên) & Chữ U ngược ở biên ý nghĩa
(\textasciitilde62\%) & Kết quả 50 nền: TP 62,0\%, LM p = ,081 \\
Nhóm III (trung bình cao) & Không xác định ở cấp nhóm;
\textasciitilde47--49\% ở Trung Quốc & Kết quả 50 nền: LM p = ,26; P5
Trung Quốc TP 47--49\%*** \\
Nhóm IV (chuyển đổi TB-thấp) & \textasciitilde35--46\% FSTS & P3 Việt
Nam TP 39--46\% \\
Nhóm V (đang nổi) & Quan hệ mất cấu trúc, không xác định & Kết quả 50
nền: LM p = ,18; khoảng trống thể chế làm phần thưởng nhỏ và bất định \\
Nhóm VI (SIDS) & Không có TP / gánh nặng bắt buộc & Quốc tế hóa bắt
buộc \\
\end{longtable}

H5: Hiệu ứng điều tiết của chế độ thể chế (ICRV) lên quan hệ quốc tế hóa
và hiệu quả có dạng ba vùng: chữ U ngược định hình rõ ở vùng giữa phổ
thể chế (Nhóm IV, điểm uốn 43,0\%); duỗi thẳng gần tuyến tính dương ở
biên trên (Nhóm I--II); mất cấu trúc ở biên dưới (Nhóm V không ý nghĩa;
Nhóm VI có khuôn mẫu gánh nặng bắt buộc). Thể chế càng mạnh; ở biên trên
phổ thể chế, chi phí giao dịch thấp và năng lực điều phối cao đẩy nhánh
đi xuống ra ngoài miền quan sát thay vì tạo đỉnh sớm hơn.

\hypertarget{h6-tuxednh-khuxf4ng-ux111ux1ed3ng-nhux1ea5t-theo-thux1eddi-gian}{%
\subparagraph{H6: Tính không đồng nhất theo thời
gian}\label{h6-tuxednh-khuxf4ng-ux111ux1ed3ng-nhux1ea5t-theo-thux1eddi-gian}}

Cơ sở lý thuyết: Wu, Fan và Chen (2022) tổng hợp 20 năm bằng chứng và
tìm thấy các dịch chuyển thời gian đáng kể trong quốc tế hóa và hiệu
quả. Ba giai đoạn trong dữ liệu có đặc điểm bối cảnh khác nhau rõ ràng:
2009--2012 (hậu khủng hoảng tài chính, hạ tầng số sơ khai ở châu Á);
2013--2017 (phục hồi chuỗi giá trị toàn cầu, bùng nổ Internet di động,
thương mại điện tử lan rộng); 2018--2025 (chuyển đổi số + AI bùng nổ,
COVID-19, tái cấu hình chuỗi cung ứng).

Bằng chứng neo đậu: P3 Việt Nam: Paternoster z = 3,353 giữa 2009 và 2015
(p \textless{} 0,001), quan hệ quốc tế hóa và hiệu quả dịch chuyển đáng
kể. P5 Trung Quốc: TCI Paternoster p = 0,011, hiệu ứng TCI tăng cường
2012 xuống 2024. P5 Trung Quốc: FSTS Paternoster p = 0,545, điểm uốn ổn
định (bền vững cấu trúc H2b xác nhận).

H6: Tác động của quốc tế hóa lên hiệu quả hoạt động kinh doanh của doanh
nghiệp thể hiện tính không đồng nhất theo thời gian:
β(FSTS×Year\_bucket) và β(FSTS²×Year\_bucket) khác nhau có ý nghĩa thống
kê giữa ba giai đoạn 2009--2012, 2013--2017, và 2018--2025. Cụ thể: tác
động TCI tăng cường theo thời gian; tác động DAI biến thiên theo giai
đoạn; điểm uốn có thể dịch chuyển trong một số nhóm ICRV nhưng ổn định
trong các nhóm khác.

\hypertarget{tux1ed5ng-hux1ee3p-hux1ec7-giux1ea3-thuyux1ebft}{%
\subparagraph{Tổng hợp hệ giả
thuyết}\label{tux1ed5ng-hux1ee3p-hux1ec7-giux1ea3-thuyux1ebft}}

Bảng 2.5. Hệ giả thuyết H1--H6, tóm tắt cơ chế, biến, kỳ vọng, và bằng
chứng neo đậu.

\begin{longtable}[]{@{}
  >{\raggedright\arraybackslash}p{(\columnwidth - 10\tabcolsep) * \real{0.0811}}
  >{\raggedright\arraybackslash}p{(\columnwidth - 10\tabcolsep) * \real{0.0676}}
  >{\raggedright\arraybackslash}p{(\columnwidth - 10\tabcolsep) * \real{0.2297}}
  >{\raggedright\arraybackslash}p{(\columnwidth - 10\tabcolsep) * \real{0.2162}}
  >{\raggedright\arraybackslash}p{(\columnwidth - 10\tabcolsep) * \real{0.1216}}
  >{\raggedright\arraybackslash}p{(\columnwidth - 10\tabcolsep) * \real{0.2838}}@{}}
\toprule\noalign{}
\begin{minipage}[b]{\linewidth}\raggedright
Giả thuyết
\end{minipage} & \begin{minipage}[b]{\linewidth}\raggedright
Tên
\end{minipage} & \begin{minipage}[b]{\linewidth}\raggedright
Cơ chế lý thuyết
\end{minipage} & \begin{minipage}[b]{\linewidth}\raggedright
Biến kiểm định
\end{minipage} & \begin{minipage}[b]{\linewidth}\raggedright
Kỳ vọng
\end{minipage} & \begin{minipage}[b]{\linewidth}\raggedright
Bằng chứng P3/P4/P5
\end{minipage} \\
\midrule\noalign{}
\endhead
\bottomrule\noalign{}
\endlastfoot
H1 & Phi tuyến I--P & Uppsala 3 giai đoạn; chi phí phối hợp & FSTS,
\(\text{FSTS}^2\), \(\text{FSTS}^3\) & \(\beta_1 > 0\); \(\beta_2 < 0\)
& P3 TP 39--46\%; P5 TP 47--49\% \\
H2 & TCI nâng mặt bằng & RBV năng lực hấp thụ & TCI,
\(\text{FSTS} \times \text{TCI}\) & \(\beta(\text{TCI}) > 0\); độ cong:
khám phá & P3 IV nhân quả; P5 +0,260 đến +0,426 \\
H3a & DAI nhóm đang nổi/chuyển đổi (+) & CDCM × Digital Adoption Lens &
DAI, \(\text{FSTS} \times \text{DAI}\) &
\(\beta(\text{FSTS} \times \text{DAI}) > 0\) (Nhóm III--V) & P3 IV null;
bằng chứng hạn chế \\
H3b & DAI Advanced (null/−, conditional +) & CDCM × Digital Adoption
Lens & DAI, \(\text{FSTS}^2 \times \text{DAI}\) & Direct: null/--; cond:
\(\beta(\text{FSTS}^2 \times \text{DAI}) > 0\) ở FSTS cao & P4 +3,119;
bão hòa số \\
H4a & Kinh nghiệm QL điều tiết & Upper Echelons & exp\_manager,
\(\text{FSTS} \times \text{exp}\) &
\(\beta(\text{FSTS} \times \text{exp}) > 0\) (\(\delta_\text{exp} > 0\))
& (WBES b5) \\
H4b & Học vấn QL điều tiết & Upper Echelons & educ\_manager,
\(\text{FSTS} \times \text{educ}\) &
\(\beta(\text{FSTS} \times \text{educ}) > 0\)
(\(\delta_\text{educ} > 0\)) & (WBES b7a) \\
H4c & Giới tính QL, khám phá & Upper Echelons & gender\_manager,
\(\text{FSTS} \times \text{gender}\) & Khám phá, hai phía & (WBES b7) \\
H5 & Phổ ICRV & Institutional Theory &
\(\text{ICRV}_j \times \text{FSTS}\) & TP phổ thể chế; Nhóm VI gánh nặng
bắt buộc & Phân tán LP tăng theo chế độ (CĐ1) \\
H6 & Không đồng nhất thời gian & Thay đổi cấu trúc &
\(\text{Year\_bucket} \times \text{FSTS}\) &
\(\beta(\text{FSTS} \times \text{Year}) \neq 0\) & P3 Paternoster p
\textless{} 0,001; P5 TCI p = 0,011 \\
\end{longtable}

Mười một mệnh đề con dưới sáu giả thuyết chính H1--H6 (H1 với hai cơ chế
biên tham gia và biên cường độ, H2, H3a--H3b với H3c tính phân biệt,
H4a--H4c, H5, H6) tạo thành hệ thống logic nhất quán: H1 cùng hai cơ chế
biên tham gia và biên cường độ (phi tuyến) là trọng tâm; H2 (TCI nâng
mặt bằng) và H3a/H3b (DAI phụ thuộc regime) phân rã cơ chế năng lực
doanh nghiệp với dấu kỳ vọng khác nhau theo ngữ cảnh; H4a--H4c (Quản
trị) bổ sung tầng lý thuyết thượng tầng quản trị với dấu kỳ vọng phân
biệt theo đặc điểm nhà quản trị; H5 (ICRV phổ thể chế) là tầng thể chế
bao trùm; H6 (Thời gian) là tầng động.

Tích hợp cả sáu giả thuyết trong M7 (điều tiết ba chiều) cho phép kiểm
định xem TCI và DAI là bổ sung (β(TCI×DAI×FSTS) \textgreater{} 0) hay
thay thế (β \textless{} 0) trong điều tiết quốc tế hóa và hiệu quả ở
Nhóm I, câu hỏi chưa được trả lời trong bất kỳ nghiên cứu quốc tế hóa và
hiệu quả nào trước đây.

\hypertarget{thiux1ebft-kux1ebf-muxf4-huxecnh-thux1ef1c-nghiux1ec7m-m0m7}{%
\paragraph{2.3.5 Thiết kế mô hình thực nghiệm
(M0--M7)}\label{thiux1ebft-kux1ebf-muxf4-huxecnh-thux1ef1c-nghiux1ec7m-m0m7}}

\hypertarget{nguyuxean-tux1eafc-muxf4-huxecnh}{%
\subparagraph{Nguyên tắc mô
hình}\label{nguyuxean-tux1eafc-muxf4-huxecnh}}

Tám mô hình M0--M7 được xây dựng theo cấu trúc lồng nhau: mỗi mô hình
sau bổ sung thêm biến so với mô hình trước, cho phép so sánh AIC/BIC và
R² tăng thêm, và kiểm định từng giả thuyết bằng F-test theo khối. Tất cả
mô hình sử dụng: biến phụ thuộc là ln(LP) = ln(doanh thu hằng năm trong
nội tệ / lao động thường xuyên), chuẩn hóa z trong cụm nền kinh tế × năm
(xem Mục 2.3.6); sai số chuẩn HC1 robust (Long \& Ervin, 2000) làm
chuẩn; hiệu ứng cố định quốc gia (\(\alpha_j\)) và năm (\(\delta_t\));
FSTS được tính trung tâm theo mean trước khi tính \(\mathrm{FSTS}^2\) để
giảm đa cộng tuyến.

Tiêu chí lựa chọn mô hình: \(\Delta R^2 > 0{,}01\) để thêm biến điều
tiết; \(\Delta\mathrm{AIC} < -2\) ủng hộ mô hình phức tạp hơn;
\(\Delta\mathrm{BIC} < -10\) là bằng chứng mạnh; F-test \(p < 0{,}05\)
để thêm khối biến. M7 được báo cáo nếu ΔAIC(M7 vs M4+M3 kết hợp)
\textless{} −2 VÀ F-test p \textless{} 0,05.

\hypertarget{m0-cux1a1-sux1edf-tuyux1ebfn-tuxednh}{%
\subparagraph{M0: Cơ sở tuyến
tính}\label{m0-cux1a1-sux1edf-tuyux1ebfn-tuxednh}}

\begin{equation}
\ln(\mathrm{LP})_i = \beta_0 + \beta_1 \cdot \mathrm{FSTS}_i + \boldsymbol{\gamma}'\mathbf{X}_i + \alpha_j + \delta_t + \varepsilon_i
\label{eq:M0}
\end{equation}

Mục đích: thiết lập cơ sở và kiểm tra hướng tác động tuyến tính thuần
túy. Kỳ vọng: \(\beta_1 > 0\) (tác động dương trung bình của xuất khẩu).
Hạn chế: không nắm bắt được phi tuyến tính, cung cấp ước lượng cận dưới
cho tác động quốc tế hóa.

\hypertarget{m1-chux1eef-u-ngux1b0ux1ee3c-bux1eadc-hai-kiux1ec3m-ux111ux1ecbnh-h1-bux1eadc-hai}{%
\subparagraph{M1: Chữ U ngược bậc hai (kiểm định H1 bậc
hai)}\label{m1-chux1eef-u-ngux1b0ux1ee3c-bux1eadc-hai-kiux1ec3m-ux111ux1ecbnh-h1-bux1eadc-hai}}

\begin{equation}
\ln(\mathrm{LP})_i = \beta_0 + \beta_1 \cdot \mathrm{FSTS}_{c,i} + \beta_2 \cdot \mathrm{FSTS}_{c,i}^{2} + \boldsymbol{\gamma}'\mathbf{X}_i + \alpha_j + \delta_t + \varepsilon_i
\label{eq:M1}
\end{equation}

Mục đích: kiểm định chữ U ngược trong H1, dạng hàm phổ biến nhất trong
tài liệu. Kỳ vọng: \(\beta_1 > 0\), \(\beta_2 < 0\). Điểm uốn:
\(\mathrm{TP}_1 = -\beta_1/(2\beta_2)\), mức FSTS tối ưu (trên thang
mean-centered). Kiểm định bổ sung: Lind--Mehlum (2010) để xác nhận chữ U
ngược thực sự.

\hypertarget{m2-ux111ux1b0ux1eddng-cong-chux1eef-s-bux1eadc-ba-kiux1ec3m-ux111ux1ecbnh-h1-ux111ux1ea7y-ux111ux1ee7}{%
\subparagraph{M2: đường cong chữ S bậc ba (kiểm định H1 đầy
đủ)}\label{m2-ux111ux1b0ux1eddng-cong-chux1eef-s-bux1eadc-ba-kiux1ec3m-ux111ux1ecbnh-h1-ux111ux1ea7y-ux111ux1ee7}}

\begin{equation}
\ln(\mathrm{LP})_i = \beta_0 + \beta_1 \cdot \mathrm{FSTS}_{c,i} + \beta_2 \cdot \mathrm{FSTS}_{c,i}^{2} + \beta_3 \cdot \mathrm{FSTS}_{c,i}^{3} + \boldsymbol{\gamma}'\mathbf{X}_i + \alpha_j + \delta_t + \varepsilon_i
\label{eq:M2}
\end{equation}

Mục đích: kiểm định đường cong chữ S ba giai đoạn (Lu \& Beamish, 2004;
Contractor et al., 2003). Kỳ vọng: \(\beta_1 > 0\), \(\beta_2 < 0\),
\(\beta_3 > 0\). Lựa chọn mô hình: so sánh M1 và M2 bằng AIC, BIC, và
F-test \(\beta_3 = 0\).

\hypertarget{m3-ux111iux1ec1u-tiux1ebft-tci-kiux1ec3m-ux111ux1ecbnh-h2}{%
\subparagraph{M3: + Điều tiết TCI (kiểm định
H2)}\label{m3-ux111iux1ec1u-tiux1ebft-tci-kiux1ec3m-ux111ux1ecbnh-h2}}

\begin{align}
\ln(\mathrm{LP})_i &= \underbrace{\beta_0 + \beta_1\,\mathrm{FSTS}_{c,i} + \beta_2\,\mathrm{FSTS}_{c,i}^{2}}_{M1} \notag \\
  &\quad + \beta_3 \cdot \mathrm{TCI}_i + \beta_4 \cdot (\mathrm{FSTS}_{c,i} \times \mathrm{TCI}_i) + \beta_5 \cdot (\mathrm{FSTS}_{c,i}^{2} \times \mathrm{TCI}_i) \notag \\
  &\quad + \boldsymbol{\gamma}'\mathbf{X}_i + \alpha_j + \delta_t + \varepsilon_i
\label{eq:M3}
\end{align}

Mục đích: kiểm định H2, TCI là nâng mặt bằng (\(\beta_3 > 0\)) và/hoặc
điều tiết độ cong (\(\beta_4\), \(\beta_5\)). Kỳ vọng chính:
\(\beta_3 > 0\) (hiệu ứng TCI trực tiếp). Kiểm định: F-test
(\(\beta_4 = \beta_5 = 0\)) = kiểm định điều tiết kết hợp (joint
moderation); nếu \(p > 0{,}10\) sang TCI là nâng mặt bằng thuần túy.

\hypertarget{m4-ux111iux1ec1u-tiux1ebft-dai-kiux1ec3m-ux111ux1ecbnh-h3}{%
\subparagraph{M4: + Điều tiết DAI (kiểm định
H3)}\label{m4-ux111iux1ec1u-tiux1ebft-dai-kiux1ec3m-ux111ux1ecbnh-h3}}

\begin{align}
\ln(\mathrm{LP})_i &= \underbrace{\text{M1}}_{\text{cơ sở}} \notag \\
  &\quad + \beta_3 \cdot \mathrm{DAI}_i + \beta_4 \cdot (\mathrm{FSTS}_{c,i} \times \mathrm{DAI}_i) + \beta_5 \cdot (\mathrm{FSTS}_{c,i}^{2} \times \mathrm{DAI}_i) \notag \\
  &\quad + \boldsymbol{\gamma}'\mathbf{X}_i + \alpha_j + \delta_t + \varepsilon_i
\label{eq:M4}
\end{align}

Mục đích: kiểm định H3, DAI là nguồn lực mở rộng tình huống. Kỳ vọng:
\(\beta_5(\mathrm{FSTS}^2 \times \mathrm{DAI}) > 0\) trong Nhóm I (nhất
quán P4 Singapore \(+3{,}119\)); \(\beta_3\) (DAI trực tiếp) rồi \(0\)
khi kiểm soát chọn lựa (nhất quán P3 Việt Nam IV null). Lưu ý IV: trong
mẫu con Nhóm IV, nếu có biến công cụ hợp lệ, chạy 2SLS để phân biệt DAI
nhân quả vs.~phụ thuộc chọn lựa.

\hypertarget{m5-ux111iux1ec1u-tiux1ebft-nhuxe0-quux1ea3n-trux1ecb-kiux1ec3m-ux111ux1ecbnh-h4ah4bh4c}{%
\subparagraph{M5: + Điều tiết nhà quản trị (kiểm định
H4a/H4b/H4c)}\label{m5-ux111iux1ec1u-tiux1ebft-nhuxe0-quux1ea3n-trux1ecb-kiux1ec3m-ux111ux1ecbnh-h4ah4bh4c}}

\begin{align}
\ln(\mathrm{LP})_i &= \underbrace{\text{M1}}_{\text{cơ sở}} \notag \\
  &\quad + \beta_3 \cdot \mathrm{Manager}_i + \beta_4 \cdot (\mathrm{FSTS}_{c,i} \times \mathrm{Manager}_i) \notag \\
  &\quad + \boldsymbol{\gamma}'\mathbf{X}_i + \alpha_j + \delta_t + \varepsilon_i
\label{eq:M5}
\end{align}

Biến Manager theo thứ tự ưu tiên: exp\_manager (b5, H4a), biến chính với
kỳ vọng β₄ \textgreater{} 0; educ\_manager (b7a, H4b), biến phụ với kỳ
vọng β₄ \textgreater{} 0; gender\_manager (b7, H4c), khám phá, kiểm định
hai phía. Kỳ vọng:
\(\beta_4(\mathrm{FSTS} \times \text{exp\_manager}) > 0\) {[}H4a{]};
\(\beta_4(\mathrm{FSTS} \times \text{educ\_manager}) > 0\) {[}H4b{]};
\(\beta_4(\mathrm{FSTS} \times \text{gender\_manager})\), không có kỳ
vọng tiên nghiệm {[}H4c{]}.

\hypertarget{m6-chux1ebf-ux111ux1ed9-thux1ec3-chux1ebf-icrv-kiux1ec3m-ux111ux1ecbnh-h5}{%
\subparagraph{M6: + Chế độ thể chế ICRV (kiểm định
H5)}\label{m6-chux1ebf-ux111ux1ed9-thux1ec3-chux1ebf-icrv-kiux1ec3m-ux111ux1ecbnh-h5}}

\begin{align}
\ln(\mathrm{LP})_i &= \underbrace{\text{M1}}_{\text{cơ sở}} \notag \\
  &\quad + \sum_{j=2}^{6} \Bigl[\beta_{3j} \cdot \mathrm{ICRV}_{j,i}
      + \beta_{4j} \cdot (\mathrm{FSTS}_{c,i} \times \mathrm{ICRV}_{j,i})
      + \beta_{5j} \cdot (\mathrm{FSTS}_{c,i}^{2} \times \mathrm{ICRV}_{j,i})\Bigr] \notag \\
  &\quad + \boldsymbol{\gamma}'\mathbf{X}_i + \delta_t + \varepsilon_i
\label{eq:M6}
\end{align}

Mục đích: kiểm định H5, phổ thể chế ICRV và gánh nặng bắt buộc (Nhóm
VI). Nhóm tham chiếu: Nhóm I (tiên tiến đổi mới). Kỳ vọng: cấu trúc ba
vùng: hệ số bậc hai âm và có ý nghĩa chỉ ở Nhóm IV (và biên ý nghĩa ở
Nhóm II); không ý nghĩa ở Nhóm I, III, V; Nhóm VI không có chữ U ngược;
Nhóm VI có β(FSTS×ICRV\_VI) ≤ 0 (neo đậu P8: FIP β = −1,339, SE 0,386, p
\textless{} 0,001). Lưu ý: Country FE bị omit trong M6 do ICRV\_j là
biến cấp quốc gia bất biến theo thời gian, thay bằng Year FE + Region FE
(6 vùng).

\hypertarget{m7-ux111iux1ec1u-tiux1ebft-ba-chiux1ec1u-tux1ed5ng-lux1ef1c-kiux1ec3m-ux111ux1ecbnh-h2-h3-h6-ux111ux1ed3ng-thux1eddi}{%
\subparagraph{M7: Điều tiết ba chiều tổng lực (kiểm định H2 + H3 + H6
đồng
thời)}\label{m7-ux111iux1ec1u-tiux1ebft-ba-chiux1ec1u-tux1ed5ng-lux1ef1c-kiux1ec3m-ux111ux1ecbnh-h2-h3-h6-ux111ux1ed3ng-thux1eddi}}

\begin{align}
\ln(\mathrm{LP})_i
  &= \beta_0 + \beta_1\,\mathrm{FSTS}_{c,i} + \beta_2\,\mathrm{FSTS}_{c,i}^{2}
     + \beta_3\,\mathrm{TCI}_i + \beta_4\,\mathrm{DAI}_i \notag \\
  &\quad + \beta_5\,(\mathrm{FSTS}_{c,i} \times \mathrm{TCI}_i)
     + \beta_6\,(\mathrm{FSTS}_{c,i}^{2} \times \mathrm{TCI}_i) \notag \\
  &\quad + \beta_7\,(\mathrm{FSTS}_{c,i} \times \mathrm{DAI}_i)
     + \beta_8\,(\mathrm{FSTS}_{c,i}^{2} \times \mathrm{DAI}_i) \notag \\
  &\quad + \beta_9\,(\mathrm{FSTS}_{c,i} \times \mathrm{TCI}_i \times \mathrm{DAI}_i) \notag \\
  &\quad + \beta_{10}\,(\mathrm{FSTS}_{c,i} \times YB_{2013\text{-}17,i})
     + \beta_{11}\,(\mathrm{FSTS}_{c,i}^{2} \times YB_{2013\text{-}17,i}) \notag \\
  &\quad + \beta_{12}\,(\mathrm{FSTS}_{c,i} \times YB_{2018\text{-}25,i})
     + \beta_{13}\,(\mathrm{FSTS}_{c,i}^{2} \times YB_{2018\text{-}25,i}) \notag \\
  &\quad + \boldsymbol{\gamma}'\mathbf{X}_i + \alpha_j + \delta_t + \varepsilon_i
\label{eq:M7}
\end{align}

trong đó \(\mathbf{X}_i\) gồm: ln\_employees, firm\_age, foreign\_own
(\(\ge 10\%\)), sector dummies; \(\alpha_j\) = country FE; \(\delta_t\)
= year FE.

Mục đích: mô hình tổng lực kiểm định H2 (TCI), H3 (DAI), và H6 (thời
gian) đồng thời. Kỳ vọng trọng tâm:
\(\beta_8(\mathrm{FSTS}^2 \times \mathrm{DAI}) > 0\) trong Nhóm I;
\(\beta_3(\mathrm{TCI}) > 0\);
\(\beta_{12}/\beta_{13} \ne \beta_{10}/\beta_{11}\) (dịch chuyển thời
gian). Yêu cầu mẫu: với k = 13 tham số trọng tâm + controls, n ≥ 5.000
cần thiết cho power \textgreater{} 0,80 tại f² = 0,02, dễ dàng đáp ứng
với mẫu gộp 88.869 (50 nền, gồm Nhật Bản).

\hypertarget{bux1ea3ng-tux1ed5ng-hux1ee3p-muxf4-huxecnh}{%
\subparagraph{Bảng tổng hợp mô
hình}\label{bux1ea3ng-tux1ed5ng-hux1ee3p-muxf4-huxecnh}}

Bảng 2.7. Tám mô hình M0--M7, cấu trúc, giả thuyết kiểm định, và biến bổ
sung.

\begin{longtable}[]{@{}
  >{\raggedright\arraybackslash}p{(\columnwidth - 6\tabcolsep) * \real{0.1250}}
  >{\raggedright\arraybackslash}p{(\columnwidth - 6\tabcolsep) * \real{0.4722}}
  >{\raggedright\arraybackslash}p{(\columnwidth - 6\tabcolsep) * \real{0.1667}}
  >{\raggedright\arraybackslash}p{(\columnwidth - 6\tabcolsep) * \real{0.2361}}@{}}
\toprule\noalign{}
\begin{minipage}[b]{\linewidth}\raggedright
Mô hình
\end{minipage} & \begin{minipage}[b]{\linewidth}\raggedright
Biến bổ sung so với mô hình trước
\end{minipage} & \begin{minipage}[b]{\linewidth}\raggedright
Giả thuyết
\end{minipage} & \begin{minipage}[b]{\linewidth}\raggedright
Kiểm định chính
\end{minipage} \\
\midrule\noalign{}
\endhead
\bottomrule\noalign{}
\endlastfoot
M0 & FSTS (tuyến tính) & Baseline & \(\beta_1 > 0\) \\
M1 & + \(\text{FSTS}^2\) & H1 (chữ U ngược) & \(\beta_2 < 0\);
Lind--Mehlum \\
M2 & + \(\text{FSTS}^3\) & H1 (chữ S) & \(\beta_3 > 0\); F-test
\(\beta_3 = 0\) \\
M3 & + TCI, \(\text{FSTS} \times \text{TCI}\),
\(\text{FSTS}^2 \times \text{TCI}\) & H2 & \(\beta_3 > 0\); F-test
(\(\beta_4 = \beta_5 = 0\)) \\
M4 & + DAI, \(\text{FSTS} \times \text{DAI}\),
\(\text{FSTS}^2 \times \text{DAI}\) & H3a/H3b & H3a:
\(\beta(\text{FSTS} \times \text{DAI}) > 0\) Nhóm III--V; H3b:
\(\beta(\text{FSTS}^2 \times \text{DAI}) > 0\) ở FSTS cao Nhóm I \\
M5 & + Manager, \(\text{FSTS} \times \text{Manager}\) & H4a/H4b/H4c &
H4a: \(\beta_4(\text{FSTS} \times \text{exp}) > 0\); H4b:
\(\beta_4(\text{FSTS} \times \text{educ}) > 0\); H4c: hai phía \\
M6 & + \(\text{ICRV}_j \times \text{FSTS}\),
\(\text{ICRV}_j \times \text{FSTS}^2\) & H5 & Phổ TP; gánh nặng bắt buộc
Nhóm VI \\
M7 & + \(\text{TCI} \times \text{DAI} \times \text{FSTS}\),
\(\text{Year\_bucket} \times \text{FSTS}\) & H2+H3+H6 & F-test block ba
chiều \\
\end{longtable}

\hypertarget{dux1eef-liux1ec7u-vuxe0-chiux1ebfn-lux1b0ux1ee3c-nhux1eadn-dux1ea1ng-nhuxe2n-quux1ea3}{%
\paragraph{2.3.6 Dữ liệu và chiến lược nhận dạng nhân
quả}\label{dux1eef-liux1ec7u-vuxe0-chiux1ebfn-lux1b0ux1ee3c-nhux1eadn-dux1ea1ng-nhuxe2n-quux1ea3}}

\hypertarget{nguux1ed3n-dux1eef-liux1ec7u-vuxe0-cux1ea5u-truxfac-mux1eabu-gux1ed9p}{%
\subparagraph{Nguồn dữ liệu và cấu trúc mẫu
gộp}\label{nguux1ed3n-dux1eef-liux1ec7u-vuxe0-cux1ea5u-truxfac-mux1eabu-gux1ed9p}}

Dữ liệu gốc: World Bank Enterprise Surveys (WBES), cơ sở dữ liệu điều
tra vi mô doanh nghiệp lớn nhất thế giới, phủ hơn 170 quốc gia với hơn
400.000 cuộc phỏng vấn từ 2006 đến nay. WBES áp dụng phương pháp lấy mẫu
ngẫu nhiên phân tầng theo quy mô doanh nghiệp và ngành.

Mẫu gộp chuyên đề: khung 50 nền kinh tế châu Á và Thái Bình Dương (gồm
Nhật Bản 2025), 88.869 doanh nghiệp, 103 cặp nền kinh tế--năm trong giai
đoạn 2009--2026; mẫu hồi quy chính của vòng ước lượng chính thức: M2 N =
81.022, M5 N = 79.080.

Bảng 2.8. Cấu trúc mẫu gộp theo nhóm ICRV và mốc khảo sát.

\begin{longtable}[]{@{}
  >{\raggedright\arraybackslash}p{(\columnwidth - 6\tabcolsep) * \real{0.2143}}
  >{\raggedright\arraybackslash}p{(\columnwidth - 6\tabcolsep) * \real{0.2143}}
  >{\centering\arraybackslash}p{(\columnwidth - 6\tabcolsep) * \real{0.2857}}
  >{\centering\arraybackslash}p{(\columnwidth - 6\tabcolsep) * \real{0.2857}}@{}}
\toprule\noalign{}
\begin{minipage}[b]{\linewidth}\raggedright
ICRV Nhóm (icrv\_label)
\end{minipage} & \begin{minipage}[b]{\linewidth}\raggedright
Nền kinh tế
\end{minipage} & \begin{minipage}[b]{\linewidth}\centering
Số nền
\end{minipage} & \begin{minipage}[b]{\linewidth}\centering
N doanh nghiệp
\end{minipage} \\
\midrule\noalign{}
\endhead
\bottomrule\noalign{}
\endlastfoot
I, Adv. đổi mới (\texttt{Advanced\_innovation}) & Singapore, HongKong,
Hàn Quốc, Đài Loan, Israel, Nhật Bản¹ & 6 & 6.390 \\
II, Adv. tài nguyên (\texttt{Advanced\_resource}) & Saudi Arabia, Qatar,
Kuwait, Bahrain, Brunei, Oman & 6 & 2.231 \\
III, Trung bình cao (\texttt{Upper\_mid}) & Trung Quốc, Malaysia, Thái
Lan, Kazakhstan, Armenia, Georgia & 6 & 13.993 \\
IV, Chuyển đổi TB-thấp (\texttt{Lower\_mid\_transition}) & Việt Nam, Ấn
Độ, Indonesia, Philippines, Mông Cổ, Bangladesh, Pakistan & 7 &
45.003 \\
V, Đang nổi (\texttt{Emerging}) & 17 nền (gồm Sri Lanka, Jordan, Nepal,
Cambodia, Laos, Uzbekistan\ldots) & 17 & 18.957 \\
VI, SIDS (\texttt{SIDS\_small}) & Fiji, Kiribati, PNG, Samoa, Solomon,
Tonga, Vanuatu, Timor-Leste & 8 & 2.295 \\
\textbf{Tổng} & & \textbf{50} & \textbf{88.869} (mẫu gộp mô tả; N hồi
quy M2 = 81.022) \\
\end{longtable}

\emph{Ghi chú: ¹Nhật Bản được WBES khảo sát lần đầu năm 2025 (n=2.168)
và là thành viên thứ sáu của Nhóm I trong khung 50 nền; vòng tái ước
lượng 50 nền gồm Nhật Bản đã hoàn tất nên mô tả và ước lượng (P7) nay
dùng chung MỘT khung 50 nền/88.869 quan sát (hồi quy M2 N=81.022; M5
N=79.080). Nhãn nhóm theo biến phân loại \texttt{icrv\_label}, dùng
thống nhất trong luận án và Chuyên đề 1, Nhóm IV
\texttt{Lower\_mid\_transition} (7 nền đông dân), Nhóm V
\texttt{Emerging} (17 nền), Nhóm VI \texttt{SIDS\_small} (8 nền). Nghiên
cứu chuyên sâu SIDS (P8) lấy mẫu chính 7 nền Thái Bình Dương và kiểm
định bền vững mở rộng 9 nền (thêm Comoros + Timor-Leste). Mẫu gộp phân
loại đầy đủ = 96.415 DN/52 nền; loại 2 nền ngoài châu Á (Turkey, Cyprus)
sang khung phân tích 50 nền Á--TBD (gồm Nhật Bản) = 88.869 DN.}

Ba thế hệ lược đồ WBES và giao thức hòa hợp: PICS3 (2006--2012): biến
FSTS từ d3b/d3c; TCI từ h1/h8/b8; DAI từ c22b only. Standardized
(2013--2017): biến tái định dạng nhưng tương thích. BREADY/BEE
(2018--2025): lược đồ mở rộng; DAI đầy đủ (c22b + k33 + k38); phân hệ
nhà quản trị chi tiết hơn.

\emph{Ghi chú về trọng số khảo sát:} Chuyên đề không áp dụng trọng số
mẫu WBES (\texttt{wt\_final}) trong các mô hình hồi quy vì ba lý do. Thứ
nhất, mục tiêu phân tích là xác định hệ số quan hệ quốc tế hóa và hiệu
quả tại cấp doanh nghiệp (conditional on observables), không phải ước
lượng thống kê mô tả đại diện tổng thể. Thứ hai, hai chiều hiệu ứng cố
định (quốc gia × năm) đã hấp thụ phần lớn biến thiên liên quan đến thiết
kế lấy mẫu phân tầng. Thứ ba, tích hợp trọng số mẫu vào mô hình hiệu ứng
cố định không đồng nhất gây phức tạp ước lượng mà không đem lại cải
thiện tính hợp lệ nội tại (Solon et al., 2015). Phân tích độ nhạy với
trọng số được đề xuất như kiểm định tương lai (xem Mục 2.4.2) để xác
nhận tính vững của kết quả so với mẫu không có trọng số.

Giao thức hòa hợp bao gồm: crosswalk biến giữa ba lược đồ; xử lý giá trị
thiếu; loại bỏ ngoại lệ tại phân vị 1/99 trong cụm nền kinh tế × năm; và
chuẩn hóa năng suất xuyên tiền tệ (xem dưới).

\hypertarget{ux111o-lux1b0ux1eddng-biux1ebfn-chi-tiux1ebft}{%
\subparagraph{Đo lường biến chi
tiết}\label{ux111o-lux1b0ux1eddng-biux1ebfn-chi-tiux1ebft}}

Biến phụ thuộc, ln(LP):

\[\ln(LP)_i = \ln\left(\frac{\text{doanh thu hằng năm}_i}{\text{lao động thường xuyên}_i}\right)\]

Doanh thu (d2) được báo cáo bằng \textbf{nội tệ} của từng nền kinh tế,
nên mức \(\ln(\mathrm{LP})\) thô không so sánh được trực tiếp xuyên quốc
gia. Để trung hòa khác biệt đơn vị tiền tệ và mặt bằng giá, năng suất
lao động được \textbf{chuẩn hóa z trong từng cụm nền kinh tế × năm}
trước khi đưa vào mô hình, và mọi mô hình đều có \textbf{hiệu ứng cố
định hai chiều (nền kinh tế × năm)} hấp thụ toàn bộ thành phần bất biến
theo nước (tiền tệ, mặt bằng giá, chế độ kế toán). Cách này bảo đảm hệ
số quan hệ quốc tế hóa và hiệu quả được nhận dạng từ biến thiên
\emph{nội bộ} nền kinh tế--năm, nhất quán với Phụ lục A của luận án (Mục
A.5) và nghiên cứu đa quốc gia P7. Mọi so sánh \textbf{mức} xuyên quốc
gia dùng tỷ suất không thứ nguyên (ROS) hoặc giá trị đã hiệu chỉnh sức
mua tương đương (PPP), không dùng mức \(\ln(\mathrm{LP})\) thô. Giá trị
được cắt đuôi (winsorize) tại phân vị 1/99 trong cụm nền kinh tế × năm
để giảm ảnh hưởng ngoại lệ. Lựa chọn năng suất lao động làm thước đo
hiệu quả phù hợp với dữ liệu WBES (Combs et al., 2005; Cusolito \&
Maloney, 2018; Richard et al., 2009).

Biến độc lập, FSTS:

\[FSTS_i = \frac{\text{doanh thu xuất khẩu trực tiếp + gián tiếp}_i}{\text{tổng doanh thu hằng năm}_i} = \frac{d3b_i + d3c_i}{100}\]

d3b là tỉ lệ phần trăm doanh thu xuất khẩu gián tiếp (qua trung gian) và
d3c là xuất khẩu trực tiếp, theo nhãn biến chính thức WBES; tổng hai
thành phần là định nghĩa tham gia xuất khẩu chuẩn của WBES, dùng thống
nhất với khung ước lượng đa quốc gia của luận án (P7). FSTS\_c = FSTS −
mean(FSTS) trong mỗi ô sóng-quốc gia trước khi tính FSTS². Tỷ trọng
doanh thu xuất khẩu là thước đo cường độ quốc tế hóa chuẩn của tài liệu
quốc tế hóa và hiệu quả (Lu \& Beamish, 2001; Sullivan, 1994); hai
nghiên cứu thành phần đơn quốc gia (P3 Việt Nam, P5 Trung Quốc) dùng
riêng d3c với luận chứng kênh xuất khẩu trực tiếp/gián tiếp (Hessels \&
Terjesen, 2010; Peng \& Ilinitch, 1998).

Bảng 2.9. Đo lường TCI theo mã WBES.

\begin{longtable}[]{@{}
  >{\raggedright\arraybackslash}p{(\columnwidth - 6\tabcolsep) * \real{0.1622}}
  >{\raggedright\arraybackslash}p{(\columnwidth - 6\tabcolsep) * \real{0.2432}}
  >{\raggedright\arraybackslash}p{(\columnwidth - 6\tabcolsep) * \real{0.1892}}
  >{\raggedright\arraybackslash}p{(\columnwidth - 6\tabcolsep) * \real{0.4054}}@{}}
\toprule\noalign{}
\begin{minipage}[b]{\linewidth}\raggedright
Item
\end{minipage} & \begin{minipage}[b]{\linewidth}\raggedright
Mã WBES
\end{minipage} & \begin{minipage}[b]{\linewidth}\raggedright
Mô tả
\end{minipage} & \begin{minipage}[b]{\linewidth}\raggedright
Schema có sẵn
\end{minipage} \\
\midrule\noalign{}
\endhead
\bottomrule\noalign{}
\endlastfoot
Chứng nhận ISO & b8 & Doanh nghiệp có chứng nhận ISO không? (nhị phân) &
PICS3, Std, BEE \\
R\&D & h8 & Có thực hiện R\&D nội bộ không? (nhị phân) & PICS3, Std,
BEE \\
Đổi mới sản phẩm & h1 & Đã ra sản phẩm/dịch vụ mới 3 năm qua? (nhị phân)
& PICS3, Std, BEE \\
Công nghệ nước ngoài & e6 & Sử dụng công nghệ nước ngoài có bản quyền?
(nhị phân) & PICS3, Std, BEE \\
\end{longtable}

TCI\_z = z-chuẩn hóa mean của ≥3/4 items không thiếu (Cronbach α ước
tính 0,55--0,65, nhất quán với tổng hợp formative, Coltman et al.,
2008). Construct phản ánh chiều sâu năng lực công nghệ nội tại theo
truyền thống năng lực hấp thụ và năng lực công nghệ hóa công nghiệp
(Cohen \& Levinthal, 1990; Lall, 1992).

\begin{quote}
\emph{Lưu ý phân biệt biến thể TCI:} Chuyên đề 2 dùng biến thể
\textbf{đầy đủ 4 mục} (≥3/4 của b8, h8, h1, e6) khả dụng trên mẫu gộp mô
tả; trong khi nghiên cứu đa quốc gia P7 và Mục 3.4.5.4 của luận án dùng
biến thể \textbf{thin 2 mục} (b8 chứng chỉ chất lượng + e6 công nghệ
nước ngoài) để tối đa hóa độ phủ trên 50 nền kinh tế. Cả hai cùng đo
construct formative \emph{năng lực công nghệ}; khác biệt về bộ mục phản
ánh ràng buộc khả dụng dữ liệu, không phải định nghĩa mâu thuẫn. Hệ số
trong chuyên đề được tính từ chính bộ 4 mục này.
\end{quote}

Bảng 2.10. Đo lường DAI theo mã WBES.

\begin{longtable}[]{@{}
  >{\raggedright\arraybackslash}p{(\columnwidth - 6\tabcolsep) * \real{0.1622}}
  >{\raggedright\arraybackslash}p{(\columnwidth - 6\tabcolsep) * \real{0.2432}}
  >{\raggedright\arraybackslash}p{(\columnwidth - 6\tabcolsep) * \real{0.1892}}
  >{\raggedright\arraybackslash}p{(\columnwidth - 6\tabcolsep) * \real{0.4054}}@{}}
\toprule\noalign{}
\begin{minipage}[b]{\linewidth}\raggedright
Item
\end{minipage} & \begin{minipage}[b]{\linewidth}\raggedright
Mã WBES
\end{minipage} & \begin{minipage}[b]{\linewidth}\raggedright
Mô tả
\end{minipage} & \begin{minipage}[b]{\linewidth}\raggedright
Schema có sẵn
\end{minipage} \\
\midrule\noalign{}
\endhead
\bottomrule\noalign{}
\endlastfoot
Có website & c22b & Doanh nghiệp có website không? (nhị phân) & PICS3,
Std, BEE \\
Thanh toán điện tử khách hàng & k33 & \% doanh thu thu qua e-payment
(liên tục) & BEE only \\
Thanh toán điện tử nhà cung cấp & k38 & \% mua hàng thanh toán qua
e-payment & BEE only \\
\end{longtable}

DAI\_z\_full (Bậc 1+2): z-mean(c22b + k33 + k38), chỉ có trong lược đồ
BEE (2019 trở đi). DAI\_z\_tier1 (Bậc 1 only): z-mean(c22b), có xuyên
tất cả lược đồ. Phân tầng đo lường số hóa theo Verhoef et al.~(2021);
vai trò công cụ số trong quốc tế hóa theo Banalieva \& Dhanaraj (2019).

Chiến lược đo lường DAI: sử dụng DAI\_z\_full làm biến chính trong mẫu
con BEE; DAI\_z\_tier1 cho so sánh xuyên lược đồ và kiểm định độ vững.

\hypertarget{phuxe2n-loux1ea1i-icrv-6-nhuxf3m}{%
\subparagraph{Phân loại ICRV 6
nhóm}\label{phuxe2n-loux1ea1i-icrv-6-nhuxf3m}}

Bảng 2.11. Phân loại 50 nền kinh tế theo ICRV, đặc điểm và tiêu chí.

\begin{longtable}[]{@{}
  >{\raggedright\arraybackslash}p{(\columnwidth - 6\tabcolsep) * \real{0.2895}}
  >{\raggedright\arraybackslash}p{(\columnwidth - 6\tabcolsep) * \real{0.1579}}
  >{\raggedright\arraybackslash}p{(\columnwidth - 6\tabcolsep) * \real{0.2368}}
  >{\raggedright\arraybackslash}p{(\columnwidth - 6\tabcolsep) * \real{0.3158}}@{}}
\toprule\noalign{}
\begin{minipage}[b]{\linewidth}\raggedright
Nhóm ICRV
\end{minipage} & \begin{minipage}[b]{\linewidth}\raggedright
Nhãn
\end{minipage} & \begin{minipage}[b]{\linewidth}\raggedright
Tiêu chí
\end{minipage} & \begin{minipage}[b]{\linewidth}\raggedright
Nền kinh tế
\end{minipage} \\
\midrule\noalign{}
\endhead
\bottomrule\noalign{}
\endlastfoot
I & Tiên tiến đổi mới sáng tạo dẫn dắt & GNI/người \textgreater25.000
PPP + Innovation-driven (WIPO) & Singapore, Hong Kong SAR, Hàn Quốc, Đài
Loan, Israel, Nhật Bản \\
II & Tiên tiến tài nguyên dẫn dắt & GNI/người \textgreater25.000 PPP +
Doanh thu tài nguyên \textgreater40\% GDP & Saudi Arabia, Qatar, Kuwait,
Bahrain, Brunei, Oman \\
III & Trung bình cao & GNI/người 10.000--25.000 PPP; chuyển đổi công
nghiệp & Trung Quốc, Malaysia, Thái Lan, Kazakhstan, Armenia, Georgia \\
IV & Chuyển đổi TB-thấp (\texttt{Lower\_mid\_transition}) & GNI/người
1.500--5.000 PPP; đông dân, GVC đang tăng & Việt Nam, Ấn Độ, Indonesia,
Philippines, Mông Cổ, Bangladesh, Pakistan \\
V & Đang nổi (\texttt{Emerging}) & GNI/người \textless4.000 PPP; khoảng
trống thể chế đáng kể & Sri Lanka, Jordan, Lào, Campuchia, Myanmar,
Nepal, Bhutan, Maldives, Uzbekistan, Tajikistan, Kyrgyzstan,
Turkmenistan, Afghanistan, Azerbaijan, Iraq, Lebanon, Yemen (17 nước) \\
VI & SIDS (\texttt{SIDS\_small}) & Đảo nhỏ; phụ thuộc cao; quốc tế hóa
bắt buộc & Fiji, Kiribati, Papua New Guinea, Samoa, Solomon Islands,
Tonga, Vanuatu, Timor-Leste (8 nền; P8 mẫu chính 7 Pacific) \\
\end{longtable}

Tổng: 6 + 6 + 6 + 7 + 17 + 8 = \textbf{50 nền kinh tế châu Á--Thái Bình
Dương} (gồm Nhật Bản; căn theo nhãn \texttt{icrv\_label} dữ liệu P7).
\emph{(Mẫu gộp phân loại đầy đủ 52 nền/96.415 DN gồm thêm Turkey, Cyprus
ngoài châu Á, loại khỏi khung 50. P8 mẫu chính 7 nền Pacific; độ vững 9
nền gồm Comoros + Timor-Leste.)}

Phân loại dựa trên GNI/người PPP (trung bình 2020--2024, World Bank),
Chỉ số Đổi mới Sáng tạo Toàn cầu (WIPO GII), và tỉ trọng doanh thu tài
nguyên (IMF Fiscal Monitor). Tất cả 50 nền kinh tế không thay đổi nhóm
ICRV trong giai đoạn 2009--2025 (ngưỡng kiểm tra: chuyển nhóm nếu
GNI/người vượt ngưỡng ±15\% trong ≥3 năm liên tiếp, không có trường hợp
nào).

\hypertarget{chiux1ebfn-lux1b0ux1ee3c-nhux1eadn-dux1ea1ng-nhuxe2n-quux1ea3}{%
\subparagraph{Chiến lược nhận dạng nhân
quả}\label{chiux1ebfn-lux1b0ux1ee3c-nhux1eadn-dux1ea1ng-nhuxe2n-quux1ea3}}

Tầng 1, OLS với SE robust và hiệu ứng cố định (tất cả mô hình M0--M7):
HC1 SE (chuẩn), điều chỉnh không đồng nhất phương sai cấp doanh nghiệp;
hiệu ứng cố định hai chiều (quốc gia × năm); kiểm tra VIF, kỳ vọng
\textless{} 10 cho các biến trọng tâm.

Tầng 2, Biến công cụ (DAI trong M4, Nhóm IV): Ứng viên biến công cụ: (1)
mean DAI trung bình ngành loại trừ chính doanh nghiệp đó; (2) khoảng
cách từ thành phố của doanh nghiệp đến trạm cáp ngầm biển gần nhất ×
năm. F thống kê bậc một: phải ≥ 16 (ngưỡng Stock \& Yogo cho biến công
cụ yếu). Tham chiếu: P3 Việt Nam đã chứng minh F bậc một = 34,6 với biến
công cụ hợp lệ.

Tầng 3, Kiểm định phi tuyến:

Lind--Mehlum (2010) U-test: xác nhận chữ U ngược thực sự (không chỉ dấu
đúng). Quy trình: kiểm tra β₁ \textgreater{} 0 và β₂ \textless{} 0; tính
TP = −β₁/(2β₂) và 95\% CI bằng delta method; kiểm định H₀: độ dốc tại
min(FSTS) ≤ 0 HOẶC độ dốc tại max(FSTS) ≥ 0, bác bỏ H₀ sang chữ U ngược
thực sự.

Paternoster z-test cho so sánh điểm uốn:

\[z = \frac{\hat{\beta}_1 - \hat{\beta}_2}{\sqrt{SE_1^2 + SE_2^2}}\]

Áp dụng: so sánh TP(Nhóm I) vs TP(Nhóm IV) cho H5; so sánh hệ số FSTS²
giữa Year\_bucket cho H6.

LOWESS phi tham số: smoother cục bộ (bandwidth = 0,4) vẽ ln(LP)
\textasciitilde{} FSTS, nếu hình dạng visual nhất quán với chữ U ngược
sang bằng chứng xác nhận không phụ thuộc giả định phân phối.

Tầng 4, Nhân rộng mẫu con: Mẫu gộp đủ lớn (88.869 / 50 nền, gồm Nhật
Bản) cho phép chạy M0--M7 riêng cho từng nhóm ICRV (Nhóm I: 6.390; Nhóm
IV: 45.003) để kiểm định tính không đồng nhất.

\hypertarget{kux1ebf-houx1ea1ch-kiux1ec3m-ux111ux1ecbnh-ux111ux1ed9-vux1eefng}{%
\paragraph{2.3.7 Kế hoạch kiểm định độ
vững}\label{kux1ebf-houx1ea1ch-kiux1ec3m-ux111ux1ecbnh-ux111ux1ed9-vux1eefng}}

Kế hoạch kiểm định độ vững được thiết kế theo sáu nhóm, mỗi nhóm trả lời
một mối lo cụ thể.

Bảng 2.12. Kế hoạch kiểm định độ vững, sáu nhóm.

\begin{longtable}[]{@{}
  >{\raggedright\arraybackslash}p{(\columnwidth - 6\tabcolsep) * \real{0.1500}}
  >{\raggedright\arraybackslash}p{(\columnwidth - 6\tabcolsep) * \real{0.2000}}
  >{\raggedright\arraybackslash}p{(\columnwidth - 6\tabcolsep) * \real{0.4250}}
  >{\raggedright\arraybackslash}p{(\columnwidth - 6\tabcolsep) * \real{0.2250}}@{}}
\toprule\noalign{}
\begin{minipage}[b]{\linewidth}\raggedright
Nhóm
\end{minipage} & \begin{minipage}[b]{\linewidth}\raggedright
Mối lo
\end{minipage} & \begin{minipage}[b]{\linewidth}\raggedright
Kiểm định cụ thể
\end{minipage} & \begin{minipage}[b]{\linewidth}\raggedright
Kỳ vọng
\end{minipage} \\
\midrule\noalign{}
\endhead
\bottomrule\noalign{}
\endlastfoot
R1 Thay đổi thước đo & Biến phụ thuộc/độc lập/điều tiết khác nhau & (a)
Biến phụ thuộc sang ROS; (b) Biến độc lập sang biến giả xuất khẩu (dummy
exporter); (c) TCI 3-item (loại h1); (d) DAI Bậc 1 only (c22b) xuyên tất
cả schema & Dấu hệ số giữ nguyên; độ lớn có thể khác \\
R2 Lọc mẫu & Thiên lệch chọn lựa từ thành phần mẫu & (a) Loại doanh
nghiệp vi mô (\textless5 lao động); (b) Loại SOE (sở hữu nhà nước
\textgreater50\%); (c) Chỉ sản xuất (ISIC 10--33); (d) Chỉ doanh nghiệp
xuất khẩu (FSTS \textgreater{} 0) & Kết quả cốt lõi robust với các mẫu
khác nhau \\
R3 Thay đổi phương pháp & Vi phạm giả định OLS & (a) Hồi quy phân vị tại
p25/p50/p75; (b) HC3 thay HC1; (c) Bootstrap SE (1.000 lần); (d)
Clustered SE tại cấp quốc gia & Bảo toàn dấu; CI có thể rộng hơn \\
R4 Kiểm tra dạng hàm & Sai lệch dạng hàm & (a) LOWESS phi tham số; (b)
RESET test; (c) Đa thức phân số; (d) Hồi quy spline với 3 knots & LOWESS
nhất quán với fit bậc hai \\
R5 Placebo tests & Tương quan giả & (a) Thay FSTS bằng biến ngẫu nhiên
cho kết quả không tìm thấy phi tuyến; (b) Thay TCI bằng nhiễu cho kết
quả không có điều tiết; (c) Kiểm tra nhân quả đảo chiều & Placebo tests
không cho kết quả có ý nghĩa \\
R6 Nhạy cảm thiết kế & Tự do lựa chọn của nhà nghiên cứu & (a) Ngưỡng
ICRV khác (GNI thay PPP GDP); (b) Giai đoạn thời gian:
2009--11/12--16/17--25; (c) Loại từng quốc gia (jackknife 47 lần); (d)
Ngưỡng Tier DAI khác nhau & Kết quả cốt lõi ổn định \\
\end{longtable}

Ngưỡng chấp nhận cho các tham số quan trọng qua kiểm định độ vững:
β₁(FSTS) trong khoảng {[}0,1; 3,0{]}; β₂(FSTS²) trong khoảng {[}−2,5;
−0,1{]}; điểm uốn trong khoảng {[}20\%; 90\%{]}; β(TCI) trong khoảng
{[}0,05; 0,5{]}; β(FSTS²×DAI) Nhóm I trong khoảng {[}0,5; 6,0{]}.

Quy trình báo cáo: nếu ≥4/6 nhóm R1--R6 cho kết quả nhất quán với chuẩn
(dấu đúng, p \textless{} 0,10), kết quả được coi là ``xác nhận mạnh''.
Nếu chỉ 2--3/6 nhóm nhất quán, kết quả là ``hỗ trợ có điều kiện''. Nếu
\textless{} 2/6 nhóm nhất quán, kết quả cần tái xem xét và không được
tuyên bố giả thuyết được xác nhận.

\textbf{Cập nhật kết quả thực hiện (khung 50 nền P7).} Bốn nhóm kiểm
định cốt lõi của kế hoạch trên đã được thực hiện trên khung ước lượng 50
nền của luận án, tái lập đúng mô hình canonical (M2: điểm uốn 51,5\%,
\(\beta_2 = -1{,}399\), \(p < {,}001\)). \emph{Về suy luận cụm nhỏ
(R3):} bootstrap cụm hoang dã (wild-cluster, Cameron--Gelbach--Miller, B
= 9.999, cụm theo nền kinh tế) cho thấy độ cong theo từng nhóm ICRV nhạy
cảm với số cụm; độ cong vùng chuyển đổi (Nhóm IV) nới từ \(p = 0{,}001\)
(CRV1) lên \(\approx 0{,}10\), xác nhận suy luận phi tuyến hình thức
được neo ở cấp gộp (49 cụm) và phân tích tổng hợp P6, còn các điểm uốn
theo nhóm được đọc như một phổ mô tả chứ không phải kiểm định độc lập đủ
công suất. \emph{Về tự do lựa chọn của nhà nghiên cứu (R6):} hiệu chỉnh
đa kiểm định Holm/Bonferroni cho cả họ bốn tương tác TCI/DAI giữ nguyên
không có ý nghĩa, nên kết luận năng lực là yếu tố nâng mặt bằng không
phải âm tính giả do kiểm định lặp. \emph{Về bất biến đo lường (R1):} các
mục lõi TCI và DAI hiện diện khoảng \(99\%\) ở cả ba thế hệ lược đồ và
trung bình DAI tăng đơn điệu (\(0{,}30 \to 0{,}46 \to 0{,}53\)) phản ánh
khuếch tán số thực chất, nên so sánh chéo đợt là có cơ sở. \emph{Về nội
sinh chọn lọc (R4):} thêm tỷ số Mills nghịch đảo từ probit tham gia xuất
khẩu giữ nguyên độ cong chữ U ngược và điểm uốn
(\(43{,}6\% \to 43{,}5\%\)) với tỷ số Mills không có ý nghĩa
(\(p = {,}317\)), cho thấy không có thiên lệch tự chọn lọc kiểu Melitz
trong độ cong. Chi tiết phương pháp và số liệu ở
\texttt{data\_wbes/analysis/p7\_robustness\_suite\_2026-06.md}.

\hypertarget{thux1ea3o-luux1eadn-ux111uxf3ng-guxf3p-luxfd-thuyux1ebft-vuxe0-thux1ef1c-tiux1ec5n}{%
\paragraph{2.3.8 Thảo luận: đóng góp lý thuyết và thực
tiễn}\label{thux1ea3o-luux1eadn-ux111uxf3ng-guxf3p-luxfd-thuyux1ebft-vuxe0-thux1ef1c-tiux1ec5n}}

\hypertarget{khung-ux111uxe1nh-giuxe1-ux111uxf3ng-guxf3p}{%
\subparagraph{Khung đánh giá đóng
góp}\label{khung-ux111uxe1nh-giuxe1-ux111uxf3ng-guxf3p}}

Tám đóng góp dưới đây được định vị theo hai tiêu chí chuẩn để đánh giá
đóng góp lý thuyết trong nghiên cứu quản trị và kinh doanh quốc tế:
\textbf{tính nguyên gốc và tính hữu dụng} (originality × utility; Corley
\& Gioia, 2011) và bộ ba \textbf{``cái gì -- như thế nào -- vì sao bây
giờ''} (Whetten, 1989). Cần phân biệt rõ giữa \textbf{đóng góp} và
\textbf{nhân rộng} (replication): việc tái xác nhận chữ U ngược ở từng
bối cảnh đơn lẻ tự nó là nhân rộng có giá trị xác nhận, còn đóng góp nằm
ở việc \emph{cấu trúc hóa} sự biến thiên xuyên bối cảnh của dạng hàm
thành một quy luật kiểm định được. Vì chuyên đề này là chuyên đề
\textbf{xây dựng mô hình}, mỗi đóng góp được phát biểu dưới dạng
\textbf{mệnh đề có thể bác bỏ} kèm \textbf{điều kiện phạm vi} tường
minh, để khung CDCM--ICRV--FIP có thể sai nếu dữ liệu không khớp --- đặc
điểm phân biệt một mô hình lý thuyết với một mô tả hậu nghiệm.

\hypertarget{ux111uxf3ng-guxf3p-luxfd-thuyux1ebft}{%
\subparagraph{Đóng góp lý
thuyết}\label{ux111uxf3ng-guxf3p-luxfd-thuyux1ebft}}

Đóng góp 1: Khung tích hợp 4 tầng + Khung chấp nhận số nền tảng cho châu
Á. Không có khung tương đương trong tài liệu IB cho khu vực 50 nền kinh
tế (gồm Nhật Bản) bao gồm cả quốc đảo nhỏ Thái Bình Dương. Khung tích
hợp Uppsala + lý thuyết dựa trên nguồn lực + lý thuyết thể chế + lý
thuyết thượng tầng quản trị + Khung chấp nhận số nền tảng cho phép giải
thích đồng thời tính không đồng nhất quan sát được, điều mà các mô hình
đơn lý thuyết không làm được.

Đóng góp 2: CDCM, Hình thức hóa Mô hình Áp dụng số và Năng lực phụ thuộc
bối cảnh. CDCM được phát hiện trong phân tích mô tả Chuyên đề Tiến sĩ số
1 và được làm sắc nét trong chuyên đề này qua H3: DAI là nguồn lực mở
rộng tình huống (không phải phần thưởng đồng nhất), với tác động phụ
thuộc ba chiều ngữ cảnh (chế độ × mức FSTS × bão hòa số). CDCM lấp đầy
khoảng trống tại giao điểm của lý thuyết IB số và lý thuyết thể chế,
chưa có trong tài liệu.

Đóng góp 3: Phân loại nhánh Advanced thành Nhóm I và Nhóm II. Phát hiện
từ phân tích mô tả về tính không đồng nhất trong Advanced, Nhóm I (tiên
tiến đổi mới: Singapore, Korea, HK, Đài Loan) vs.~Nhóm II (tiên tiến tài
nguyên: Saudi, Qatar, Kuwait, Bahrain), được tích hợp vào H5 như dự đoán
lý thuyết: cùng GNI/người cao nhưng cơ chế quốc tế hóa và hiệu quả khác
nhau do nguồn lực dẫn dắt khác nhau.

Đóng góp 4: Giả thuyết gánh nặng bắt buộc trong khung quốc tế hóa và
hiệu quả. Dạng gánh nặng bắt buộc (forced penalty) ở SIDS (Nhóm VI) được
tích hợp vào H5/H1b như trường hợp cực đoan của phổ thể chế, SIDS không
có chữ U ngược điển hình do quốc tế hóa bắt buộc (Briguglio, 1995). Phát
hiện vững của nghiên cứu chuyên sâu SIDS (P8, trên \textbf{mẫu đầy đủ
bảy cụm Thái Bình Dương}) là \textbf{việc dạng chữ U ngược mất cấu
trúc}: độ dốc tuyến tính (−0,085; p\_wild = 0,66) không có ý nghĩa và độ
cong (+0,696; p\_wild = 0,082) chỉ đạt mức biên. Hệ số âm mạnh đặc trưng
cho FIP (β = −1,339) chỉ xuất hiện trên một bản dựng dữ liệu hạn chế ba
cụm, mang tính \textbf{minh họa và nhạy cảm với phiên bản dữ liệu}; do
đó FIP được phát biểu chủ yếu như một \textbf{đóng góp khái niệm} (điều
kiện biên cực trị) chứ không phải một ước lượng điểm chắc chắn. Tích hợp
này mở rộng tài liệu quốc tế hóa và hiệu quả từ châu Á lục địa sang các
quốc đảo nhỏ Thái Bình Dương.

\hypertarget{ux111uxf3ng-guxf3p-vux1ec1-muxf4-huxecnh-vuxe0-phux1b0ux1a1ng-phuxe1p}{%
\subparagraph{Đóng góp về mô hình và phương
pháp}\label{ux111uxf3ng-guxf3p-vux1ec1-muxf4-huxecnh-vuxe0-phux1b0ux1a1ng-phuxe1p}}

Đóng góp 5: Tám mô hình M0--M7 với điều tiết ba chiều (M7). Chuỗi mô
hình lồng nhau cho phép kiểm định từng cơ chế riêng biệt trước khi tích
hợp. M7 là mô hình đầu tiên trong tài liệu châu Á kiểm định đồng thời
I×TCI×DAI điều tiết ba chiều kết hợp với tính không đồng nhất thời gian,
cho phép kiểm tra xem TCI và DAI là bổ sung hay thay thế trong điều tiết
quốc tế hóa và hiệu quả.

Đóng góp 6: Mẫu gộp 88.869 doanh nghiệp (50 nền, gồm Nhật Bản), rộng
nhất cho nghiên cứu quốc tế hóa và hiệu quả châu Á. Ba thế hệ lược đồ
WBES được hòa hợp qua giao thức chi tiết. Mẫu gộp này đủ lớn để chạy
nhân rộng mẫu con theo ICRV và giai đoạn thời gian với power đầy đủ.

Đóng góp 7: Chiến lược nhận dạng đa tầng tích hợp biến công cụ. Kết hợp
HC1/HC3 SE robust, hiệu ứng cố định hai chiều, nhận dạng IV (F bậc một
\textgreater{} 16), và Paternoster z-test cho so sánh xuyên nhóm tạo ra
chiến lược nhận dạng mạnh nhất có thể với dữ liệu WBES cắt ngang. Đặc
biệt, việc tách DAI nhân quả vs.~phụ thuộc chọn lựa (từ bằng chứng P3
IV) là bước tiến quan trọng trong phân tích áp dụng số trong IB.

Đóng góp 8: Thiết kế giai đoạn thời gian ba giai đoạn. Phân chia 14 sóng
WBES thành ba giai đoạn (2009--12 / 2013--17 / 2018--25) và kiểm định
tương tác FSTS×Year\_bucket là thiết kế chưa được áp dụng trong tài liệu
quốc tế hóa và hiệu quả cho châu Á. Bằng chứng từ P3 Việt Nam và P5
Trung Quốc neo đậu kỳ vọng thực nghiệm.

\hypertarget{huxe0m-uxfd-chuxednh-suxe1ch}{%
\subparagraph{Hàm ý chính sách}\label{huxe0m-uxfd-chuxednh-suxe1ch}}

Cho doanh nghiệp Việt Nam và châu Á mới nổi:

Đầu tư vào TCI trước DAI: bằng chứng biến công cụ (P3 Việt Nam) cho thấy
chỉ TCI là nhân quả, trong khi DAI dương trong OLS là phụ thuộc chọn
lựa. Doanh nghiệp nên ưu tiên xây dựng năng lực công nghệ thực chất
(ISO, R\&D, công nghệ nước ngoài) trước khi đầu tư vào áp dụng số.

Ngưỡng điểm uốn 39--46\% FSTS cho Việt Nam: doanh nghiệp tiếp cận 40\%
FSTS cần đánh giá lại chiến lược quốc tế hóa, tăng cường năng lực điều
phối thay vì chỉ tăng tỉ lệ xuất khẩu.

DAI phát huy mạnh khi FSTS cao và môi trường thể chế tốt: Nhóm I
(Singapore) cho thấy DAI×FSTS² = +3,119, khi doanh nghiệp đã ở mức FSTS
cao và hoạt động trong môi trường thể chế tốt, DAI mới trở thành đòn bẩy
mở rộng.

Cho nhà hoạch định chính sách:

Phổ ICRV cho thấy chính sách hỗ trợ quốc tế hóa cần phân biệt theo chế
độ thể chế, những gì hiệu quả ở Singapore không nhất thiết hiệu quả ở
Bangladesh.

Tính không đồng nhất thời gian H6, chính sách năm 2025 không thể sao
chép chính sách năm 2009; AI bùng nổ tạo ra cơ hội mới cho DAI trong
giai đoạn 2018 trở đi.

SIDS gánh nặng bắt buộc, các quốc đảo nhỏ Thái Bình Dương cần khung
chính sách khác hoàn toàn, không áp dụng logic chữ U ngược chuẩn.

\begin{center}\rule{0.5\linewidth}{0.5pt}\end{center}

\hypertarget{kux1ebft-luux1eadn-vuxe0-ux111ux1ec1-xuux1ea5t}{%
\subsubsection{2.4 Kết luận và đề
xuất}\label{kux1ebft-luux1eadn-vuxe0-ux111ux1ec1-xuux1ea5t}}

\hypertarget{kux1ebft-luux1eadn}{%
\paragraph{2.4.1 Kết luận}\label{kux1ebft-luux1eadn}}

Chuyên đề này xây dựng một mô hình nghiên cứu toàn diện về quan hệ quốc
tế hóa -- hiệu quả (I--P) cho châu Á và Thái Bình Dương. Mô hình tích
hợp năm tầng lý thuyết, Uppsala, lý thuyết dựa trên nguồn lực, lý thuyết
thể chế, lý thuyết thượng tầng quản trị, Khung chấp nhận số nền tảng,
vào tám mô hình thực nghiệm M0--M7 có thể kiểm định trên mẫu gộp 88.869
doanh nghiệp WBES xuyên 50 nền kinh tế (gồm Nhật Bản) và 103 cặp quốc
gia-năm. Hệ giả thuyết H1--H6 được phát triển có neo đậu thực nghiệm từ
ba bản thảo đồng hành (P3 Việt Nam, P4 Singapore, P5 Trung Quốc), xác
nhận tính khả nghiệm của khung đề xuất trước khi áp dụng cho toàn mẫu
gộp.

Kết quả dự kiến từ bằng chứng neo đậu: H1 (phi tuyến) xác nhận chữ U
ngược ở hầu hết nhóm ICRV với phổ điểm uốn từ Nhóm I đến Nhóm VI; H2
(TCI nâng mặt bằng) TCI dương bền vững và nhân quả; H3a (DAI
Emerging/Frontier) khám phá; H3b (DAI Advanced) β(FSTS²×DAI) dương trong
Nhóm I (P4 xác nhận), null/phụ thuộc chọn lựa trong Nhóm IV (P3 IV
null); H4a (Kinh nghiệm QL) exp\_manager×FSTS dương; H4b (Học vấn QL)
educ\_manager×FSTS dương; H4c (Giới tính) khám phá; H5 (ICRV phổ thể
chế) điểm uốn phổ thể chế từ Nhóm I đến VI, Nhóm VI gánh nặng bắt buộc;
H6 (Thời gian) DAI dịch chuyển 2009 xuống 2023, TCI tăng cường theo thời
gian.

Đóng góp trung tâm, CDCM, giải thích tại sao DAI hoạt động như nguồn lực
mở rộng tình huống chứ không phải phần thưởng đồng nhất: tác động phụ
thuộc đồng thời vào chế độ thể chế, mức FSTS, và mức độ bão hòa số của
nền kinh tế. Đây là bước tiến lý thuyết quan trọng, kết nối lý thuyết IB
số với lý thuyết thể chế trong một khung thống nhất và có thể kiểm định.

Thông điệp chính sách cốt lõi: không có một công thức duy nhất cho quốc
tế hóa thành công, chế độ thể chế, năng lực công nghệ nội tại, và mức độ
số hóa cùng quyết định hiệu quả của chiến lược xuất khẩu. Chính sách và
chiến lược doanh nghiệp cần được hiệu chỉnh theo vị trí trên bản đồ
ICRV.

\hypertarget{hux1ea1n-chux1ebf-nghiuxean-cux1ee9u}{%
\paragraph{2.4.2 Hạn chế nghiên
cứu}\label{hux1ea1n-chux1ebf-nghiuxean-cux1ee9u}}

Chuyên đề này có ba giới hạn suy diễn cần được thừa nhận rõ ràng theo
cấu trúc ba phần: điều không thể kết luận, phương pháp khắc phục, và
nguồn tham chiếu phương pháp.

\textbf{(a) Điều không thể kết luận từ khung nghiên cứu này:}

Thứ nhất, dữ liệu WBES không phải dữ liệu bảng chuẩn, mỗi sóng là thiết
kế cắt ngang với mẫu mới, không theo dõi cùng một doanh nghiệp qua thời
gian. Do đó, mô hình M0--M7 xác lập \emph{liên kết có điều kiện} chứ
không phải \emph{tác động nhân quả} giữa quốc tế hóa và hiệu quả; không
thể loại trừ khả năng doanh nghiệp vốn hiệu quả hơn sẽ xuất khẩu nhiều
hơn (reverse causality) mà không có biến công cụ hợp lệ hoặc thiết kế
thực nghiệm tự nhiên.

Thứ hai, DAI bị giới hạn ở Bậc 1--2 (website và e-payment) do hạn chế
lược đồ WBES. Không thể kết luận về vai trò của năng lực số Bậc 3--4,
tích hợp ERP, triển khai AI, quản trị dữ liệu xuyên biên giới, đối với
mối quan hệ quốc tế hóa và hiệu quả. Kết luận H3 chỉ hợp lệ trong ranh
giới đo lường Bậc 1--2; ngoài ranh giới này, CDCM chưa được kiểm định
thực nghiệm.

Thứ ba, FSTS (Foreign Sales to Total Sales) chỉ nắm bắt xuất khẩu hàng
hóa và dịch vụ trực tiếp, bỏ qua FDI, licensing, franchising, và các
hình thức quốc tế hóa phi xuất khẩu khác. H1 (phi tuyến) có thể có dạng
hàm khác nhau nếu đo bằng chỉ số quốc tế hóa đa chiều hơn như DOI hoặc
entropy index.

Thứ tư, biến nhà quản trị trong H4a--H4c bị hạn chế do WBES không đo
khoảng cách tâm lý, gắn kết mạng lưới quốc tế, hay cấu trúc hội đồng
quản trị đầy đủ. Kiểm định H4a--H4c vì vậy là \emph{chỉ báo gián tiếp}
của cơ chế lý thuyết thượng tầng quản trị (Hambrick \& Mason, 1984),
không phải kiểm định trực tiếp các cấu trúc lý thuyết gốc.

\textbf{(b) Phương pháp khắc phục cho nghiên cứu tương lai:}

Để giải quyết vấn đề nhân quả, các nghiên cứu tiếp theo có thể sử dụng:
(i) biến công cụ dựa trên tỷ lệ áp dụng trung bình ngành × vùng × sóng
như IV cho TCI và DAI (chiến lược đã áp dụng thành công trong P3 với
first-stage F = 22--35, vượt ngưỡng Staiger--Stock); (ii) thiết kế
difference-in-differences khai thác các cú sốc chính sách ngoại sinh (ví
dụ: CPTPP 2018, chiến tranh thương mại Mỹ-Trung 2018-2019) như ranh giới
thực nghiệm tự nhiên; hoặc (iii) phân tích panel dài với dữ liệu WBES
core panel 2023-2027 khi có sẵn.

Để mở rộng đo lường số hóa, nghiên cứu tương lai có thể kết hợp WBES với
dữ liệu nền tảng số cấp ngành (World Bank Digital Progress and Trends
Report, ITU DDI) hoặc thu thập dữ liệu Bậc 3--4 từ báo cáo thường niên
doanh nghiệp niêm yết. Khi lược đồ B-READY bổ sung thêm mô-đun về AI và
ERP, CDCM có thể được kiểm định đầy đủ hơn.

Để mở rộng thước đo quốc tế hóa vượt FSTS, nghiên cứu có thể tích hợp dữ
liệu UNCTAD FDI Bilateral Investment Matrix với WBES theo mã quốc
gia-ngành, hoặc dùng entropy index đa thị trường như trong Contractor et
al.~(2003).

\textbf{(c) Phương pháp và nguồn tham chiếu:}

Về nhân quả hóa quốc tế hóa và hiệu quả: Angrist \& Pischke (2009);
Oster (2019) về δ-stability bounds; Wolfolds \& Siegel (2019) về điều
kiện hợp lệ của Heckman. Về tổng hợp bằng chứng xuyên bối cảnh: Sutton
\& Higgins (2008), dữ liệu bảng cho phân tích tổng hợp; Borenstein et
al.~(2009), hồi quy phân tích tổng hợp theo hiệu ứng ngẫu nhiên. Về mở
rộng bậc số hóa: Verhoef et al.~(2021), phân cấp Bậc 1--4. Về mở rộng
quốc tế hóa đa chiều: Contractor et al.~(2003); Sullivan (1994) về DOI
chỉ số hợp thành.

\begin{center}\rule{0.5\linewidth}{0.5pt}\end{center}

\hypertarget{hux1b0ux1edbng-nghiuxean-cux1ee9u-tiux1ebfp-theo}{%
\paragraph{2.4.3 Hướng nghiên cứu tiếp
theo}\label{hux1b0ux1edbng-nghiuxean-cux1ee9u-tiux1ebfp-theo}}

Chuyên đề này đặt nền tảng lý thuyết và thực nghiệm cho năm hướng nghiên
cứu cụ thể trong chuỗi luận án và xa hơn:

\textbf{Hướng 1, P6: Tổng quan định lượng ba tầng (three-level MARA, mục
tiêu JWB):} Bản thảo P6 đã hoàn thành, mở rộng ICBEF 2025 cơ sở (Đỗ \&
Phan, 2024, k=113) lên \textbf{k=238} (K=288 effect sizes, 49 nền kinh
tế) với hệ thống tìm kiếm PRISMA 2020 trên Web of Science + Scopus
(1977--2026). Ba biến điều tiết mới: (i) cDAI, mức độ số hóa cấp quốc
gia; (ii) ICRV 5 nhóm; (iii) DPL giai đoạn. Kết quả MARA xác nhận: gộp
\emph{r}=0,074 (CI {[}0,060; 0,088{]}), I²=62,4\%; ICRV Q\_M=17,35
(p=.002), H1 điều tiết thể chế được hỗ trợ; cDAI và DPL không có ý nghĩa
(H2, H3 không được hỗ trợ). Phát hiện chính: lệch lạc công bố xác nhận
qua Begg--Mazumdar τ=−0,132 (p=.001), cắt và điền k=57 imputed,
\emph{r}\_adj=0,035, gộp effect thực tế thấp hơn khoảng 50\% so với ước
lượng thô. Đóng góp của P6 trong khuôn khổ CD2: xác nhận phổ thể chế
ICRV tổng thể (Q\_M có ý nghĩa) nhưng không đủ k ở mỗi nhóm để xác nhận
đường phổ thể chế đơn điệu E1a/E1b.

\textbf{Hướng 2, P7: Kiểm định H1--H6 toàn mẫu gộp 50 nền kinh tế (mục
tiêu IBR):} P7 đã ước lượng trên N=81.022 doanh nghiệp, 50 nền kinh tế
châu Á và Thái Bình Dương, 103 cặp nền kinh tế và năm, 6 nhóm ICRV. Kết
quả xác nhận: H1 (phi tuyến) điểm uốn M5 = 43,6\% (Lind--Mehlum
p\textless,001); H2 TCI hiệu ứng mức +0,141 (p\textless,001); H3 DAI
hiệu ứng mức +0,201 (p\textless,001), tương tác độ cong cùng chiều nhưng
không ý nghĩa; H5 ba vùng thể chế, chữ U ngược \textbf{rõ} từ
SIDS/Frontier (\textasciitilde55\%) xuống tiên tiến đổi mới
(\textasciitilde28\%); DAI×ICRV p=.012 (digital shield strongest khi thể
chế yếu). P3 (Nhóm IV), P4 (Nhóm I), P5 (Nhóm III) cung cấp neo đậu
cross-context đã được xác nhận bởi gộp 50-economy results (gồm Nhật
Bản).

\textbf{Hướng 3, P8: SIDS Thái Bình Dương, gánh nặng bắt buộc (mục tiêu
World Development):} P8 đã hoàn thành, \textbf{mẫu chính 7 nền Pacific
SIDS (mẫu gộp N = 1.450; 212 cơ sở có xuất khẩu, 14,6\%), β(FSTS\_c) =
−1,339} (bản dựng ba cụm, minh họa); kiểm định độ vững mở rộng 9 nền
(thêm Comoros + Timor-Leste, N≈1.469, β≈−0,4), minh họa cho gánh nặng
quốc tế hóa cưỡng bức như trường hợp giới hạn lý thuyết (trên mẫu đầy đủ
7 nền, dạng chữ U ngược mất cấu trúc; hệ số âm mạnh chỉ ở bản dựng hạn
chế ba cụm), nhất quán với giả thuyết gánh nặng bắt buộc (Briguglio,
1995; Bertram, 2006). Khoảng trống về điều kiện biên FIP (Nhóm VI SIDS)
trong CĐ2 đã được lấp đầy bởi P8.

\textbf{Hướng 4, CDCM Bậc 3--4 extension (MIS Quarterly hoặc Information
Systems Research):} Khi lược đồ B-READY (World Bank, 2023--2027) và WBES
mới bổ sung thêm mô-đun AI, ERP, và quản trị dữ liệu, CDCM có thể được
kiểm định đầy đủ bốn tầng Verhoef et al.~(2021). Câu hỏi lý thuyết: liệu
DAI Bậc 3--4 có chuyển vai trò từ nguồn lực tình huống (chỉ phát huy ở
FSTS cao) thành nguồn lực nền tảng (nâng mặt bằng đồng nhất), và điều
kiện chuyển đổi là gì (ngưỡng thể chế, quy mô doanh nghiệp, hay sóng
công nghệ)?

\textbf{Hướng 5, Dữ liệu bảng dọc (panel longitudinal) WBES core
2023--2027 (JIBS hoặc SMJ):} tập dữ liệu bảng lõi (WBES core panel) theo
dõi một nhóm doanh nghiệp trong ít nhất hai sóng liên tiếp, cho phép
kiểm định nhân quả TCI sang năng suất và DAI sang FSTS (theo hướng: số
hóa có thúc đẩy xuất khẩu không, thay vì chỉ làm tăng năng suất?) không
phụ thuộc vào giả định IV. Đây là hướng nghiên cứu thiết yếu để chuyển
H2 và H3 từ liên kết có điều kiện thành tác động nhân quả được xác nhận.

\begin{center}\rule{0.5\linewidth}{0.5pt}\end{center}

\hypertarget{tuxe0i-liux1ec7u-tham-khux1ea3o}{%
\subsection{TÀI LIỆU THAM KHẢO}\label{tuxe0i-liux1ec7u-tham-khux1ea3o}}

Angrist, J. D., \& Pischke, J.-S. (2009). \emph{Mostly harmless
econometrics: An empiricist's companion}. Princeton University Press.

Banalieva, E. R., \& Dhanaraj, C. (2019). Internalization theory for the
digital economy. \emph{Journal of International Business Studies,
50}(8), 1372--1387. \url{https://doi.org/10.1057/s41267-019-00243-7}

Barney, J. (1991). Firm resources and sustained competitive advantage.
\emph{Journal of Management, 17}(1), 99--120.
\url{https://doi.org/10.1177/014920639101700108}

Bausch, A., \& Krist, M. (2007). The effect of context-related
moderators on the internationalization-performance relationship:
Evidence from meta-analysis. \emph{Management International Review,
47}(3), 319--347.

Bertram, G. (2006). Introduction: The MIRAB model in the twenty-first
century. \emph{Asia Pacific Viewpoint, 47}(1), 1--13.

Beugelsdijk, S., Kostova, T., Kunst, V. E., Spadafora, E., \& van Essen,
M. (2018). Cultural distance and firm internationalization: A
meta-analytical review and theoretical implications. \emph{Journal of
Management, 44}(1), 89--130.
\url{https://doi.org/10.1177/0149206317729027}

Bhandari, K. R., Ranta, M., \& Salo, J. (2023). The resource-based view,
stakeholder capitalism, ESG, and sustainable competitive advantage: The
firm's embeddedness into ecology, society, and governance.
\emph{Business Strategy and the Environment, 32}(4), 1525--1541.

Bharadwaj, A., El Sawy, O. A., Pavlou, P. A., \& Venkatraman, N. (2013).
Digital business strategy: Toward a next generation of insights.
\emph{MIS Quarterly, 37}(2), 471--482.
\url{https://doi.org/10.25300/MISQ/2013/37:2.3}

Borenstein, M., Hedges, L. V., Higgins, J. P. T., \& Rothstein, H. R.
(2009). \emph{Introduction to meta-analysis}. Wiley.
\url{https://doi.org/10.1002/9780470743386}

Briguglio, L. (1995). Small island developing states and their economic
vulnerabilities. \emph{World Development, 23}(9), 1615--1632.
\url{https://doi.org/10.1016/0305-750X(95)00065-K}

Buckley, P. J., \& Casson, M. C. (1976). \emph{The future of the
multinational enterprise}. Macmillan.

Cannella, A. A., Jr., Park, J.-H., \& Lee, H.-U. (2008). Top management
team functional background diversity and firm performance: Examining the
roles of team member colocation and environmental uncertainty.
\emph{Academy of Management Journal, 51}(4), 768--784.
\url{https://doi.org/10.5465/amj.2008.33665310}

Caves, R. E. (1971). International corporations: The industrial
economics of foreign investment. \emph{Economica, 38}(149), 1--27.

Chen, S., \& Tan, H. (2012). Region effects in the
internationalization--performance relationship in Chinese firms.
\emph{Journal of World Business, 47}(1), 73--80.
\url{https://doi.org/10.1016/j.jwb.2010.10.019}

Cohen, W. M., \& Levinthal, D. A. (1990). Absorptive capacity: A new
perspective on learning and innovation. \emph{Administrative Science
Quarterly, 35}(1), 128--152. \url{https://doi.org/10.2307/2393553}

Coltman, T., Devinney, T. M., Midgley, D. F., \& Venaik, S. (2008).
Formative versus reflective measurement models: Two applications of
formative measurement. \emph{Journal of Business Research, 61}(12),
1250--1262. \url{https://doi.org/10.1016/j.jbusres.2008.01.013}

Combs, J. G., Crook, T. R., \& Shook, C. L. (2005). The dimensionality
of organizational performance and its implications for strategic
management research. In D. J. Ketchen \& D. D. Bergh (Eds.),
\emph{Research methodology in strategy and management} (Vol. 2,
pp.~259--286). Emerald.
\url{https://doi.org/10.1016/S1479-8387(05)02011-4}

Contractor, F. J. (2012). Why do multinational firms exist? A theory
note about the effect of multinational expansion on performance and
recent methodological critiques. \emph{Global Strategy Journal, 2}(4),
318--331.

Contractor, F. J., Kundu, S. K., \& Hsu, C.-C. (2003). A three-stage
theory of international expansion: The link between multinationality and
performance in the service sector. \emph{Journal of International
Business Studies, 34}(1), 5--18.
\url{https://doi.org/10.1057/palgrave.jibs.8400003}

Corley, K. G., \& Gioia, D. A. (2011). Building theory about theory
building: What constitutes a theoretical contribution? \emph{Academy of
Management Review, 36}(1), 12--32.
\url{https://doi.org/10.5465/amr.2009.0486}

Cronbach, L. J. (1951). Coefficient alpha and the internal structure of
tests. \emph{Psychometrika, 16}(3), 297--334.
\url{https://doi.org/10.1007/BF02310555}

Cusolito, A. P., \& Maloney, W. F. (2018). \emph{Productivity revisited:
Shifting paradigms in analysis and policy}. World Bank.
\url{https://doi.org/10.1596/978-1-4648-1334-4}

Đỗ, T. H., \& Phan, A. T. (2026a). \emph{Revisiting the
internationalisation--performance relationship in an emerging market:
The roles of technological capability and foundational digital adoption}
{[}Bản thảo chưa nộp{]}. \emph{Journal of Economics and Development}.
{[}P3, Việt Nam{]}

Đỗ, T. H., \& Phan, A. T. (2026b). \emph{Technological capability,
digital adoption, and internationalization performance in China:
Evidence from two WBES waves} {[}Bản thảo chưa nộp{]}.
\emph{International Journal of Emerging Markets}. {[}P5, Trung Quốc{]}

Đỗ, T. H., \& Phan, A. T. (2026c). Unveiling the impact of Chinese
manufacturing SMEs' internationalization on performance. \emph{Journal
of Finance and Accounting Research}. Advance online publication. {[}P2,
Trung Quốc SMEs, đã công bố{]}

Dubin, R. (1978). \emph{Theory building} (Rev.~ed.). Free Press.

Dunning, J. H. (1988). The eclectic paradigm of international
production: A restatement and some possible extensions. \emph{Journal of
International Business Studies, 19}(1), 1--31.
\url{https://doi.org/10.1057/palgrave.jibs.8490372}

Geringer, J. M., Beamish, P. W., \& daCosta, R. C. (1989).
Diversification strategy and internationalization: Implications for MNE
performance. \emph{Strategic Management Journal, 10}(2), 109--119.

Glaum, M., \& Oesterle, M.-J. (2007). 40 years of research on
internationalization and firm performance: More questions than answers?
\emph{Management International Review, 47}(3), 307--317.

Gomes, L., \& Ramaswamy, K. (1999). An empirical examination of the form
of the relationship between multinationality and performance.
\emph{Journal of International Business Studies, 30}(1), 173--187.
\url{https://doi.org/10.1057/palgrave.jibs.8490065}

Greene, W. H. (2018). \emph{Econometric analysis} (8th ed.). Pearson.

Hambrick, D. C. (2007). Upper echelons theory: An update. \emph{Academy
of Management Review, 32}(2), 334--343.
\url{https://doi.org/10.5465/amr.2007.24345254}

Hambrick, D. C., \& Mason, P. A. (1984). Upper echelons: The
organization as a reflection of its top managers. \emph{Academy of
Management Review, 9}(2), 193--206.
\url{https://doi.org/10.5465/amr.1984.4277628}

Hessels, J., \& Terjesen, S. (2010). Resource dependency and
institutional theory perspectives on direct and indirect export choices.
\emph{Small Business Economics, 34}(2), 203--220.
\url{https://doi.org/10.1007/s11187-008-9156-4}

Hitt, M. A., Hoskisson, R. E., \& Kim, H. (1997). International
diversification: Effects on innovation and firm performance in
product-diversified firms. \emph{Academy of Management Journal, 40}(4),
767--798. \url{https://doi.org/10.2307/256948}

Hsu, C.-C., \& Boggs, D. J. (2003). Internationalization and
performance: Traditional measures and their decomposition.
\emph{Multinational Business Review, 11}(3), 23--50.

Hymer, S. H. (1976). \emph{The international operations of national
firms: A study of direct foreign investment}. MIT Press.

Johanson, J., \& Vahlne, J.-E. (1977). The internationalization process
of the firm: A model of knowledge development and increasing foreign
market commitments. \emph{Journal of International Business Studies,
8}(1), 23--32. \url{https://doi.org/10.1057/palgrave.jibs.8490676}

Johanson, J., \& Vahlne, J.-E. (2009). The Uppsala internationalization
process model revisited: From liability of foreignness to liability of
outsidership. \emph{Journal of International Business Studies, 40}(9),
1411--1431. \url{https://doi.org/10.1057/jibs.2009.24}

Kafouros, M., Aliyev, M., \& Krammer, S. (2023). Do firms benefit from
country-level institutional reforms? The role of internal firm resources
and capabilities. \emph{Global Strategy Journal, 13}(1), 70--100.
\url{https://doi.org/10.1002/gsj.1433}

Khanna, T., \& Palepu, K. G. (2010). \emph{Winning in emerging markets:
A road map for strategy and execution}. Harvard Business Press.

Kim, G., Kumar, T., Ramalho, R., \& Russell, S. (2026, January).
\emph{Institutional capacity for policy implementation: An analytical
framework} (Policy Research Working Paper No.~11279). World Bank Group.

Kirca, A. H., Hult, G. T. M., Roth, K., Cavusgil, S. T., Perry, M. Z.,
Akdeniz, M. B., Deligonul, S. Z., Mena, J. A., Pollitte, W. A., Hoppner,
J. J., Miller, J. C., \& White, R. C. (2012). Firm-specific assets,
multinationality, and financial performance: A meta-analytic review and
theoretical integration. \emph{Academy of Management Journal, 55}(1),
47--73. \url{https://doi.org/10.5465/amj.2009.0703}

Knight, G. A., \& Cavusgil, S. T. (2004). Innovation, organizational
capabilities, and the born-global firm. \emph{Journal of International
Business Studies, 35}(2), 124--141.

Lall, S. (1992). Technological capabilities and industrialization.
\emph{World Development, 20}(2), 165--186.
\url{https://doi.org/10.1016/0305-750X(92)90097-F}

Li, J., Chen, L., Yi, J., Mao, J., \& Liao, J. (2022).
Ecosystem-specific advantages in international digital commerce.
\emph{Journal of International Business Studies, 53}(4), 677--695.

Lind, J. T., \& Mehlum, H. (2010). With or without U? The appropriate
test for a U-shaped relationship. \emph{Oxford Bulletin of Economics and
Statistics, 72}(1), 109--118.
\url{https://doi.org/10.1111/j.1468-0084.2009.00569.x}

Long, J. S., \& Ervin, L. H. (2000). Using heteroscedasticity consistent
standard errors in the linear regression model. \emph{The American
Statistician, 54}(3), 217--224.
\url{https://doi.org/10.1080/00031305.2000.10474549}

Lu, J. W., \& Beamish, P. W. (2001). The internationalization and
performance of SMEs. \emph{Strategic Management Journal, 22}(6--7),
565--586. \url{https://doi.org/10.1002/smj.184}

Lu, J. W., \& Beamish, P. W. (2004). International diversification and
firm performance: The S-curve hypothesis. \emph{Academy of Management
Journal, 47}(4), 598--609. \url{https://doi.org/10.2307/20159604}

Luo, Y., \& Tung, R. L. (2007). International expansion of emerging
market enterprises: A springboard perspective. \emph{Journal of
International Business Studies, 38}(4), 481--498.

Mar, T. T., Đỗ, T. H., \& Phan, A. T. (2026). \emph{Technological
capability, digital adoption, and the internationalization--performance
relationship: A firm-level study of Singapore} {[}Bản thảo chưa nộp{]}.
\emph{Management International Review}. {[}P4, Singapore{]}

Marano, V., Arregle, J.-L., Hitt, M. A., Spadafora, E., \& van Essen, M.
(2016). Home country institutions and the
internationalization-performance relationship: A meta-analytic review.
\emph{Journal of Management, 42}(5), 1075--1110.
\url{https://doi.org/10.1177/0149206315624963}

Mathews, J. A. (2002). \emph{Dragon multinationals: Towards a new model
of global growth}. Oxford University Press.

Melitz, M. J. (2003). The impact of trade on intra-industry
reallocations and aggregate industry productivity. \emph{Econometrica,
71}(6), 1695--1725.

Meyer, K. E., Mudambi, R., \& Narula, R. (2017). Multinational
enterprises and local contexts: The opportunities and challenges of
multiple embeddedness. \emph{Journal of Management Studies, 48}(2),
235--252.

Nielsen, B. B., \& Nielsen, S. (2011). The role of top management team
international orientation in international strategic decision-making:
The choice of foreign entry mode. \emph{Journal of World Business,
46}(2), 185--193. \url{https://doi.org/10.1016/j.jwb.2010.05.003}

North, D. C. (1990). \emph{Institutions, institutional change and
economic performance}. Cambridge University Press.
\url{https://doi.org/10.1017/CBO9780511808678}

Paternoster, R., Brame, R., Mazerolle, P., \& Piquero, A. (1998). Using
the correct statistical test for the equality of regression
coefficients. \emph{Criminology, 36}(4), 859--866.
\url{https://doi.org/10.1111/j.1745-9125.1998.tb01268.x}

Peng, M. W. (2003). Institutional transitions and strategic choices.
\emph{Academy of Management Review, 28}(2), 275--296.
\url{https://doi.org/10.5465/amr.2003.9416341}

Peng, M. W., \& Ilinitch, A. Y. (1998). Export intermediary firms: A
note on export development research. \emph{Journal of International
Business Studies, 29}(3), 609--620.
\url{https://doi.org/10.1057/palgrave.jibs.8490011}

Peng, M. W., Wang, D. Y. L., \& Jiang, Y. (2008). An institution-based
view of international business strategy: A focus on emerging economies.
\emph{Journal of International Business Studies, 39}(5), 920--936.
\url{https://doi.org/10.1057/palgrave.jibs.8400377}

Pisani, N., Garcia-Bernardo, J., \& Heemskerk, E. (2020). Does it pay to
be a multinational? A large-sample, cross-national replication assessing
the multinationality--performance relationship. \emph{Strategic
Management Journal, 41}(1), 152--172.
\url{https://doi.org/10.1002/smj.3087}

Riahi-Belkaoui, A. (1998). The effects of the degree of
internationalization on firm performance. \emph{International Business
Review, 7}(3), 315--321.

Richard, P. J., Devinney, T. M., Yip, G. S., \& Johnson, G. (2009).
Measuring organizational performance: Towards methodological best
practice. \emph{Journal of Management, 35}(3), 718--804.
\url{https://doi.org/10.1177/0149206308330560}

Ruigrok, W., \& Wagner, H. (2003). Internationalization and performance:
An organizational learning perspective. \emph{Management International
Review, 43}(1), 63--83.

Schwens, C., Zapkau, F. B., Brouthers, K. D., \& Hollender, L. (2018).
Limits to outside-in open innovation in SMEs: When it helps and when it
does not help with international performance. \emph{Journal of World
Business, 53}(2), 168--177.

Solon, G., Haider, S. J., \& Wooldridge, J. M. (2015). What are we
weighting for? \emph{Journal of Human Resources, 50}(2), 301--316.
\url{https://doi.org/10.3368/jhr.50.2.301}

Stallkamp, M., \& Schotter, A. P. J. (2021). Platforms without borders?
The international strategies of digital platform firms. \emph{Global
Strategy Journal, 11}(1), 58--80. \url{https://doi.org/10.1002/gsj.1336}

Sullivan, D. (1994). Measuring the degree of internationalization of a
firm. \emph{Journal of International Business Studies, 25}(2), 325--342.
\url{https://doi.org/10.1057/palgrave.jibs.8490203}

Sutton, A. J., \& Higgins, J. P. T. (2008). Recent developments in
meta-analysis. \emph{Statistics in Medicine, 27}(5), 625--650.
\url{https://doi.org/10.1002/sim.2982}

Teece, D. J., Pisano, G., \& Shuen, A. (1997). Dynamic capabilities and
strategic management. \emph{Strategic Management Journal, 18}(7),
509--533.

Torraco, R. J. (2005). Writing integrative literature reviews:
Guidelines and examples. \emph{Human Resource Development Review, 4}(3),
356--367.

Verhoef, P. C., Broekhuizen, T., Bart, Y., Bhatt, A., Dong, J. Q.,
Fabian, N., \& Haenlein, M. (2021). Digital transformation: A
multidisciplinary reflection and research agenda. \emph{Journal of
Business Research, 122}, 889--901.
\url{https://doi.org/10.1016/j.jbusres.2019.09.022}

Vernon, R. (1966). International investment and international trade in
the product cycle. \emph{The Quarterly Journal of Economics, 80}(2),
190--207.

Wernerfelt, B. (1984). A resource-based view of the firm.
\emph{Strategic Management Journal, 5}(2), 171--180.
\url{https://doi.org/10.1002/smj.4250050207}

Whetten, D. A. (1989). What constitutes a theoretical contribution?
\emph{Academy of Management Review, 14}(4), 490--495.
\url{https://doi.org/10.5465/amr.1989.4308371}

Wolfolds, S. E., \& Siegel, J. (2019). Misaccounting for endogeneity:
The peril of relying on the Heckman two-step method without a valid
instrument. \emph{Strategic Management Journal, 40}(3), 432--462.
\url{https://doi.org/10.1002/smj.2995}

Wooldridge, J. M. (2010). \emph{Econometric analysis of cross section
and panel data} (2nd ed.). MIT Press.

Wu, J., Fan, D., \& Chen, X. (2022). Revisiting the
internationalization--performance relationship: A twenty-year
meta-analysis of emerging market multinationals. \emph{Management
International Review, 62}(2), 199--231.
\url{https://doi.org/10.1007/s11575-022-00462-x}

Xiao, S. S., Jeong, I., Moon, J. J., Chung, C. C., \& Chung, J. (2013).
Internationalization and performance of firms in China. \emph{Journal of
International Management, 19}(2), 118--137.
\url{https://doi.org/10.1016/j.intman.2013.01.003}

Xu, D. (2024). From de jure to de facto: Institutional distance and firm
internationalization. \emph{Global Strategy Journal, 14}(1), 1--28.

Yang, Z., Zhao, Y., \& Wei, Z. (2025). Digital capability as a moderator
of internationalization--performance in emerging Asia: Evidence from
WBES 2018--2023. \emph{Asia Pacific Journal of Management} {[}sắp xuất
bản{]}.

\begin{center}\rule{0.5\linewidth}{0.5pt}\end{center}

\hypertarget{phux1ee5-lux1ee5c}{%
\subsection{PHỤ LỤC}\label{phux1ee5-lux1ee5c}}

\hypertarget{phux1ee5-lux1ee5c-a-danh-suxe1ch-50-nux1ec1n-kinh-tux1ebf-theo-nhuxf3m-icrv}{%
\subsubsection{Phụ lục A: Danh sách 50 nền kinh tế theo nhóm
ICRV}\label{phux1ee5-lux1ee5c-a-danh-suxe1ch-50-nux1ec1n-kinh-tux1ebf-theo-nhuxf3m-icrv}}

ICRV Nhóm I, Tiên tiến đổi mới sáng tạo dẫn dắt (6 nền kinh tế):
Singapore, Hong Kong SAR, Hàn Quốc, Đài Loan, Israel, Nhật Bản (đợt khảo
sát WBES đầu tiên năm 2025).

ICRV Nhóm II, Tiên tiến tài nguyên dẫn dắt (6 nền kinh tế): Saudi
Arabia, Qatar, Kuwait, Bahrain, Brunei, Oman.

ICRV Nhóm III, Trung bình cao (6 nền kinh tế): Trung Quốc, Malaysia,
Thái Lan, Kazakhstan, Armenia, Georgia.

ICRV Nhóm IV, Chuyển đổi thu nhập trung bình--thấp (7 nền kinh tế): Việt
Nam, Ấn Độ, Indonesia, Philippines, Mông Cổ, Bangladesh, Pakistan.

ICRV Nhóm V, Đang nổi (17 nền kinh tế): Sri Lanka, Jordan, Lào,
Campuchia, Myanmar, Nepal, Bhutan, Maldives, Uzbekistan, Tajikistan,
Kyrgyzstan, Turkmenistan, Afghanistan, Azerbaijan, Iraq, Lebanon, Yemen.

ICRV Nhóm VI, SIDS, trường hợp biên (8 nền kinh tế): Fiji, Kiribati,
Papua New Guinea, Samoa, Solomon Islands, Tonga, Vanuatu, Timor-Leste.
\emph{(Ghi chú: mẫu chính của nghiên cứu chuyên sâu SIDS là 7 nền Thái
Bình Dương; kiểm định độ vững mở rộng 9 nền, bổ sung Comoros.)}

Tổng: 6 + 6 + 6 + 7 + 17 + 8 = 50 nền kinh tế.

\hypertarget{phux1ee5-lux1ee5c-b-tuxf3m-tux1eaft-ba-bux1ea3n-thux1ea3o-ux111ux1ed3ng-huxe0nh}{%
\subsubsection{Phụ lục B: Tóm tắt ba bản thảo đồng
hành}\label{phux1ee5-lux1ee5c-b-tuxf3m-tux1eaft-ba-bux1ea3n-thux1ea3o-ux111ux1ed3ng-huxe0nh}}

P3, Việt Nam (mục tiêu JED): Đỗ, T. H., \& Phan, A. T. (2026a). Phân
tích 2.958 doanh nghiệp WBES ba sóng 2009/2015/2023. Kết quả chính: chữ
U ngược xác nhận (Lind--Mehlum p \textless{} 0,001 tổng hợp); điểm uốn
46,2\% (2009), 39,3\% (2015), 41,6\% (2023), 39,7\% (tổng hợp); TCI nhân
quả (β = 1,639, F bậc một = 22,1); DAI phụ thuộc chọn lựa (β\_IV =
0,018, F bậc một = 34,6); Paternoster z = 3,353 giữa điểm uốn 2009 và
2015 (p \textless{} 0,001).

P4, Singapore (mục tiêu MIR): Mar, T. T., Đỗ, T. H., \& Phan, A. T.
(2026). Phân tích 623 doanh nghiệp WBES 2023. Kết quả chính: đường cong
I--P chủ yếu dương; điểm uốn hàm ý \textasciitilde88,6\% (vùng thưa dữ
liệu; Lind--Mehlum p = 0,303); DAI×FSTS² = +3,119 (p = 0,005), DAI là
nguồn lực mở rộng tình huống; TCI β = 0,153.

P5, Trung Quốc (mục tiêu IJOEM): Đỗ, T. H., \& Phan, A. T. (2026b). Phân
tích 4.544 doanh nghiệp WBES hai sóng 2012/2024. Kết quả chính: chữ U
ngược bền vững cấu trúc, điểm uốn 49,4\% (2012), 47,2\% (2024),
Paternoster FSTS² p = 0,545 (không có sự khác biệt đáng kể); TCI tăng
cường theo thời gian (+0,260 sang +0,426, Paternoster p = 0,011).

\hypertarget{phux1ee5-lux1ee5c-c-quy-ux1b0ux1edbc-buxe1o-cuxe1o-thux1ed1ng-kuxea}{%
\subsubsection{Phụ lục C: Quy ước báo cáo thống
kê}\label{phux1ee5-lux1ee5c-c-quy-ux1b0ux1edbc-buxe1o-cuxe1o-thux1ed1ng-kuxea}}

Hệ số hồi quy trình bày không chuẩn hóa với HC1 SE trong ngoặc. Điểm uốn
trình bày trên thang FSTS gốc (0--1) với 95\% CI bằng delta method. Mức
ý nghĩa: *** p \textless{} 0,001; ** p \textless{} 0,01; * p \textless{}
0,05; † p \textless{} 0,10. R² báo cáo cho từng mô hình cộng với ΔR²
tăng thêm.

Kiểm định Lind--Mehlum: báo cáo p-value cho H₀ (không có chữ U ngược),
điểm uốn, và 95\% CI điểm uốn.

Kiểm định Paternoster: báo cáo z-statistic và p-value hai chiều cho so
sánh hệ số giữa hai nhóm hoặc hai giai đoạn.

\hypertarget{phux1ee5-lux1ee5c-d-thiux1ebft-lux1eadp-stata-sai-sux1ed1-chuux1ea9n-cluster}{%
\subsubsection{Phụ lục D: Thiết lập Stata: Sai số chuẩn
cluster}\label{phux1ee5-lux1ee5c-d-thiux1ebft-lux1eadp-stata-sai-sux1ed1-chuux1ea9n-cluster}}

Toàn bộ pipeline phân tích sử dụng Stata 18. Mô hình chuẩn dùng HC1
robust SE cho điều chỉnh phương sai không đồng nhất cấp doanh nghiệp. Để
kiểm định độ vững với clustering ở cấp quốc gia × ngành (country ×
industry cluster), thêm bước sau trước khi chạy mô hình hồi quy:

\begin{Shaded}
\begin{Highlighting}[]
\NormalTok{* === Tạo biến }\KeywordTok{cluster}\NormalTok{ country × industry ===}
\KeywordTok{gen}\NormalTok{ cluster\_id = country\_iso3 + }\StringTok{"\_"}\NormalTok{ + }\FunctionTok{string}\NormalTok{(sector\_main)}

\NormalTok{* === M1 với }\KeywordTok{cluster} \KeywordTok{SE}\NormalTok{ (kiểm định H1) ===}
\KeywordTok{reg}\NormalTok{ ln\_lp fsts\_c fsts\_c2 ln\_emp firm\_age foreign\_own i.country i.}\FunctionTok{year}\NormalTok{, }\CommentTok{///}
    \KeywordTok{vce}\NormalTok{(}\KeywordTok{cluster}\NormalTok{ cluster\_id)}

\NormalTok{* === M2 với }\KeywordTok{cluster} \KeywordTok{SE}\NormalTok{ (kiểm định H2: TCI) ===}
\KeywordTok{reg}\NormalTok{ ln\_lp fsts\_c fsts\_c2 tci\_z c.fsts\_c\#c.tci\_z }\CommentTok{///}
\NormalTok{    ln\_emp firm\_age foreign\_own i.country i.}\FunctionTok{year}\NormalTok{, }\CommentTok{///}
    \KeywordTok{vce}\NormalTok{(}\KeywordTok{cluster}\NormalTok{ cluster\_id)}

\NormalTok{* === M3 với }\KeywordTok{cluster} \KeywordTok{SE}\NormalTok{ (kiểm định H3: DAI) ===}
\KeywordTok{reg}\NormalTok{ ln\_lp fsts\_c fsts\_c2 dai\_z c.fsts\_c\#c.dai\_z c.fsts\_c2\#c.dai\_z }\CommentTok{///}
\NormalTok{    ln\_emp firm\_age foreign\_own i.country i.}\FunctionTok{year}\NormalTok{, }\CommentTok{///}
    \KeywordTok{vce}\NormalTok{(}\KeywordTok{cluster}\NormalTok{ cluster\_id)}

\NormalTok{* === M5 với }\KeywordTok{cluster} \KeywordTok{SE}\NormalTok{ (kiểm định H4a–H4c) ===}
\NormalTok{* H4a: kinh nghiệm}
\KeywordTok{reg}\NormalTok{ ln\_lp fsts\_c fsts\_c2 exp\_manager c.fsts\_c\#c.exp\_manager }\CommentTok{///}
\NormalTok{    ln\_emp firm\_age foreign\_own i.country i.}\FunctionTok{year}\NormalTok{, }\CommentTok{///}
    \KeywordTok{vce}\NormalTok{(}\KeywordTok{cluster}\NormalTok{ cluster\_id)}

\NormalTok{* H4b: học vấn}
\KeywordTok{reg}\NormalTok{ ln\_lp fsts\_c fsts\_c2 educ\_manager c.fsts\_c\#c.educ\_manager }\CommentTok{///}
\NormalTok{    ln\_emp firm\_age foreign\_own i.country i.}\FunctionTok{year}\NormalTok{, }\CommentTok{///}
    \KeywordTok{vce}\NormalTok{(}\KeywordTok{cluster}\NormalTok{ cluster\_id)}

\NormalTok{* H4c: giới tính (khám phá, hai phía, không đặt kỳ vọng)}
\KeywordTok{reg}\NormalTok{ ln\_lp fsts\_c fsts\_c2 gender\_manager c.fsts\_c\#c.gender\_manager }\CommentTok{///}
\NormalTok{    ln\_emp firm\_age foreign\_own i.country i.}\FunctionTok{year}\NormalTok{, }\CommentTok{///}
    \KeywordTok{vce}\NormalTok{(}\KeywordTok{cluster}\NormalTok{ cluster\_id)}

\NormalTok{* === Lind{-}Mehlum utest sau mỗi mô hình có FSTS² ===}
\NormalTok{* (cài đặt: }\KeywordTok{ssc}\NormalTok{ install utest)}
\NormalTok{utest fsts\_c fsts\_c2}

\NormalTok{* === Kiểm tra số }\KeywordTok{cluster}\NormalTok{ ===}
\NormalTok{* Số }\KeywordTok{cluster}\NormalTok{ tối thiểu khuyến nghị: ≥ 50 (Cameron \& Miller, 2015)}
\NormalTok{* Mẫu gộp 50 nền kinh tế × 5 ngành chính ≈ 250 cụm, thỏa mãn}
\KeywordTok{tab}\NormalTok{ cluster\_id }\KeywordTok{if} \FunctionTok{e}\NormalTok{(}\KeywordTok{sample}\NormalTok{) == 1, }\KeywordTok{sum}\NormalTok{(ln\_lp)}
\end{Highlighting}
\end{Shaded}

\textbf{Lưu ý kỹ thuật}: \texttt{vce(cluster\ cluster\_id)} điều chỉnh
cho tương quan trong cùng nhóm quốc gia-ngành (intra-cluster
correlation), khác với HC1 ở chỗ cho phép phương sai không đồng nhất lẫn
tương quan trong cluster. Kết quả kỳ vọng: SE có thể tăng lên so với HC1
trong các cluster có nhiều doanh nghiệp tương tự (ví dụ, nhóm Ấn Độ sản
xuất). Hệ số điểm ước lượng (point estimates) không thay đổi.

\textbf{Tham chiếu}: Cameron, A. C., \& Miller, D. L. (2015). A
practitioner's guide to cluster-robust inference. \emph{Journal of Human
Resources, 50}(2), 317--372.

\end{document}
